\documentclass[11pt]{article}
\usepackage[T1]{fontenc}
\usepackage[utf8]{inputenc}
\usepackage{lmodern}
\usepackage[margin=1in]{geometry}
\usepackage{amsmath,amssymb,amsthm,mathtools}
\usepackage{enumitem,booktabs,array,longtable,xcolor,microtype,etoolbox,url,needspace}
\usepackage[colorlinks=true,linkcolor=blue!50!black,citecolor=green!35!black,urlcolor=blue!60!black]{hyperref}
\usepackage[capitalize,nameinlink]{cleveref}
% ---------------------------------------------------------------------------
% Theorem environments
% ---------------------------------------------------------------------------
\theoremstyle{plain}
\newtheorem{theorem}{Theorem}[section]
\newtheorem{lemma}[theorem]{Lemma}
\newtheorem{corollary}[theorem]{Corollary}
\newtheorem{proposition}[theorem]{Proposition}
\newtheorem{conjecture}[theorem]{Conjecture}
\newtheorem{question}[theorem]{Question}
\newtheorem{problem}[theorem]{Problem}
\theoremstyle{definition}
\newtheorem{definition}[theorem]{Definition}
\newtheorem{remark}[theorem]{Remark}
\newtheorem{example}[theorem]{Example}

% Keep a theorem heading, formal-status line and opening statement together.
\AtBeginEnvironment{theorem}{\Needspace{8\baselineskip}}
\AtBeginEnvironment{lemma}{\Needspace{8\baselineskip}}
\AtBeginEnvironment{corollary}{\Needspace{8\baselineskip}}
\AtBeginEnvironment{proposition}{\Needspace{8\baselineskip}}
\AtBeginEnvironment{example}{\Needspace{8\baselineskip}}

% The Lean proof is the proof; the text gives a faithful outline of it.
\newenvironment{outline}[1][Outline of the formal proof]{\begin{proof}[#1]}{\end{proof}}
% Results that are not formalized get complete proofs, labelled as such.
\newenvironment{informalproof}{\begin{proof}[Proof (not formalized)]}{\end{proof}}

% ---------------------------------------------------------------------------
% Lean references
%   \lean{name}            a declaration, relative to the namespace SharedDecoders
%   \formal{file}{names}   the formal counterpart of a numbered result: a file relative to
%                          lean/SharedDecoders/ and a comma-separated list of declarations
%   \partformal{file}{names}{gap}
%                          a numbered result whose formal counterpart covers all but the
%                          part described in {gap}; the gap is proved or cited in the text
%   \notformalized         a numbered result that has no formal counterpart
%   \cited{where}          a result of an earlier paper used without proof: where it is proved,
%                          and whether it is Lean-verified there
% ---------------------------------------------------------------------------
% Lean names are typeset like paths, with line breaks allowed after "_", "." and "/".
\DeclareUrlCommand\leanname{\def\UrlFont{\ttfamily\upshape}}
\renewcommand{\UrlBreaks}{\do\_\do\.\do\/\do\-}
\newcommand{\lean}[1]{\leanname{#1}}
\newcommand{\leanfile}[1]{\leanname{#1}}
\newcommand{\leanlist}[1]{\begingroup\def\leansep{}%
  \renewcommand*{\do}[1]{\leansep\leanname{##1}\def\leansep{, }}\docsvlist{#1}\endgroup}
\newcommand{\formal}[2]{\par\noindent{\small\upshape\raggedright\textsf{Lean:} \leanlist{#2}\unskip
  \nobreak\hfil\penalty50\hskip1em\hbox{}\nobreak\hfill[\leanfile{#1}]\parfillskip=0pt\par}%
  \smallskip\noindent\ignorespaces}
\newcommand{\partformal}[3]{\par\noindent{\small\upshape\raggedright\textsf{Lean (partial):} \leanlist{#2}\unskip
  \nobreak\hfil\penalty50\hskip1em\hbox{}\nobreak\hfill[\leanfile{#1}]\parfillskip=0pt\par}%
  \par\noindent{\small\upshape\textsf{Not formalized:} #3}\par\smallskip\noindent\ignorespaces}
\newcommand{\notformalized}{\par\noindent{\small\upshape\textsf{Not formalized;} a complete proof is given
  below.}\par\smallskip\noindent\ignorespaces}
\newcommand{\cited}[1]{\par\noindent{\small\upshape\textsf{Cited:} #1}\par\smallskip\noindent\ignorespaces}

% ---------------------------------------------------------------------------
% Notation
% ---------------------------------------------------------------------------
\newcommand{\N}{\mathbb{N}}
\newcommand{\Z}{\mathbb{Z}}
\newcommand{\Q}{\mathbb{Q}}
\newcommand{\R}{\mathbb{R}}
\newcommand{\C}{\mathbb{C}}
\newcommand{\F}{\mathbb{F}}
\newcommand{\T}{\mathsf{T}}
\newcommand{\clog}[1]{\lceil \log_2 #1 \rceil}
\newcommand{\abs}[1]{\lvert #1 \rvert}
\newcommand{\card}[1]{\lvert #1 \rvert}
\newcommand{\eps}{\varepsilon}
\newcommand{\prev}{\mathrm{prev}}
\newcommand{\cost}{\mathrm{cost}}
\newcommand{\gn}{\mathrm{gn}}
\newcommand{\Sh}{\mathrm{Sh}}
\newcommand{\cl}{\mathrm{cl}}
\newcommand{\Aut}{\mathrm{Aut}}
\newcommand{\minrk}{\mathrm{minrk}}
\newcommand{\MAIS}{\mathrm{MAIS}}
\newcommand{\Crux}{\mathrm{Crux}}
\newcommand{\TCP}{\mathrm{TCP}}
\DeclareMathOperator{\rank}{rank}
\DeclareMathOperator{\dist}{dist}
\DeclareMathOperator{\poly}{poly}

\title{Capacity and Composition of Shared Transposition Decoders}
\author{Hans Brende \and GPT-6 Astra (OpenAI)}
\date{November 2026}
\begin{document}
\maketitle
\begin{abstract}
A shared transposition decoder sees a cyclic column rotation of the complete
matrix and its receiver class. We study the dimension of linear sources
on which such a rule is correct and the composition of complete rules.
For two classes, arbitrary deterministic decoders have normalized source
dimension at most $41/64+5/(3\sqrt d)$, over every field and with coupled
entry lanes. For $M\le d/2$ classes, a difference-cover argument gives a
deficit of order $M^{-1/2}$. A cyclic lift preserves a fixed class count and
gives two-class binary sources of dimension $b(d^2/2+7d/6)$ on an infinite
family of dimensions, with a paid two-work-tape realization at bit and
canonical logarithmic widths. Complete composition has different constraints:
two linear factors with bounded class counts cannot transpose in unbounded
dimension. We prove quantitative rank and separately scoped nonlinear
obstructions using balance, literal periods and equivariant outer maps.
One shared decoder can also be correct on an almost full-entropy arbitrary
set, showing why entropy alone cannot replace source linearity.
These results strengthen the algebraic restrictions and clarify what a
native execution would have to supply; they do not identify arbitrary tape
computations with this model. All numbered results are Lean-verified except
the native two-work-tape implementation theorem, which has a complete written
proof and an explicit formalization todo.

\end{abstract}
\begin{quote}\small
\textbf{Contributions.} Hans Brende chose the problem and directed the research.
GPT-6 Astra (OpenAI) developed the mathematics, Lean formalizations and text
under his direction.
\end{quote}
\tableofcontents
\clearpage
\section{Introduction}\label{sec:introduction}
Matrix transposition exchanges row and column addresses while preserving every
entry. The series studies whether a fixed deterministic multitape Turing machine
can perform this operation in time linear in the bit length, the question
\textup{(T)} of~\cite{FirstPaper2026}. A candidate computation must both move
information and organize the local rules that use it. This paper isolates the
second obligation in the cyclic shared-decoder model introduced
in~\cite{SeventeenthPaper2026}.

One rule can be correct on a large source without being correct everywhere.
Nor does a large source for each stage guarantee that the stages compose to
transposition. We therefore keep two questions separate. Part I bounds and
constructs linear correctness sources for a single rule, allowing its decoders
to be arbitrary deterministic functions. Part II studies exact complete-map
factorizations, distinguishing linear factors from the nonlinear architectures
for which additional symmetry arguments apply.

\paragraph{Results and their use.}
The two-class ceiling improves the earlier $3/4+o(1)$ bound for linear
decoders to $41/64+o(1)$, now allowing arbitrary deterministic decoders.
The proof intersects translated sources inside the smaller invariant spaces
created by earlier intersections and retains a negative doubled-step term
before averaging. For $M\le d/2$, the many-class estimate gives an
inverse-square-root deficit. These are quantitative advances on the arbitrary-class
capacity question in \cite[Problem~9.1(a)]{SeventeenthPaper2026}; the half-rate
target remains open. Both bounds restrict the linear sources a construction
may use; neither is a native time lower bound.

On the construction side, a cyclic lift retains the number of receiver classes.
It gives a growing additive advantage over the interval bound with two classes,
while its normalized advantage tends to zero. The revised seventeenth paper
already gives unbounded additive excess for every fixed number of classes,
over every field, including $\F_2$\cite{SeventeenthPaper2026}.
Our contribution here is the general lift preserving the seed classes,
its exact periodic-capacity formula, and the full rule kernel of its binary
specialization; binary fixed-class excess itself is already known.
A paid realization also separates a finite algebraic description from a native
implementation. It remains correct only on its specified source.

The principal composition obstruction addresses the many-hop gap identified in
\cite[Section~7.4]{SeventeenthPaper2026}. It excludes exact products of two complete
linear rules with fixed class counts for all sufficiently large dimensions,
over every field and with coupled lanes. Its common-message lemma and the
rank-deficit extension explain the additional obligation imposed by a complete
intermediate state. Nonlinear statements require different hypotheses:
equivariant surrounding maps, actual commutation identities, or a specified
first factor. No unrestricted nonlinear two-stage exclusion is claimed.

The extensions test which premises a native reduction must preserve.
Even one shared decoder can work on an arbitrary correctness set with only
logarithmic entropy deficit, while deterministic line rules with a fixed
number of classes have asymptotic rate at most one half on linear sources.
Thus extracting a large fibre does not
by itself make the capacity theorem applicable. The exhibited cyclic lift's dimension is
now its entire agreement-kernel dimension. An all-width reflection family
and a sharper extension-linear residual theorem separate two different
linearity obligations. A complete nonlinear first factor can satisfy the
one-stage requirements yet force every class of its remaining factor to be
a singleton. These statements narrow specific interfaces; none identifies
the algebraic model with every native tape execution.

\paragraph{Proofs and formal status.}
A line marked \textsf{Lean} names the authoritative formal statement and proof;
the English outline is a reading aid. A partial line identifies the exact missing
step, and a written-only result has a complete written proof. Every numbered
mathematical result is formalized except the native two-work-tape implementation
in \cref{thm:paid-lift}; its transition-level execution proof remains a formalization
todo. The checked results include orbit enumeration, finite probability estimates, the actual
reflection/Fourier adapter, and complete residual accounting. All formal files use the shared
libraries of this repository and Mathlib~\cite{Mathlib2026}; standard algebra,
finite-field and Ramsey theorems are credited where used.

\paragraph{Organization.}
We first fix the common model and its two correctness notions. Part I develops
capacity bounds, structural reductions and constructions. Part II turns to
linear composition and nonlinear balance and period constraints. The final
main section states the remaining mathematical and native-machine interfaces.
For the source-linearity premise, read \cref{sec:set-capacity}; for decoder
linearity, read \cref{sec:field-linear-residual,sec:linearization,sec:prime-reflection}.
The actual diagonal defects and mask-method limits refine the capacity search
in \cref{sec:block-defects,sec:mask-methods,sec:structural-bounds}.
The remaining appendices test complete first factors and balanced nonlinear
maps. Finally \cref{sec:residuals} gives a conditional operational route:
an exhibited complete affine baseline turns its operator defect into a
lower bound on the joint nonlinear correction record.

\section{The shared-decoder model}\label{sec:model}
Let $d\ge1$ and write $I=\Z/d\Z$. For an entry alphabet $A$, the complete
state space is $\mathcal X=A^{I\times I}$. In the linear setting $A=F^b$,
where $F$ is a field and $b\ge1$; thus $N=bd^2$ is the dimension over $F$.
For binary tape inputs the bit length is $L=d^2b$. The canonical width in
\textup{(T)} is $b(d)=\max(1,\lceil\log_2d\rceil)$, but algebraic statements
will quantify over every positive $b$.

Define the column and row shifts and the transpose by
\[
 (T^sX)_{i,j}=X_{i,j+s},\qquad
 (R^sX)_{i,j}=X_{i+s,j},\qquad
 (\tau X)_{i,j}=X_{j,i}.
\]
On a row or column word, $\sigma^s$ shifts its single cyclic index in the same
orientation. Thus $\tau T^s=R^s\tau$, $\tau R^s=T^s\tau$, and $T^sR^t=R^tT^s$.

\begin{definition}[A shared rule and its correctness source]\label{def:shared-rule}
\formal{Model.lean}{ClassUniform,IsSharedCode}
A receiver partition $p:I\to\{1,\ldots,M\}$ uses every label.
A complete shared rule $K:\mathcal X\to\mathcal X$ has total functions
$f_J:\mathcal X\to A^I$ such that
\[
 \operatorname{col}_c K(X)=f_{p(c)}(T^cX).
\]
The decoder has no additional receiver-index argument. In the linear setting,
a subspace $C\le\mathcal X$ is a correctness source if
$f_{p(c)}(T^cX)=\operatorname{row}_cX$ for every $X\in C$ and every $c$.
Its dimension is denoted $D$.
\end{definition}

The predicates are distinct: a correctness source need not contain arbitrary
intermediate states of a composition. A linear source does not imply that its
decoders are linear. A \emph{linear complete rule} means that $K$ itself is
$F$-linear on all of $\mathcal X$. All definitions permit mixing between lanes.
A map $H$ is column-equivariant if $HT^s=T^sH$ for every $s$; one-class
complete rules are exactly these maps. Equivariance does not assert that the
map is cheap to compute on tapes.

All complete factorizations in Part II act on this same state space and use
this same column-view action, unless a theorem explicitly says otherwise.
For a finite alphabet, exact factorization of the bijection $\tau$ by total
square maps forces every factor to be bijective. Correctness on a proper
source, or a larger workspace with a promised initialization, does not imply
this conclusion. Algebraic row or column operations incur no native time
charge until a separately stated implementation supplies it.

\part{Capacity and constructions}
\section{Capacity of deterministic shared decoders}\label{sec:capacity}

Throughout this section $F$ is a field, $d,b\ge1$, and
$V=(F^b)^{(\Z/d)\times(\Z/d)}$ has dimension $N=bd^2$.
A partition $p$ of the receiver indices has $M$ nonempty classes.
The source $C\le V$ is linear, of dimension $D$; the shared decoders
$f_J:V\to(F^b)^d$ are arbitrary total functions satisfying
$f_{p(c)}(T^cX)=\operatorname{row}_cX$ for $X\in C$.
All lanes may be coupled. This distinction between a linear source and a
linear decoder will remain explicit.

\subsection{Intersect inside the covariance space}

For a nonempty set $E\subseteq\Z/d$, put
\[
 C_E=\bigcap_{e\in E}T^{-e}C.
\]
Join distinct indices $c,c'$ when $p(c)=p(c')$ and $c'-c\in E-E$;
write $\chi(E)$ for the number of components. Let $H_E$ be the subspace
defined by
\[
 \operatorname{row}_{c'}X=\sigma^{c'-c}\operatorname{row}_cX
 \quad\text{on every edge}.
\]

\begin{lemma}[Deterministic translate pigeonhole]\label{lem:deterministic-tp}
\formal{CapacityAmbient.lean}{IsSharedCode.component_finrank_le, finrank_covarianceSpace, covarianceSpace_rot}
For every nonempty $E$, $C_E\subseteq H_E$, and
$\dim H_E=bd\chi(E)$. Consequently
\[
 |E|D\le (|E|-1)N+bd\chi(E).
\]
Every $H_E$ is invariant under every column rotation $T^t$.
\end{lemma}
\begin{outline}
If $c+e=c'+e'$ with $e,e'\in E$ and $p(c)=p(c')$, apply the two
correct decoder equations to $T^eX,T^{e'}X\in C$. Their arguments are
the same full matrix, so
$\sigma^e\operatorname{row}_cX=\sigma^{e'}\operatorname{row}_{c'}X$.
This gives the edge equation without using linearity of the decoder.
After replacing row $c$ by $\sigma^{-c}\operatorname{row}_cX$, the
equations say simply that normalized rows are constant on each component.
One arbitrary vector of dimension $bd$ per component gives all solutions;
there is no further condition around a cycle. This proves the dimension.
Codimension is subadditive under intersection, and each $T^e$ is invertible,
so $\dim C_E\ge |E|D-(|E|-1)N$. Finally, applying the same rotation
$\sigma^t$ to all rows preserves every edge equation because the rotations
commute. This proves invariance.
\end{outline}

The covariance argument is that of the translate-pigeonhole theorem in
\cite{SeventeenthPaper2026}. Its use of decoder values does not require those
decoders to be linear. The smaller invariant ambient in the next theorem
is the additional resource.

\begin{theorem}[Iterated translate bound]\label{thm:iterated-tp}
\formal{CapacityIteration.lean}{IsSharedCode.iterated_finrank_le, componentCount_zero}
Choose shifts $t_1,\ldots,t_k$, allowing repetitions, and set
\[
 E_0=\{0\},\qquad E_j=E_{j-1}\cup(E_{j-1}+t_j),\qquad
 n_j=bd\chi(E_j).
\]
Then $n_0=N$ and
\begin{equation}\label{eq:iterated-tp}
 D\le\sum_{j=0}^{k-1}\frac{n_j}{2^{j+1}}+\frac{n_k}{2^k}.
\end{equation}
\end{theorem}
\begin{outline}
Write $C_j=C_{E_j}$. The set identity gives
$C_j=C_{j-1}\cap T^{-t_j}C_{j-1}$. Both spaces on the right lie in
$H_{E_{j-1}}$, by the preceding lemma and rotation invariance. They have
the same dimension. Grassmann's formula inside this ambient gives
$2\dim C_{j-1}\le n_{j-1}+\dim C_j$. Iterate these inequalities and
use $\dim C_k\le n_k$. Repeated subset sums are counted once in $E_j$;
their collisions do not affect either the set identity or the argument.
\end{outline}

The theorem bounds actual simultaneous intersections. It assumes neither
independence of the constraints nor correctness of $H_E$ itself. It can
be combined with every applicable direct translate bound.

\subsection{A uniform bound for two classes}

Suppose now that $p$ has at most two labels. For $s\in\Z/d$, let
$g_s=\gcd(d,s)$, where $g_0=d$, and
$Q_s=|\{c:p(c)\ne p(c+s)\}|$.
On the cycles of addition by $s$, let $R_j(s)$ count proper monochromatic
runs of length at least $j$. A proper run is preceded by a color change;
a monochromatic cycle contributes no runs. Thus $R_1(s)=Q_s$.

\begin{lemma}[Run components and doubled steps]\label{lem:binary-runs}
\formal{CapacityRunsComponents.lean}{componentCount_le_runs_add_two_gcd, componentCount_le_changes_add_gcd}
\formal{CapacityRuns.lean}{CapacityRuns.changeCount_double, CapacityRuns.weighted_run_budget}
For every integer $r\ge1$, the same-color graph with edges at $s$-cycle distance at most $r$ has
at most $R_r(s)+2g_s$ components. At radius one it has at most $Q_s+g_s$.
For every $s$, including zero and steps with short cycles,
\begin{align}
 Q_{2s}&=2R_2(s),\label{eq:doubled-runs}\\
 Q_s+R_2(s)+2R_4(s)+2R_{13}(s)&\le d.\label{eq:run-budget}
\end{align}
\end{lemma}
\begin{outline}
Work on one $s$-cycle. If its length is at most $2r$, the radius-$r$
graph is complete within each color and has at most two components.
Otherwise list the positions of one occurring color cyclically. Consecutive
positions are connected when their intervening gap is at most $r$.
Every other short edge is a path of these consecutive short edges, so it
adds no component merger. If $B>0$ gaps are longer than $r$, there are
$B$ components; if no gap is long, there is one. A singleton has its
one long cyclic gap and one component. For a nonconstant binary cycle,
these long gaps are exactly the opposite-color runs of length at least
$r$. Summing the two colors proves the first estimate. At radius one a
nonconstant cycle has exactly one component per run; a monochromatic
cycle adds one component and no change. Summing cycles proves both claims.

For the identity, write $z_c=1_{p(c)\ne p(c+s)}$.
Then $1_{p(c)\ne p(c+2s)}=z_c\mathbin{\mathrm{xor}}z_{c+s}$.
In a nonconstant, nonalternating cycle, the changes of $z$ number twice
its nonempty runs of ones, each ending at the start of a $p$-run of length
at least two. Constant and alternating cycles give zero on both sides.
This also handles cycles of lengths one and two. Finally a proper run
of length $L$ contributes
$1+1_{L\ge2}+2\,1_{L\ge4}+2\,1_{L\ge13}$ to the left of
\eqref{eq:run-budget}. Its values are $1,2,4,6$ on the successive
ranges $L=1$, $2\le L\le3$, $4\le L\le12$, $L\ge13$; each is at most
$L$. Monochromatic cycles contribute zero. Sum the run lengths.
\end{outline}

\begin{theorem}[Theorem C: the universal two-class ceiling]\label{thm:universal-capacity}
\formal{CapacityUniversal.lean}{IsSharedCode.two_class_capacity_bound}
For every field, every $d,b\ge1$, every partition into at most two classes,
and every linear source with deterministic shared decoders,
\begin{equation}\label{eq:universal-capacity}
 D\le\frac{41}{64}bd^2+\frac53 b d^{3/2}.
\end{equation}
In particular the normalized dimension is at most $41/64+o(1)$,
uniformly in the partition, field, word width and decoders.
\end{theorem}
\begin{outline}
Apply \cref{thm:iterated-tp} with successive shifts $s,3s,9s$.
The three difference sets contain every multiple of $s$ with coefficient
in $[-1,1]$, $[-4,4]$, and $[-13,13]$, respectively: this follows by
choosing difference digits from $\{-1,0,1\}$ in bases $1,3,9$.
Reduction modulo $d$ may identify elements, which causes no problem.
The component estimates give
\[
 D\le\frac12bd^2+\frac14bdQ_s+
 \frac18bd\bigl(R_4(s)+R_{13}(s)\bigr)+\frac34bdg_s.
\]
Use \eqref{eq:run-budget} and \eqref{eq:doubled-runs} to obtain
\begin{equation}\label{eq:one-step-universal}
 D\le\frac9{16}bd^2+\frac3{16}bdQ_s-\frac1{32}bdQ_{2s}
       +\frac34bdg_s.
\end{equation}
The negative term is retained before averaging.

Choose $a$ uniformly from all $d$ residues and put
$A_r=\mathbb E_a Q_{2^ra}$ and $B_r=\mathbb E_a\gcd(d,2^ra)$.
If the label sizes are $n_0,n_1$, ordered-pair counting gives
$A_0=2n_0n_1/d\le d/2$. Average \eqref{eq:one-step-universal} at
$s=2^ra$, then weight these inequalities by $w_r=(5/6)6^{-r}$ for
$r\ge0$. The weights sum to one and all terms are bounded for fixed
$d$. Since $(3/16)w_r=(1/32)w_{r-1}$ for $r\ge1$, the internal
$A_r$ terms cancel, leaving
\[
 D\le\frac{41}{64}bd^2+\frac34bd\sum_{r\ge0}w_r B_r.
\]
Equivalently, sum through a finite index and then let the index grow;
the terminal term has absolute value at most $bd^2w_r/32$ and tends to zero.
No uniformity of $2^ra$ on the whole cycle is assumed.

Let $a_r=\gcd(d,2^r)$ and $n=d/a_r$. Its actual distribution is uniform
on the multiples of $a_r$, so
$B_r=a_r\mathbb E_{u\in\Z/n}\gcd(n,u)$.
For each divisor $e$ of $n$, at most $n/e$ residues have gcd equal to
$e$. Thus the mean gcd is at most the number of divisors of $n$,
which is at most $2\sqrt n$ by pairing divisors across $\sqrt n$.
This includes $u=0$ and $n=1$. Consequently
\[
 B_r\le2\sqrt d\,2^{r/2},\qquad
 \sum_{r\ge0}w_rB_r\le
 \frac{5\sqrt d}{3(1-\sqrt2/6)}\le\frac{20}{9}\sqrt d.
\]
Substitution proves \eqref{eq:universal-capacity}. The proof works with
rows in $(F^b)^d$ throughout; no lane decomposition of $C$ was used.
\end{outline}

The bound improves the general two-class $3/4+o(1)$ ceiling for linear
decoders stated in \cite[Open problems]{SeventeenthPaper2026}, and allows
arbitrary deterministic decoders. It leaves the conjectured
$1/2+o(1)$ ceiling open. The parameter being bounded is the dimension of
a linear correctness source, not the density of native algorithms.

\subsection{Arbitrarily many classes}

\begin{theorem}[Theorem M: a square-root class deficit]\label{thm:many-classes}
\formal{CapacityDifferenceCover.lean}{IsSharedCode.difference_cover_real_bound, IsSharedCode.square_cover_bound, IsSharedCode.saturated_cover_bound}
For positive integers $u,v$, set
$R=\min(uv,d-1)$ and $U=\min(d,M\lfloor d/(R+1)\rfloor)$.
Then
\begin{equation}\label{eq:many-classes}
 D\le N-\frac{bd(d-U)}{u+v}.
\end{equation}
If $M\le d/2$ and $t=\lceil\sqrt{2M}\rceil$, then
\begin{equation}\label{eq:sqrt-classes}
 D\le\left(1-\frac1{4t}\right)N.
\end{equation}
For every $d\ge2$ and every $M$, writing $t_d=\lceil\sqrt{d-1}\rceil$,
$D\le N-bd(d-M)/(2t_d)$.
\end{theorem}
\begin{outline}
Use the shift set
$E=\{0,1,\ldots,u-1\}\cup\{u,2u,\ldots,vu\}$, reduced modulo $d$.
Its size is at most $u+v$. Every integer $1\le h\le uv$ equals $ju-i$
with $j=\lceil h/u\rceil\le v$ and $0\le i<u$. Thus $E-E$ contains
every signed difference through $R$.

For one class, list its vertices cyclically. The positive gaps sum to
$d$, including the sole gap $d$ for a singleton. Keep consecutive edges
whose gap is at most $R$. If there are $\ell\ge1$ larger gaps, there
are at most $\ell$ components and $\ell(R+1)\le d$. If there is no
large gap, there is one component, also at most $\lfloor d/(R+1)\rfloor$
because $R\le d-1$. Hence $\chi(E)\le U$.
The translate bound gives $D\le N-bd(d-\chi(E))/|E|$.
Both numerators are nonnegative; replacing $\chi(E)$ by $U$ and
$|E|$ by its upper bound proves \eqref{eq:many-classes}.

For $u=v=t$, either $t^2\ge d-1$, giving $U=M\le d/2$, or
$U\le Md/(t^2+1)<d/2$. Equation \eqref{eq:sqrt-classes} follows.
For $u=v=t_d$ the radius is $d-1$ and $U=M$, giving the last claim.
Reduction of the shift set may create repetitions; the proof uses only
its upper cardinality bound. It includes $d=1$ in the first inequality.
\end{outline}

This application uses the elementary sparse difference cover above, not a
new difference-cover theorem. It changes the general deficit order from
$1/M$ to $1/\sqrt M$, while stronger fixed-$M$ bounds remain available.
For example, if $D\ge(1-\eps)N$ and $0<\eps\le1/8$, then
\eqref{eq:sqrt-classes} gives the precise dichotomy
\[
 M>d/2\quad\text{or}\quad M>1/(128\eps^2).
\]
Indeed $\lceil\sqrt{2M}\rceil\ge1/(4\eps)$ forces
$\sqrt{2M}>1/(4\eps)-1\ge1/(8\eps)$.
The saturation at $d$ is necessary because singleton classes can decode
the full source. Nothing here extracts these decoder interfaces from a
native tape execution or closes them under composition.

\section{Completion and exact interval capacity}\label{sec:structure}

\subsection{A normal form without loss of dimension}

Choose a representative $r_J$ from each class $J$. Define
\begin{align*}
 (PX)_c&=\sigma^{c-r_J}\operatorname{row}_{r_J}X\quad(c\in J),\\
 H_*&=\operatorname{im}P,
 &W&=\ker P=\{X:\operatorname{row}_{r_J}X=0\text{ for every }J\}.
\end{align*}
Here $P$ acts on rows, and each row has $bd$ scalar coordinates.
The class-covariant space $H_*$ was already used in the small-code
constructions of \cite{SeventeenthPaper2026}. What matters here is that
one may require its inclusion when maximizing capacity.

\begin{theorem}[No-loss completion]\label{thm:completion}
\formal{CapacityCompletion.lean}{IsSharedCode.completedCode, finrank_completedCode_add_rank, finrank_le_completedCode, completedCode_eq_of_range_le}
The map $P$ is an idempotent linear map commuting with every $T^t$, and
$\dim H_*=Mbd$. If $C$ is a linearly supported code with deterministic
shared decoders, then
\[
 C^+=(C\cap W)\oplus H_*
\]
has deterministic shared decoders for the same partition and
\[
 \dim C^+=\dim C-\rank(P|_C)+Mbd\ge\dim C.
\]
If $H_*\subseteq C$, then $C^+=C$ as subspaces.
\end{theorem}
\begin{outline}
The representative rows are unchanged by $P$, freely parametrize its
range, and vanish exactly on its kernel. These facts prove idempotence,
the range dimension, and the direct decomposition. Row rotations commute,
so $PT^t=T^tP$. Given original decoders $f_J$, use
\begin{equation}\label{eq:completed-decoder}
 f_J^+(Y)=f_J((I-P)Y)+\sigma^{-r_J}\operatorname{row}_{r_J}Y.
\end{equation}
For $X_0\in C\cap W$, $H\in H_*$ and $c\in J$, the first argument
in \eqref{eq:completed-decoder} at $T^c(X_0+H)$ is $T^cX_0$.
The second term is $\operatorname{row}_cH$, because the representative
rows of $X_0$ vanish. Thus the output is the demanded row of $X_0+H$.
Rank-nullity for $P|_C$ gives the exact dimension. If $H_*\subseteq C$,
then $PX\in C$ for $X\in C$, giving $C=(C\cap W)\oplus H_*$.
\end{outline}

The completion need not contain the original code: it replaces it, and
may change every decoder away from required views. The residual is
$C\cap\ker P$, not $(I-P)C$. If the original decoders are linear,
\eqref{eq:completed-decoder} is linear as well.

\begin{corollary}[Residual capacity and transversal rank]\label{cor:residual-capacity}
\formal{CapacityOptimization.lean}{shared_capacity_eq_residual, linear_capacity_eq_residual, capacity_transversal_deficiency, optimal_transversal_full_rank}
In either the deterministic-decoder or linear-decoder model, let
$\operatorname{Cap}(p)$ be the maximum code dimension, and let
$\operatorname{Cap}_W(p)$ be the maximum among codes contained in $W$.
Then
\[
 \operatorname{Cap}(p)=Mbd+\operatorname{Cap}_W(p).
\]
For every code $C$ and every representative transversal,
\[
 \operatorname{Cap}(p)-\dim C\ge Mbd-\rank(P|_C).
\]
In particular every dimension-optimal code projects onto the entire tuple
of representative rows, for every such transversal.
\end{corollary}
\begin{outline}
Attainable dimensions form nonempty subsets of $\{0,\ldots,N\}$, so
the maxima exist. For any $C$, rank-nullity gives
$\dim C\le Mbd+\dim(C\cap W)\le Mbd+\operatorname{Cap}_W(p)$.
Conversely, complete a residual maximizer. Completion preserves linearity
when required, so both models obey the same equality. Applying maximality
to the exact dimension of $C^+$ proves the displayed inequality. Row
evaluation on the transversal has the same rank as $P|_C$; at a maximum
that rank must be $Mbd$.
\end{outline}

\subsection{Normal bases and disjoint residual views}

We describe a cyclic-interval partition by a transversal containing every
class boundary: there is one representative $r_J\in J$ per class, and
every $c$ with $p(c-1)\ne p(c)$ is one of these representatives.
For a single class any representative may be used. This formulation
includes intervals that cross the chosen coordinate origin.

\begin{theorem}[Extremal interval codes]\label{thm:interval-attainment}
\formal{CapacityIntervals.lean}{CapacityIntervals.exists_binary_code_dimension, CapacityIntervals.exists_optimal_interval_code, IsSharedCode.two_mul_finrank_le_interval}
Let a partition of $\Z/d$ have $M$ nonempty classes and let $d-M$ be even.
There is a binary linear source of dimension
\[
 \frac{d^2+Md}{2}
\]
with deterministic shared decoding for this partition. For cyclic-interval
classes this dimension is optimal. Repeating the construction in $b$ lanes
gives the corresponding dimension $b(d^2+Md)/2$ for every $b\ge1$.
\end{theorem}
\begin{outline}
Let $E=\F_{2^d}$ and $\Phi(x)=x^2$. By the normal-basis theorem
\cite[Theorem~5.18]{Milne2022}, choose
$\alpha$ whose conjugates form a binary basis, and use its reversed order:
$e_j=\Phi^{-j}\alpha$. The identification
$\iota(x)=\sum_jx_je_j$ satisfies
\[
 \iota(\sigma x)=\Phi(\iota(x)).
\]
Indeed replacing $j$ by $j+1$ in the sum proves this identity with the
stated rotation convention. Put $\beta=\alpha$. Its $d$ Frobenius
conjugates are distinct, since the basis elements are distinct.

Write $d-M=2s$ and pair the nonrepresentative rows as $(a_i,b_i)$.
Define a residual source $C_0$ by assigning to each pair the rows
$\iota^{-1}(x_i),\iota^{-1}(\beta x_i)$, with independent $x_i\in E$;
set every representative row to zero. Then $\dim C_0=sd$. In $T^cC_0$
the same pair has the form $(u_i,\Phi^c(\beta)u_i)$. If a matrix belongs
to both $T^cC_0$ and $T^eC_0$, $c\ne e$, equality of the two rows gives
$(\Phi^c\beta-\Phi^e\beta)u_i=0$ for every $i$. Thus all $u_i=0$.

Consequently the nonzero translated views of $C_0$ are disjoint. Define
one total function $f$ on full matrices by $f(0)=0$ and
$f(Y)=\operatorname{row}_c(T^{-c}Y)$ if $0\ne Y\in T^cC_0$; set it to
zero outside these views. The phase $c$ is unique, so this is well-defined
and decodes all residual receivers. Complete $C_0$ by Theorem
\ref{thm:completion}. Its dimension becomes $Md+sd=(d^2+Md)/2$.
This also covers $s=0$, when the residual is zero.

For an interval partition, the graph in Lemma \ref{lem:deterministic-tp}
with $E=\{0,1\}$ has one component per interval. Hence
$2\dim C\le d^2+Md$. The case $d=1$ is immediate. Independent lane
copies give the asserted all-width construction, while the same
deterministic translate bound gives its all-width interval upper bound.
\end{outline}

The normal-basis input is classical; it holds in even and composite
degrees too. No primitive normal basis, root of unity, or diagonalization
of the binary rotation is needed. The decoder identifies a phase by
examining a full residual view. This is a construction in the shared-view
model, without a tape implementation of that access.

At $d=5$, classes $\{0\},\{1\},\{2,3,4\}$ have deterministic capacity
$20$. The comparison with the reported linear-decoder maximum $19$
in \cite{SeventeenthPaper2026} retains that source's solver qualification.
The theorem proves the deterministic side and its upper bound, not a new
formal proof of the reported linear nonexistence.

\subsection{Why the hypotheses stay separate}

\begin{proposition}[Deterministic decoding need not linearize]\label{prop:strict-decoder}
\formal{CapacityControls.lean}{CapacityControls.exists_deterministic_not_linear}
There is a one-class binary linear source with deterministic shared
decoding and no globally linear shared decoder.
\end{proposition}
\begin{outline}
At $d=3$ let $X$ have row zero $(1,1,0)$ and its other rows zero, and take
$C=\operatorname{span}_{\F_2}\{X\}$. The three nonzero views
$X,TX,T^2X$ are distinct, so a deterministic map may send them respectively
to $(1,1,0),0,0$, and send zero to zero. Those are exactly the demanded
rows. But $X+TX+T^2X=0$, while the required outputs have nonzero sum.
No linear function has these values.
\end{outline}

\begin{proposition}[A high-dimensional affine marker]\label{prop:affine-marker}
\formal{CapacityMarker.lean}{CapacityMarker.exists_affine_source_decoder}
Over every field and for every $d\ge1$, a one-class deterministic shared
decoder transposes an affine source of dimension $d^2-d$.
\end{proposition}
\begin{outline}
Require row zero to be the coordinate word $e_0$ and leave all other
entries free. The first row of $T^cX$ is $\sigma^ce_0$, which identifies
$c$ uniquely. A total decoder reads that marker, applies $T^{-c}$, and
returns row $c$; define it arbitrarily on other full views. This proves
correctness on all matrices in the affine source and its dimension.
For $d>3$ this exceeds the one-class linear-source bound
$(d^2+d)/2$, so a linear-source hypothesis cannot be replaced by an
affine-dimension assertion for arbitrary deterministic decoders.
\end{outline}

Both controls concern complete shared views. The marker is not an
unpaid row-addressing operation for a tape machine, and the strictness
example does not compare the maximum capacities of two decoder models.

\section{Cyclic lifts with the same classes}\label{sec:constructions}

The original blow-up construction in \cite{ClassUniformMemo2026} gives
unbounded absolute excess over the interval bound with a growing number
of classes. Its revised memo also gives a two-class parity-coset family
for every $d=4k\ge8$ over fields where $2\ne0$
\cite[Theorem~4.8]{ClassUniformMemo2026}. The revised seventeenth paper
extends that rule to every field, including characteristic two, and uses
mixed blow-ups to obtain unbounded excess for every fixed number of
classes\cite{SeventeenthPaper2026}.
The construction below instead gives a lift for an arbitrary seed with
the same number of classes, an exact periodic-capacity formula, and a
paid realization of its exhibited binary rule. These fixed-class families
all have vanishing normalized excess; fixed-class or binary additive
excess alone is not a new conclusion here.

\begin{theorem}[Deterministic cyclic lift]\label{thm:cyclic-lift}
\formal{ConstructionLinearLift.lean}{cyclic_lift}
Let a linear source $C_s\le(F^b)^{s\times s}$ have dimension $D$ and
deterministic shared decoders for an $M$-class partition $p$ of $\Z/s$.
Let $m\ge1$ be coprime to $s$, with $m\ne0$ in $F$, and put $d=ms$.
The lifted partition $p^\uparrow(c)=p(c\bmod s)$ admits a linear source
of dimension
\begin{equation}\label{eq:cyclic-lift-dimension}
 D^\uparrow=\frac{b(d^2-ms^2)}2+mD
\end{equation}
with deterministic shared decoders. Linear seed decoders give linear
lifted decoders.
\end{theorem}
\begin{outline}
Use Chinese-remainder coordinates $(a,r)\in\Z/m\times\Z/s$ for both
rows and columns. For $z\in\Z/m$ define
\[
 Z_z(X)_{r,c}=\frac1m\sum_{a\in\Z/m}X_{(a,r),(z-a,c)},\qquad
 (EX)_{(a,r),(h,c)}=Z_{a+h}(X)_{r,c}.
\]
The map $E$ is an idempotent. Its image consists of arrays whose outer
indices enter only through their sum, and is parametrized by the $m$
independent seed matrices $Z_z$. Thus $\rank E=mbs^2$.
It commutes with column rotation and transpose. On its image transpose
acts as the inner $s\times s$ transpose on each $Z_z$.

Take the direct sum
\[
 C^\uparrow=(\ker E\cap\operatorname{Fix}\tau)\oplus
 \{X\in\operatorname{im}E:Z_z(X)\in C_s\text{ for every }z\}.
\]
The second summand has dimension $mD$. The symmetric space has dimension
$bd(d+1)/2$ in every characteristic. Its image under $E$ is exactly the
space of $m$ symmetric seed matrices, of dimension $mbs(s+1)/2$:
commutation proves inclusion, and every such image vector is already
symmetric and fixed by $E$. Rank-nullity gives the first summand's
dimension $b(d^2-ms^2)/2$.

For an arbitrary full view $Y$, define one decoder per original class by
\begin{equation}\label{eq:lifted-deterministic-decoder}
 F_J(Y)_{a,r}=(Y-EY)_{(a,r),(0,0)}+f_J(Z_a(Y))_r.
\end{equation}
At a receiver with CRT index $(u,v)$, literal rotation satisfies
$Z_a(T^{(u,v)}X)=T_s^vZ_{a+u}(X)$. Thus the complete induced map is
\[
 K^\uparrow(X)_{(a,r),(u,v)}
 =(X-EX)_{(a,r),(u,v)}+f_{p(v)}(T_s^vZ_{a+u}(X))_r.
\]
For a codeword, the first term is its symmetric part, and the second
transposes its projected part. This proves correctness. The receiver
index is not an extra argument to \eqref{eq:lifted-deterministic-decoder};
only its seed class selects the function. In particular a nonlinear
decoder is applied after averaging the seed matrix, never distributed
through that average. All displayed operations are linear when the
seed decoder is linear.
\end{outline}

For linear rules there is also an operator description. On the outer
tensor square put
$E=m^{-1}\sum_{t\in\Z/m}G^t\otimes G^{-t}$, where $G$ is rotation.
After regrouping outer and inner factors, the complete lifted rule is
\begin{equation}\label{eq:lift-operator}
 K^\uparrow=(I-E)\otimes I+E\otimes K_s.
\end{equation}
If a seed term has column factor $P_JS^e$, its lifted column factor is
$G^{-t}\otimes P_JS^e=P_{J^\uparrow}\sigma_d^u$, where CRT gives
$u\equiv-t\pmod m$ and $u\equiv e\pmod s$. Hence this is one allowed
same-class operator, not a composition with a free coordinate permutation.
If the seed code is the entire agreement kernel, the displayed lifted
code is the entire kernel of \eqref{eq:lift-operator}.

\begin{corollary}[A fixed two-class excess]\label{cor:fixed-two-excess}
\formal{ConstructionSeed.lean}{binaryCyclicCode_linear, binaryCyclicPartition_surjective, finrank_binaryCyclicCode, binaryCyclicCode_interval_excess}
For every $m\ge1$ coprime to six, $d=6m$ admits a binary two-class
linear rule and a linear source of dimension
\[
 \frac{d^2}{2}+\frac{7d}{6}.
\]
Its excess over the two-interval bound is $d/6$.
\end{corollary}
\begin{outline}
Use the published six-dimensional seed with classes
$\{0,1,3\},\{2,4,5\}$ and its verified $25$-dimensional code
from \cite{SeventeenthPaper2026,ClassUniformMemo2026}.
Denote the complete rule induced by these seed decoders by $K_6$.
Since $m$ is odd, it is invertible in $\F_2$.
Equation \eqref{eq:cyclic-lift-dimension} gives
$18m^2+7m=d^2/2+7d/6$; subtract $d^2/2+d$.
The conclusion needs a seed subcode of dimension $25$, not a proof that
$25$ is the seed's maximum capacity.
\end{outline}

\subsection{The exact linear capacity of a periodic lift}

Write $\operatorname{Cap}^{\mathrm{lin}}_F(p)$ for capacity with linear
decoders over $F$.

\begin{theorem}[Periodic capacity reduction]\label{thm:periodic-capacity}
\formal{ConstructionCapacity.lean}{periodic_linear_capacity_split, periodic_linear_capacity_lower, periodic_linear_capacity_algebraicClosure_upper}
In the scalar case $b=1$, suppose additionally that $F$ contains all
$m$-th roots of unity. Then
\[
 \operatorname{Cap}^{\mathrm{lin}}_F(p^\uparrow)
 =\frac{d^2-ms^2}{2}+m\operatorname{Cap}^{\mathrm{lin}}_F(p).
\]
Without the splitting assumption, the same expression with $F$ is a lower
bound, and the expression with its algebraic closure $\overline F$ is an
upper bound.
\end{theorem}
\begin{outline}
The lower bound is Theorem \ref{thm:cyclic-lift}. For the upper bound,
the allowed column space in CRT coordinates is
$\operatorname{Circ}_m\otimes\operatorname{span}\{P_JS^e\}$.
Simultaneously diagonalize the outer circulants, using the same basis
on both tensor factors so that the demand remains swap. This is valid
because $m$ is invertible and its roots split. Write the input as an
$m\times m$ array of $s\times s$ blocks $X_{ij}$. Correctness has the form
\[
 X_{ji}^{\mathsf T}=\sum_a L_{ia}^j(X_{aj}),
\]
where each $L_{ia}^j$ is an allowed seed rule. In particular
$\delta_j=\dim\ker(\tau_s-L_{jj}^j)\le
\operatorname{Cap}^{\mathrm{lin}}_F(p)$.

Supply all upper-triangular blocks $X_{ij}$, $i<j$, and process $j$ in
decreasing order. All off-diagonal blocks of column $j$ are known: those
above were supplied, those below belong to recovered rows. The equation
with $i=j$ determines $X_{jj}$ up to a vector in a fixed $\delta_j$-
dimensional kernel. Supply $\delta_j$ linear coordinates injective on
that kernel; this determines $X_{jj}$ uniquely even on an affine fibre.
The other equations now recover row $j$. This is an injective linear
description of size $s^2m(m-1)/2+\sum_j\delta_j$, proving the bound.
For nonsplitting $F$, embed a basis of a maximizing source into
$\overline F$ and take its $\overline F$-span. Scalar extension preserves
the basis's independence and extends every class rule by its coordinate
coefficients. The extended source has the same dimension and the same
receiver classes. Apply the split result over $\overline F$.
\end{outline}

Thus every rule on a fixed periodic seed partition has rate
$1/2+O(1/d)$, even if it is not the exhibited lift. The qualitative limit
already follows from the period obstruction in \cite{SeventeenthPaper2026};
the equality records the exact finite seed-capacity contribution.

\subsection{A paid implementation of the exhibited binary lift}

\begin{theorem}[Two-tape implementation]\label{thm:paid-lift}
\notformalized
There is a fixed finite-alphabet, finite-control machine with two initially
blank work tapes, read-only input and append-only unreadable output that
computes the binary operator in Corollary \ref{cor:fixed-two-excess}
lane by lane in $O(P)$ time, $P=d^2b$, on admitted dimensions $d=6m$,
$\gcd(m,6)=1$. There are separate machines for $b=1$ and for
$b=\max(1,\lceil\log_2d\rceil)$. Each has a transpose-correct subcode of
dimension $b(d^2/2+7d/6)$.
\end{theorem}
\begin{informalproof}
Here is a complete bounded-state schedule, including preparation. Put
$\mu=mb$. The first work tape has $m$ consecutive $b$-cell words and
$36$ binary payload tracks per physical cell. For $z\in\Z/m$ their
intended contents are
\[
 H_z[i,j]=\bigoplus_{r\equiv i\ (6),\ c\equiv j\ (6),\ r+c\equiv z\ (m)}X_{r,c},
 \qquad 0\le i,j<6.
\]
Each word has independent FIRST/LAST marker bits; HOME is at cell zero,
END at $\mu$, and START at one word's first cell. Coincident markers
when $b=1$ or $m=1$ are distinct alphabet tracks. All payload starts zero.
After preparation the second work head parks throughout both data passes.

At the start of input row $r$, START and the bank head are at word
$r\bmod m$, with finite registers $i=r\bmod6$, $j=0$, and lap count zero.
A DATA step reads the old raw bit and old bank symbol, advances both heads
right, and either toggles $H[i,j]$ by that bit (pass one) or appends
$x+H[i,j]+(K_6H)[i,j]$ (pass two), preserving the bank. The last expression
is a fixed Boolean function of $36$ old bits. An OLD LAST marker increments
$j$ modulo six and enters LAND. LAND at END rewinds physically to HOME;
else it is at a FIRST marker. Returning to START counts a completed lap.
Six laps finish one row. Importantly $j$ is not reset at a lap boundary:
$m$ need not be divisible by six. At row end clear START, physically advance
one whole word, wrapping at END if necessary, and set START there; then
increment $i$ modulo six and reset $j$ and the lap counter. Perform this
offset after the last row too. Only the row-start mode tests the raw
endmarker; housekeeping preserves payload, holds raw still, and emits
nothing.

Induction gives bank address $((r+c)\bmod m)b+t$ at lane $t$ of entry
$(r,c)$, with $j=c\bmod6$. Thus pass one computes exactly $H$. After
$d$ rows its START is HOME. Rewind raw across its entire input, preserving
the bank, and run the same schedule in emission mode. There are exactly
$P$ appends in ordinary row-major order. Because $m$ is odd, the sum in
$H$ is precisely the binary averaging projector $E$; the emitted formula
is \eqref{eq:lift-operator}. All reads are pretransition reads, and all
writes are at the old head locations.

For completeness one pass has $P$ DATA steps, $d^2$ ordinary LAND steps,
and six physical END wraps per row, costing $6d(\mu+1)$ extra steps.
Row offsets cost $db+d+6(\mu+1)$, and row-start branches cost $d$.
The total is
\[
 2P+d^2+2db+8d+6\le5P\qquad(d\ge6,b\ge1).
\]
The between-pass raw rewind costs $P+2$; the final absent transition halts.

It remains to establish this bank from blanks without size advice. Count
the raw length on the second tape using a little-endian binary counter
and a marked least-significant cell. A carry through $k$ ones and return
costs $O(k+1)$, whose total over $P$ increments is $O(P)$; hold raw fixed
during each carry, then advance it. Rewind raw. Fixed schoolbook arithmetic
on a constant number of delimited $O(\log(P+2))$-bit registers verifies
$P=d^2$ or $P=d^2\max(1,\lceil\log_2d\rceil)$ by binary search, and
checks $d=6m$, $\gcd(m,6)=1$. These functions are strictly increasing in
$d$, and all arithmetic takes $\operatorname{polylog}(P+2)$ time.
Reject other lengths before output.

The first head now writes the bank sequentially while the second head
maintains counters for $m$ words and $b$ cells. Binary decrement and return
has aggregate $O(b+\log(P+2))$ cost per word; no full-width zero test is
performed per cell. Setup, copying and register trips cost at most
$\operatorname{polylog}(P+2)$ per word. Install all markers, write END,
return to HOME, and erase the second tape's delimited scratch frame.
This costs
$O(P+m\log^c(P+2)+mb+\log^c(P+2))=O(P)$ for a fixed exponent $c$,
since $m=d/6$ and $d\log^c(d+2)=O(d^2b)$ in both width regimes.
Every unbounded counter was tape data, and all other registers range over
fixed finite sets. This establishes the promised initial boundary with
exactly two work heads and proves the uniform time bound.
\end{informalproof}

This is an implementation of one displayed operator on all inputs of its
admitted lengths. Its transpose-correct source has rate tending to one
half, so it is not a full-input transposer or a lower bound against native
machines. No general class-uniform operator is assigned a free implementation.

\part{Composition and nonlinear obstructions}
\section{Two linear stages}
\label{sec:composition}

A class bound for one stage is not closed under composition: even two
small-class maps can have many different endpoint rules.  We therefore
keep both complete operators.  Their exact composition forces a
\emph{common message} on a small rotation-invariant space.  Arithmetic
progressions then force this message space to grow.

\subsection{One message on several translates}

\begin{lemma}[Common-message growth]\label{lem:common-message}
\formal{CompositionProgression.lean}{commonMessage_finrank_le}
Let $A$ be a finite-dimensional vector space over any field, let
$R\colon A\to A$ be linear, and let $W\le A$ have dimension $D$.
Suppose that a linear map $L\colon A\to Z$ is injective on $W$ and
\[
 L(R^jx)=L(x)\qquad(x\in W,\ 0\le j<\ell),
\]
where $\ell\ge1$.  If $f=\dim\ker(R-I)$, then
\begin{equation}\label{eq:message-growth}
 \ell D\le \dim A+(\ell-1)f.
\end{equation}
Invertibility of $R$ is not required.
\end{lemma}
\begin{outline}
Put $V_t=W+RW+\cdots+R^{t-1}W$ and $V_0=0$.
For $t\le\ell$, the map $L$ vanishes on $(R-I)V_{t-1}$.
This image is disjoint from $W$, because $L$ is injective on $W$,
and their sum lies in $V_t$.  Rank--nullity therefore gives
\[
 \dim V_t\ge D+\dim V_{t-1}-f.
\]
Starting with $V_1=W$ and using $\dim V_\ell\le\dim A$ proves the
claim.  In fact telescoping gives
$V_t=W\oplus(R-I)V_{t-1}$, but the formal proof only needs inclusion.
\end{outline}

The correction $f$ matters.  Over $\F_2$, take a three-coordinate
rotation and $W=\operatorname{span}\{(1,1,1),e_0\}$.  Projection along
$\operatorname{span}\{(R-I)e_0\}$ identifies $W$ and $RW$ with the same
two-dimensional message.  Thus $\ell D=4=3+1$ for $\ell=2$; a bound
without the fixed-space term would be false.

For positive integers $M,\ell$, write $\mathsf W(M,\ell)$ for any
finite interval length such that every $M$-coloring of
$\{0,\ldots,\mathsf W(M,\ell)-1\}$ has a monochromatic ordinary
$\ell$-term arithmetic progression with positive step.
Existence is the finite van der Waerden theorem.  Our formal library
derives this version from Mathlib's finite Hales--Jewett theorem
\cite{HalesJewett1963}:
summing the symbols of a combinatorial line gives an ordinary progression
in a fixed bounded interval
(\lean{exists_progression_bound}).  We use no numerical estimate for
$\mathsf W$.

\begin{theorem}[Vanishing common-message rate]\label{thm:message-rate}
\formal{CompositionChannels.lean}{commonMessage_interval_bound}
Let $A=(F^q)^{\Z/d}$, with simultaneous cyclic rotation $T$, where
$q,d\ge1$ and $F$ is any field.  Let $W\le A$, and let
$B\colon W\to Z$ be linear and injective.  Suppose the $d$ positions
have $M$ nonempty classes and each class $J$ has a linear map
$L_J\colon A\to Z$ such that
\[
 L_J(T^cx)=B(x)\qquad(c\in J,\ x\in W).
\]
For $\ell\ge2$, put $N_0=\mathsf W(M,\ell)$.  If $d\ge N_0$, then
\begin{equation}\label{eq:message-rate}
 \dim W\le \frac{q(d+N_0-1)}{\ell}.
\end{equation}
\end{theorem}
\begin{outline}
Choose a monochromatic progression
$c_0,c_0+s,\ldots,c_0+(\ell-1)s$ among the first $N_0$ positions.
Here $s\ge1$ and $(\ell-1)s\le N_0-1$.
Apply \cref{lem:common-message} with $R=T^s$ and
$L=L_JT^{c_0}$.  A word fixed by $T^s$ is determined by its first
$s$ coordinates, by repeatedly reducing a natural-number index by $s$.
Thus the fixed space in $A$ has dimension at most $qs$ over every field.
No diagonalization is used.
Thus
\[
 \ell\dim W\le qd+(\ell-1)qs
 \le q(d+N_0-1).
\]
\end{outline}

For fixed $M$, this implies $\dim W/(qd)\to0$ uniformly in $F$, $q$,
the partition and the linear maps: choose $\ell$ first, then take
$d\ge\mathsf W(M,\ell)$, so the ratio is less than $2/\ell$.

Linearity here concerns the decoder, not merely the source.  To see the
distinction, identify binary words of length $d$ with $E=\F_{2^d}$ in a
normal basis, so rotation is Frobenius $\Phi$ \cite{Milne2022}.  Choose $\beta$ with
$d$ distinct conjugates.  The binary-linear two-channel space
$W=\{(x,\beta x):x\in E\}$ has dimension $d$, and its distinct rotated
views intersect only at zero.  A single total nonlinear decoder sends
$(u,\Phi^c(\beta)u)$, for $u\ne0$, to $\Phi^{-c}u$ and sends all other
points to zero.  It recovers the same original $x$ at every view with
rate $1/2$.  This is a source-level counterexample to removing decoder
linearity; it is not a two-stage transpose factorization.

\subsection{The complete two-stage reduction}

We now use the matrix model with $b$ coordinates per entry, allowing
arbitrary mixing of rows and lanes by every linear decoder.  Both stages
use the same column rotation.  Their partitions need not agree.

\begin{theorem}[Two linear stages]\label{thm:two-linear-stages}
\formal{CompositionLinear.lean}{twoHop_progression_gcd,twoHop_coprime_progression_dimension_le,twoHop_dimension_le,exists_twoHop_cutoff}
Let $K_1,K_2$ be linear class-uniform maps on $F^{d\times d\times b}$,
with $M_1,M_2$ nonempty classes, and suppose $K_2K_1=\tau$ on the
whole space.  If a first-stage class contains
$c_0,c_0+s,\ldots,c_0+(\ell-1)s$, with $\ell\ge1$, integer $s\ge0$
and indices read modulo $d$, then
\begin{equation}\label{eq:two-stage-ap}
 \ell d\le M_2d+(\ell-1)M_2\gcd(d,s).
\end{equation}
In particular, a progression of length $M_2+1$ with $\gcd(d,s)=1$
forces $d\le M_2^2$. Without the coprimality premise,
\begin{equation}\label{eq:two-stage-cutoff}
 d\le M_2\bigl(\mathsf W(M_1,M_2+1)-1\bigr).
\end{equation}
Thus fixed bounds on both class counts exclude exact two-stage
transposition at all sufficiently large $d$, over every field and
uniformly over every positive width $b$.
\end{theorem}
\begin{outline}
Both square factors are invertible, and $K_1^{-1}=\tau K_2$.
Let $S_0=\operatorname{Fix}(T)$, of dimension $bd$, and let $P_2\le F^d$
be the span of the second-class indicator vectors.  Their disjoint
nonempty supports give $\dim P_2=M_2$ over every field.
For $X\in S_0$, all views $T^cX$ coincide.  Hence $K_2X$ is constant
on each second receiver class.  It follows that
\[
 W:=K_1^{-1}S_0\ \le\
 A:=P_{2,\mathrm{row}}\otimes F^d_{\mathrm{column}}\otimes F^b,
 \qquad \dim W=bd,
 \qquad \dim A=bM_2d.
\]
This is an actual consequence of the global product identity, not an
independently chosen source.  The ambient $A$ is $T$-invariant and consists
of $bM_2$ cyclic channels.

The common column of $K_1X$, for $X\in W$, defines one injective linear
message $B(X)$.  At every receiver $c$ in a first class $J$,
$f_{1,J}(T^cX)=B(X)$.  Apply \cref{lem:common-message} on $A$ to the
specified progression. A word fixed by the $s$-step rotation is constant
on its $\gcd(d,s)$ position orbits. Restriction to one representative per
orbit is injective, so the fixed-space term is at most
$bM_2\gcd(d,s)$. The formal proof uses the actual translation-orbit quotient.
Canceling $b>0$ gives \cref{eq:two-stage-ap}; putting $\ell=M_2+1$ and
$\gcd(d,s)=1$ gives $d\le M_2^2$.

Put $N_0=\mathsf W(M_1,M_2+1)$.  If $d>M_2(N_0-1)$, then
$d\ge N_0$.  A monochromatic $(M_2+1)$-term progression in the first
$N_0$ positions has positive step $s$ with $M_2s\le N_0-1$.
Since $\gcd(d,s)\le s$ for this positive step,
\eqref{eq:two-stage-ap} would imply
\[
 d\le M_2^2s\le M_2(N_0-1),
\]
a contradiction.  This proves the finite cutoff.  Using fixed upper
bounds for the numbers of colors gives one cutoff for all the stated
parameters.
\end{outline}

The orbit correction can exclude particular masks before the general
Ramsey threshold applies. For example, with two second-stage classes,
a monochromatic three-term progression with step coprime to $d$ already
forces $d\le4$. The global two-class cutoff below remains $16$, since
Ramsey extraction does not require its progression step to be coprime.

The theorem concerns two complete maps on the same state space and
tensor side.  It neither multiplies endpoint class counts nor excludes
additional uniform postprocessing.  The familiar two partial-digit
swaps at $d=M_1M_2$ remain valid controls.  The bound does not prove
that their product condition is necessary for arbitrary coded factors.

\begin{corollary}[Two classes at each stage]\label{cor:two-binary-classes}
\formal{CompositionLinear.lean}{twoHop_two_classes_le_sixteen}
If each of the two linear stages uses at most two classes, exact
transposition requires $d\le16$, over every field and every positive width.
\end{corollary}
\begin{outline}
The finite statement $\mathsf W(2,3)\le9$ is checked by Lean's kernel
over all colorings of nine positions.  Its declaration is
\lean{binary_progression_nine_finite}; no minimality assertion is needed.
Substitute this bound in \cref{eq:two-stage-cutoff}.
\end{outline}

\subsection{One global rank error}

The preceding witness has dimension only $bd$.  To obtain more than a
row-scale obstruction, retain a whole band of frequencies and charge the
same global error only once.

\begin{lemma}[Recurrence growth]\label{lem:recurrence-growth}
\formal{CompositionRecurrence.lean}{recurrence_finrank_le}
Let $R\colon A\to A$ and $L\colon A\to Z$ be linear on a
finite-dimensional space, let $W\le A$ have dimension $D$, and let
$k,m\ge1$.  Let $p(z)=\sum_{i=0}^k a_i z^i$.
Suppose $(L,LR,\ldots,LR^{k-1})$ is injective on $W$ and
\[
 LR^jp(R)w=0\qquad(w\in W,\ j\ge0,\ j+2\le mk).
\]
Then, with $f=\dim\ker p(R)$,
\begin{equation}\label{eq:recurrence-growth}
 mD\le\dim A+(m-1)f.
\end{equation}
\end{lemma}
\begin{outline}
Set $V_t=\sum_{0\le j<t}R^jW$ and $G=(L,\ldots,LR^{k-1})$.
For $0\le t\le(m-1)k$, every component of $G$ vanishes on
$p(R)V_t$: its generator exponent is at most $t+k-2\le mk-2$.
Consequently $W\cap p(R)V_t=0$.
The degree bound gives $W+p(R)V_t\subseteq V_{t+k}$, so
\[
 \dim V_{t+k}\ge D+\dim V_t-f.
\]
Start with $\dim V_k\ge D$ and use $t=k,2k,\ldots,(m-1)k$.
Then $\dim V_{mk}\ge mD-(m-1)f$, proving the result.  When $m=k=1$
the displayed annihilation range is empty, as required.
\end{outline}

\begin{theorem}[Growing two-stage rank deficit]\label{thm:two-stage-rank}
\formal{CompositionRankField.lean}{twoHop_rank_interval_any_field,exists_twoHop_rank_threshold,twoHop_rank_superlinear}
Let the characteristic of $F$ not divide $d$, and let $K_1,K_2$ be
linear shared-decoder maps with $M_1,M_2$ classes on
$F^{d\times d\times b}$.  They need not be invertible.  Put
$r=\rank(\tau-K_2K_1)$.
For $k\ge1$, set
\[
 \ell=(2M_2+1)k-1,\qquad N_0=\mathsf W(M_1,\ell).
\]
If $d\ge N_0$, then
\begin{equation}\label{eq:rank-deficit}
 r\ge \frac{bk(d-N_0+1)}2.
\end{equation}
For fixed class bounds, $r/(bd)\to\infty$ uniformly in the field,
width, partitions and maps within this characteristic regime.
\end{theorem}
\begin{outline}
Extend scalars to a cyclotomic field with a primitive $d$th root
$\zeta$.  The complete maps and their receiver classes extend, and the
error rank does not increase: its extension factors through the scalar
extension of its original range.  It therefore suffices to prove the bound
there.  The characteristic hypothesis makes the roots distinct.  Let $C=\ker(\tau-K_2K_1)$ and $Y=K_1C$.
The restriction of $K_1$ to $C$ is injective, since $K_1x=0$ there
implies $\tau x=0$.  Thus $\dim Y\ge bd^2-r$, with $r$ denoting the
original error rank.

Let $S$ be the band of $k$ consecutive frequencies, consisting of
\[
 y_{:,c,:}=\sum_{\nu=0}^{k-1}\zeta^{\nu c}a_\nu.
\]
Its dimension is $bdk$, because $d\ge N_0\ge\ell\ge k$.
Pull $Y\cap S$ back through the injective restriction of $K_1$, obtaining
$W\le C$ of dimension $D\ge bdk-r$.  This intersection charges the
single global error $r$ once.

Let $H\le E^d$ be spanned by
$c\mapsto1_{p_2(c)=J}\zeta^{\nu c}$, over second classes $J$ and
$0\le\nu<k$, and put $h=\dim H\le M_2k$.
On $S$, the second operator produces columns of the form
$\sum_\nu\zeta^{\nu c}a_{p_2(c),\nu}$.
Since $w=\tau K_2K_1w$ on $W$, we have
\[
 W\le A:=H_{\mathrm{row}}\otimes E^d_{\mathrm{column}}\otimes E^b,
 \qquad \dim A=bhd.
\]
This ambient is invariant under column rotation.

Choose a monochromatic $\ell$-term progression with step $s\ge1$
among the first $N_0$ positions.  On $A$, put
$R=T^s$ and $L=f_{1,J}T^{c_0}$.  For $w\in W$,
\[
 LR^jw=\sum_{\nu=0}^{k-1}\zeta^{\nu c_0}
             (\zeta^{\nu s})^j a_\nu\qquad(0\le j<\ell).
\]
The numbers $\lambda_\nu=\zeta^{\nu s}$ are distinct:
for $\nu<\nu'$,
$0<(\nu'-\nu)s\le(k-1)s<d$.
The first $k$ displayed observations recover the $a_\nu$ by a
Vandermonde inverse, so their stacked map is injective on $W$.
The polynomial $p(z)=\prod_{\nu<k}(z-\lambda_\nu)$ annihilates these
sequences wherever all $k+1$ successive observations occur.  Taking
$m=2M_2$, the available exponents are exactly $j+2\le mk$.

The monic recurrence $p(R)x=0$ determines each cyclic channel from
its first $ks$ positions: its leading term recovers position $n$ from
positions $n-ks,n-(k-1)s,\ldots,n-s$.  Thus
$\dim\ker p(R)\le bhks$ on $A$.  This coordinate proof works without
diagonalizing the cyclic operator.
Apply \cref{lem:recurrence-growth} and $D\ge bdk-r$:
\[
 m(bdk-r)\le bhd+(m-1)bhks
 \le bM_2kd+(m-1)bM_2k^2s.
\]
Consequently
\[
 r\ge\frac{bkd}2-\frac{b(2M_2-1)k^2s}2.
\]
Now $(\ell-1)s\le N_0-1$ and
$(2M_2-1)k\le\ell-1$.  The subtracted term is at most
$bk(N_0-1)/2$, proving \cref{eq:rank-deficit}.

For each fixed $k$, once $d\ge\max\{N_0,2(N_0-1)\}$ the bound is
$r\ge bkd/4$.  Choosing $k$ arbitrarily large before $d$ proves the
uniform divergence of $r/(bd)$.
\end{outline}

This excludes repairs of uniformly $O(bd)$ rank, including all odd
binary dimensions.  It does not give $r\ge\epsilon bd^2$ for any fixed
$\epsilon>0$.  The arithmetic-progression cutoff grows with $k$, and
the proof does not permit $k$ to be a fixed positive fraction of $d$.
Nor is rank automatically a physical time charge.  Additional stages,
nonlinear factors and paid native realization require separate arguments.

\section{Complete nonlinear maps: symmetry, balance and periods}
\label{sec:nonlinear}

The arguments in this section use exactness on the complete finite state
space.  They do not linearize a decoder on a correctness source.  A complete
bijection has balanced column outputs; an arbitrary function agreeing with
transpose on a subspace need not have them.  This distinction supplies an
additional obstruction beyond the source-dimension bounds of Part I.

\subsection{One nonuniform layer and periodic first maps}

\begin{theorem}[An equivariant sandwich]\label{thm:uniform-sandwich}
\formal{UniformSandwich.lean}{uniform_sandwich_injective_classes,uniform_sandwich_finite}
Let the entry alphabet have at least two elements.  Suppose $U,V$ commute
with every column rotation, $U$ is injective, $K$ is class-uniform, and
$UKV=\tau$ on the complete state space.  Then every receiver class of $K$
is a singleton.  If the alphabet is finite, the injectivity of $U$ follows
from the factorization and need not be assumed.
\end{theorem}
\begin{outline}
Put $S=\operatorname{Fix}(T)$ and $W=KV(S)$.  The set $\tau S$ consists
of matrices with all rows equal, so it is invariant under column rotation.
Since $UW=\tau S$, equivariance and injectivity of $U$ imply that $W$ is
column-rotation invariant.  On each input in $V(S)\subseteq S$, the columns
of $K$ agree at any two receivers $c,e$ in the same class.  Applying this
fact to every rotation of every element of $W$ shows that $T^{e-c}$ fixes
all of $W$.  Equivariance of $U$ would then make $T^{e-c}$ fix all of
$\tau S$.  An arbitrary repeated row word disproves this when $c\ne e$.
For a finite alphabet, $U$ is surjective because the composite is, and
therefore injective.
\end{outline}

The following quantitative class count concerns an actual periodic first
map, and permits that map to be nonlinear.

\begin{theorem}[A periodic first factor]\label{thm:periodic-first}
\formal{SubgroupClasses.lean}{covariance_no_class_collision,orbit_classes_injective,exact_periodic_first_class_bound,exact_periodic_first_class_size_bound}
Let $A$ be a finite alphabet with at least two elements and let
$\tau=KF$ on $A^{d\times d}$, where $F$ is any total map commuting
with $T^r$ and $K$ is class-uniform with $M$ classes.  Put
$g=\gcd(d,r)$.  Every class meets each orbit of addition by $r$ in at
most one receiver.  Consequently each class has at most $g$ receivers and
\begin{equation}\label{eq:periodic-class-bound}
 M\ge d/g.
\end{equation}
The conclusion includes $r=0$, when it says only $M\ge1$.
\end{theorem}
\begin{outline}
Both factors are bijective.  From $FT^r=T^rF$ and $KF=\tau$ we obtain
$KT^r=R^rK$, hence the same covariance for every $s\in\langle r\rangle$.
If $e=c+s$ lies in the same class as $c$, the shared rule and this covariance
give, for every input $X$,
\[
 \operatorname{col}_eK(X)
 =\operatorname{col}_cK(T^sX)
 =\sigma^s\operatorname{col}_cK(X).
\]
Surjectivity of $K$ rules out this fixed relation between two distinct
output columns.  The $g$ cyclic orbits give the asserted class-size and
class-count bounds.  The formal proof derives the numerical bound through
the order of $r$ in $\Z/d\Z$.
\end{outline}

The bound is sharp for every divisor $g\mid d$.  Here is the explicit
partial-digit-swap construction, a reuse of the composition examples in
\cite{SeventeenthPaper2026}.  Write $d=gh$, $i=gw+z$ and $c=gu+v$, with
$0\le z,v<g$ and $0\le w,u<h$.  Define
\[
 (FX)_{gw+z,gu+v}=X_{gw+v,gu+z},\qquad
 (KY)_{gw+z,gu+v}=Y_{gu+z,gw+v}.
\]
Both are coordinate permutations, $F$ commutes with $T^g$, and direct
substitution gives $KF=\tau$.  The second map has the $h$ interval classes
$p(gu+v)=u$, with
\[
 f_u(Z)_{gw+z}=Z_{gu+z,\,g(w-u)\bmod d}.
\]
Indeed $Z=T^{gu+v}Y$ turns the column index in this formula into $gw+v$.
If $\gcd(r,d)=g$, then $\langle r\rangle=\langle g\rangle$, so this
also realizes equality for the prescribed subgroup.  The construction
is algebraic; a coordinate permutation is not a free tape operation.

\subsection{Agreement distance from equivariance}

The one-layer obstruction has a simple probabilistic form on an arbitrary
correctness set.  This form uses no linear-source assumption.

\begin{proposition}[Full-input agreement bound]\label{prop:agreement}
\formal{CompositionAgreement.lean}{full_input_agreement_probability,repeated_class_agreement_probability}
Let $A$ be finite, $|A|\ge2$, and put $q=|A|^d$, the cardinality of one
whole column.  If $K$ is class-uniform, with class sizes $n_J$, then
\begin{equation}\label{eq:agreement-bound}
 \Pr_X[KX=\tau X]
 \le \frac{M}{d}+
       \frac{\sum_J n_J(n_J-1)}{2dq},
\end{equation}
where $X$ is uniform on the complete state space.  If some class has two
distinct receivers, the separate bound
\begin{equation}\label{eq:agreement-pair}
 \Pr_X[KX=\tau X]\le\frac{1+q^{-1}}2
\end{equation}
also holds.
\end{proposition}
\begin{outline}
For a fixed $X$, consider all $d$ rotated inputs $T^sX$ and count the
phases on which $K$ is correct.  At receiver $c=-s$, its view is $X$.
Correctness there requires
\[
 f_{p(c)}(X)=z_c(X),\qquad
 z_c(X):=\sigma^{-c}\operatorname{row}_cX.
\]
The $z_c$ are independent uniform words from a set of cardinality $q$.
In one class, the number of correct phases is at most the largest
multiplicity of these words, hence at most one plus the number of equal
unordered pairs in that class.  Taking expectations and using rotation
invariance of the uniform input law proves \cref{eq:agreement-bound}.

For the second bound, take distinct same-class receivers $c,e$, set
$s=e-c$, and let $G=\{X:KX=\tau X\}$.  If both $X$ and $T^sX$ lie in
$G$, sharedness gives
$\operatorname{row}_eX=\sigma^s\operatorname{row}_cX$.
This event has probability $1/q$, since the two rows are independent.
Thus $\Pr(G\cap T^{-s}G)\le1/q$.  Inclusion--exclusion and equality of
the two marginal probabilities give $2\Pr(G)-1\le1/q$.
\end{outline}

For example, let $F$ be any first map, let $H$ be equivariant, and let
$\gamma=\Pr[KF\ne\tau]$ and $\epsilon=\Pr[F\ne H]$.  The composite
$KH$ has the same class partition as $K$, and is correct on every input
where $KF$ is correct and $F=H$.  Therefore
$\gamma+\epsilon$ is at least one minus either applicable bound above.
This measures equality of entire maps on an input.  It does not measure
the number of output coordinates changed by the first map.

\subsection{Balanced intertwiners}

We now allow a second nonuniform layer, but retain a concrete property
of its first nonuniform operator.  Let $\Sigma$ be a finite alphabet of
size $a\ge2$ with a distinguished symbol $0$.  For width $b\ge1$, regard
a state as a tuple of $db$ words in $\Sigma^d$, one for each row and lane.
Let $L$ be a zero-fixing permutation of $\Sigma^d$, applied identically
to all these words.  Define
\[
 Q=LTL^{-1},\qquad Q_c=T^cQT^{-c}.
\]
In the next statement the group acts on scalar words; the remaining
maps may mix all rows and lanes.

\begin{theorem}[Transitive pairs obstruct balanced outputs]
\label{thm:balanced-pair}
\formal{CompositionWordAction.lean}{transitive_pair_impossible}
Suppose $H$ is a total column-equivariant map and $K$ is a total
class-uniform map on $\Sigma^{d\times d\times b}$, where $d\ge2$.
If $KLH=\tau$ everywhere and $p(c)=p(e)$, the group
$\langle Q_c,Q_e\rangle$ cannot act transitively on
$\Sigma^d\setminus\{0\}$.
\end{theorem}
\begin{outline}
Exactness on the finite complete domain makes $H$ and $K$ bijective.
Writing $q=a^{db}$ for the whole-column alphabet, each column of $K$
has exactly $q^{d-1}$ preimages for each value.  Since input rotation is
bijective, every shared decoder $f_J$ has these same balanced fibres.

Equivariance of $H$ gives $KQ=RK$.  Evaluating at a receiver in $J$ and
changing its view yields
\begin{equation}\label{eq:decoder-covariance}
 f_J(Q_cY)=\sigma f_J(Y)\qquad(c\in J).
\end{equation}
The corresponding inverse identities also hold.  A word in $Q_c,Q_e$
and their inverses therefore rotates $f_J$ by its signed exponent sum.
If the word represents the identity in the actual scalar-word group
$G=\langle Q_c,Q_e\rangle$, it fixes every matrix.  Surjectivity of
$f_J$ and the exact order $d$ of $\sigma$ imply that its exponent sum
is zero modulo $d$.  Consequently there is a surjective homomorphism
\[
 \theta:G\longrightarrow\Z/d\Z,
 \qquad\theta(Q_c)=\theta(Q_e)=1.
\]
Let $N=\ker\theta$.  Equation~\eqref{eq:decoder-covariance} makes $f_J$
constant on all $N$-orbits of matrices.

Suppose $G$ were transitive on the nonzero words.  Normality of $N$
makes all its orbits there have one size $\lambda$, and $G/N$ acts
transitively on the set of these orbits.  Its cardinality $t$ divides
$d$, so
\[
 \lambda=(a^d-1)/t>1,\qquad \gcd(\lambda,a)=1.
\]
The strict inequality follows from $t\le d<a^d-1$ for $a,d\ge2$.
If a matrix $Y$ is nonzero, choose one nonzero row/lane word $y$.
The stabilizers satisfy $N_Y\le N_y$, hence
\[
 |N\cdot Y|=[N:N_Y]=[N:N_y][N_y:N_Y]
            =\lambda[N_y:N_Y].
\]
This uses no transitivity on tuples of words.  Choose a column value
different from $f_J(0)$.  Its fibre excludes the zero matrix and is a
union of $N$-orbits, so its size is divisible by $\lambda$.  But its
balanced size $q^{d-1}$ is a power of $a$, a contradiction.
\end{outline}

Surjectivity alone would give the cyclic quotient but would not give
the last divisibility contradiction.  Likewise, a code decoder defined
on a proper correctness source need not have balanced fibres.

\begin{proposition}[Anchored binary differences]\label{prop:anchored}
\formal{CompositionAnchoredEndpoint.lean}{anchored_adjacent_classes_distinct}
For $d\ge3$, define the binary word permutation
\[
 (Lx)_0=x_0,\qquad (Lx)_i=x_i+x_{i+1}\quad(1\le i<d),
\]
with cyclic indices.  In any complete factorization $KLH=\tau$, with
$H$ equivariant, the partition of $K$ satisfies $p(c)\ne p(c+1)$ for
every receiver $c$.  This holds at every positive width, even when
$H$ and $K$ are nonlinear and mix lanes.
\end{proposition}
\begin{outline}
The inverse is $x_0=y_0$ and
$x_i=y_0+\sum_{j=i}^{d-1}y_j$ for $1\le i<d$.
For $S=\sum_i y_i$, this gives
$Qy=(S,y_2,\ldots,y_{d-1},S+y_0)$.
Conjugating $\langle Q,TQT^{-1}\rangle$ by $L^{-1}$ gives
$\langle T,A\rangle$, where direct substitution in this inverse yields
\[
 Ax=(x_2,x_1+x_2,x_0+x_1,x_3+\eta,\ldots,x_{d-1}+\eta),
 \qquad\eta=x_0+x_1+x_2.
\]
On its first triple, $A$ has the cycle
\[
 100\to001\to110\to010\to011\to101\to111\to100.
\]
Moreover $(Ax)_j+(Ax)_2=x_j+x_2$ for every $j\ge3$.
Thus a power taking a nonzero triple to another with the same third
bit fixes every outside coordinate, regardless of its value.  In
particular the following are local moves in the group:
\[
 100\to010,\quad110\to010,\quad101\to001,
 \quad011\to001,\quad111\to001.
\]
They use powers $3,1,3,4,2$, respectively.  Conjugation by rotations
places them on any cyclic triple.

For a word with at least two ones, choose consecutive ones in cyclic
order.  If their intervening zero gap has length at least two, use
$100\to010$ to shorten it.  A gap of length zero or one then permits
one of the other moves to lower the weight.  Iteration reaches a unit
vector, which a rotation sends to coordinate zero.  Reversing such
paths proves transitivity on all nonzero words, including when $d=3$.
Conjugacy gives the same transitivity for every adjacent pair
$\langle Q_c,Q_{c+1}\rangle$.  Apply \cref{thm:balanced-pair}.
\end{outline}

\subsection{A joint-output period constraint}

Balance examines a shared output column.  A literal period of the final
mask supplies a separate constraint on the complete output map.

\begin{theorem}[Transport of a mask period]\label{thm:period-transport}
\formal{PeriodTransport.lean}{literal_period_complete,fixed_conjugate_iff,fixed_image_preserved}
Let the entry alphabet be finite.  Suppose $U,B,H$ are total
column-equivariant maps, $L$ is any total map, $K$ is class-uniform,
and
\[
 \tau=UKBLH
\]
on the complete state space.  Exactness makes $L$ bijective.  If
$p(c+s)=p(c)$ for every $c$, then its actual conjugated shift
$Q=LTL^{-1}$ commutes with $T^s$.
\end{theorem}
\begin{outline}
All five factors are bijective.  Put $F=UKB$.  Sharedness and the
literal mask period give $KT^s=T^sK$, hence $FT^s=T^sF$.
Equivariance of $H$ gives $FQ=RF$.  As $R$ commutes with $T^s$,
\[
 FQT^s=RFT^s=RT^sF=T^sRF=T^sFQ=FT^sQ.
\]
Cancel the injective map $F$.  The formal proof derives the required
equivalences and injectivity from the total factorization.
\end{outline}

In particular, $T^s$ must preserve
$\operatorname{Fix}(Q)=L\operatorname{Fix}(T)$.
For the anchored binary map of \cref{prop:anchored}, this is the space
of matrices supported on column zero: $L(1,\ldots,1)=e_0$ on each
row/lane word.  No nonzero column rotation preserves this space.

\begin{corollary}[The anchored two-class architecture]\label{cor:anchored-all}
\formal{CompositionAnchoredEndpoint.lean}{anchored_at_most_two_classes_impossible}
For every $d\ge3$ and every positive binary width, there is no complete
factorization $\tau=KBLH$ in which $L$ is the anchored difference,
$B,H$ are arbitrary equivariant total maps, and $K$ uses at most two
receiver classes.
\end{corollary}
\begin{outline}
Absorb $B$ into $K$; this preserves its shared partition.
By \cref{prop:anchored}, that partition is a proper coloring of the
cycle.  An odd cycle has no proper two-coloring.  On an even cycle,
every proper two-coloring alternates and has period two.  This is a
nonzero period for even $d\ge4$, contradicting
\cref{thm:period-transport} and the preceding fixed-space calculation.
\end{outline}

The period theorem permits an arbitrary equivariant outer $U$.  The
combined anchored corollary does not: its balance step concerns $K$
itself.  Neither result excludes two arbitrary nonlinear nonuniform
layers, gives a positive agreement-rate deficit for the anchored
architecture, or supplies a native running-time bound.

\section{Construction requirements and the remaining questions}\label{sec:open}
The results constrain two distinct resources: one-stage correctness dimension
and complete-state composition. Their combination removes several ways of
turning a large one-stage source into a complete algorithm, but it does not
supply an interface theorem for arbitrary native tape computations.

For \cite[Problem~9.1(a)]{SeventeenthPaper2026}, the next capacity question is whether every fixed number of shared deterministic
decoders has dimension at most $N/2+o(N)$ on linear sources. The two-class
$41/64$ bound leaves a constant gap; the fixed-class constructions give only
vanishing normalized excess. Nonlinear source promises and chosen-source
nonlinear-to-linear losses must not be substituted for this maximum-capacity
question. The measured covariance and symmetric-retention deficits provide
more structured intermediate targets. The exact arbitrary-set capacity shows
why a high-entropy fibre alone is insufficient: a native extraction needs the
linearity or actual affine-intersection premise used by its capacity bound.
The mask floors identify limits of the stated positive-term certificates;
improving beyond them requires more compatibility information, a smaller
actual starting ambient, or another argument.

The interval constructions attain the bound with deterministic binary
decoders when $d-M$ is even. This does not settle the original linear-decoder
attainment question for unequal pieces in
\cite[Problem~9.2]{SeventeenthPaper2026}. The diagonal-defect bound supplies
a way to retain actual block losses when testing that question.

For composition, a sharp class-product inequality for arbitrary linear masks
would improve the Ramsey cutoff. A positive normalized rank deficit would
strengthen the present growing rank loss. On the nonlinear side, the complete
masked-neighbor family remains open beyond the uniform, anchored and single-site
cases and the general period restrictions. The balance arguments apply to
actual conjugated operators and balanced fibres; an assigned covariance equation
or a surjective but unbalanced decoder is insufficient.

The complete scattered first factor is a concrete warning: total bijective
completion need not preserve a small class count in the required second map.
The residual bounds offer a different application route, if one can exhibit
a complete structured affine baseline and pay its joint correction record.
They do not supply such a baseline for an arbitrary fast computation.

Finally a native application must account for production of the representation,
access to its entries, the complete state at every intermediate stage, and
recovery of the ordered output. None of the shared-rule theorems proves that a
linear-time native transposer would possess the required fixed class count,
linear source or same-space factorization. This is the unresolved step toward
\cite[Problem~9.4]{SeventeenthPaper2026}. A native theorem giving one of those
interfaces with small paid advice could make the algebraic bounds decisive.
A construction must instead realize a complete compatible sequence of rules
within the native time budget. The two-stage theorems do not exclude a growing
number of stages, and none of the bounds supplies the multiscale time charge
asked for in \cite[Problem~9.5]{SeventeenthPaper2026}.

\appendix
\section{Exact capacity for arbitrary correctness sets}\label{sec:set-capacity}

The linear-source premise in the capacity theorems cannot be replaced by a
large correctness set. Here the alphabet $\Sigma$ need only be finite and
nonempty, with $q$ symbols. Let $d\ge1$, let $\sigma$ rotate columns as in
\cref{sec:model}, and write $x_c$ for row $c$ of $X$. A single total
deterministic decoder $f$ validates a set $C$ when
\[
 f(\sigma^cX)=x_c\qquad(X\in C,\ c\in\Z/d\Z).
\]
Let $M(d,q)$ be the largest possible $|C|$. There is no linearity or
evaluation-cost restriction on $f$.

For each matrix define subgroups
\[
 H_X=\{h:\sigma^hX=X\},\qquad
 K_X=\{a:x_{c+a}=\sigma^a x_c\text{ for every }c\}.
\]
The second rotation acts on a row word. Call the column-rotation orbit of
$X$ eligible if $H_X\subseteq K_X$. Both groups are unchanged when $X$ is
replaced by $\sigma^tX$.

\begin{theorem}[Exact one-class set capacity]\label{thm:set-capacity}
\formal{CapacitySetOrbit.lean}{SetCapacity.exact_capacity,SetCapacity.capacity_orbit_formula,SetCapacity.capacity_average_formula}
With the preceding conventions,
\begin{equation}\label{eq:set-capacity}
 M(d,q)=\sum_{\text{eligible orbits }O}\frac{|K_O|}{|H_O|}
       =\frac1d\sum_X\mathbf 1_{H_X\subseteq K_X}|K_X|.
\end{equation}
One total decoder attains the bound simultaneously on all orbits.
\end{theorem}

\begin{outline}
$H_X$ is a stabilizer. The covariance defining $K_X$ is preserved under
addition and negation of shifts, so $K_X$ is also a subgroup. Column
rotations commute, proving invariance of both groups along an orbit.

A single matrix $X$ is admissible exactly when repeated views prescribe
the same output. Repetition at rows $c,c+h$ means $h\in H_X$, and then
the two demanded rows must agree. Since such $h$ fixes every row word,
this is exactly $H_X\subseteq K_X$.

In an eligible orbit write $Y_a=\sigma^aX$, with $a$ taken modulo $H_X$.
Two orbit words $Y_a,Y_b$ are compatible exactly when $a-b\in K_X$.
Indeed their views at $c,c'=c+a-b$ coincide. Equality of the required
outputs is
$\sigma^a x_c=\sigma^b x_{c+a-b}$, precisely the covariance for $a-b$.
Conversely, any equality of views has $c'=c+a-b+h$ with $h\in H_X$;
membership of $a-b,h$ in $K_X$ makes the outputs agree. Thus all collisions,
including collisions within one word, are accounted for.

A nonempty compatible subset of the orbit lies in one coset of $K_X/H_X$,
and the whole coset is compatible. Its maximum size is $|K_X|/|H_X|$.
Choose one such coset in every eligible orbit. Views from different orbits
cannot collide, so all their output requirements define one partial
function. Extend it arbitrarily to a total decoder. This proves the first
formula and attainment. Each orbit has $d/|H_X|$ elements; summing the
constant value $|K_X|$ on it gives the second formula. For $d=1$ the value
is $q$, and for $q=1$ it is $1$, as the same formula asserts.
\end{outline}

The formal construction chooses an actual representative in every rotation orbit,
takes its entire covariance coset when eligible, and proves all its view
requirements consistent before totalization. Its cardinal and the universal
upper bound are both obtained from actual finite fibres.
The proof uses consistency of deterministic function values. It supplies
neither the additional additive consistency needed for a linear decoder
nor a small description or cheap evaluator for the attaining function.

\begin{corollary}[Prime dimensions and the all-dimension rate]\label{cor:set-capacity-rate}
\formal{CapacitySetCounting.lean}{SetCapacity.prime_capacity,SetCapacity.fixed_H_card,SetCapacity.fixed_K_card}
\formal{CapacitySetRate.lean}{SetCapacity.relative_error,SetCapacity.capacity_lower,SetCapacity.logarithmic_rate_tendsto}
For prime $d$,
\[
 M(d,q)=\frac{q^{d^2}+(d-2)q^d+q}{d}.
\]
For all $d\ge2,q\ge2$,
\[
 \left|\frac{dM(d,q)}{q^{d^2}}-1\right|
 \le(d-1)q^{-d\lceil d/2\rceil},\qquad
 M(d,q)\ge\frac{q^{d^2}}{2d}.
\]
Consequently, at fixed $q\ge2$ along all integer dimensions,
$\log_q M(d,q)=d^2-\log_q d+o(1)$.
\end{corollary}

\begin{outline}
For prime $d$, each subgroup is either trivial or the whole group. There
are $q^d$ matrices with full $H$ (constant rows), $q^d$ with full $K$
(choose $x_0$ and set $x_c=\sigma^cx_0$), and $q$ with both full. The
four cases $(H,K)$ therefore contribute: $q$ words when both are full;
$q^d-q$ when only $K$ is full; zero when only $H$ is full; and
$(q^{d^2}-2q^d+q)/d$ when both are trivial. Their sum is the formula.

For a nonzero shift $h$, put $g=\gcd(d,h)$. The condition $h\in H_X$
permits $g$ free symbols in each row, hence $q^{dg}$ matrices. The
condition $h\in K_X$ has the same count: choose one whole row in each of
the $g$ row cycles and propagate by rotation. On returning around a cycle
the total rotation is a multiple of $d$, so there is no extra constraint.
Put $S=\sum_{h\ne0}q^{d\gcd(d,h)}$. A union bound leaves at least
$q^{d^2}-S$ matrices with trivial $H$; each contributes at least one to
the second sum in \cref{eq:set-capacity}. Dropping eligibility and summing
$|K_X|$ gives the upper bound $dM(d,q)\le q^{d^2}+S$.
Since $\gcd(d,h)\le\lfloor d/2\rfloor$,
$S\le(d-1)q^{d\lfloor d/2\rfloor}$. Division gives the displayed relative
error. It is at most $1/2$ and is bounded by $d/2^d$, which tends to zero.
The lower bound follows, and continuity of $\log_q$ at $1$ gives the
logarithmic remainder. The formal limit runs through dimensions $n+1$,
so it covers every positive integer dimension. No extrapolation from
small dimensions is used.
\end{outline}

\subsection{Line views, linear sources and affine repairs}

For a line, rows and columns are $0,\ldots,d-1$ and the full view is
\[
 \operatorname{view}_c(X)_{i,e}=
 \begin{cases}X_{i,e+c}&0\le e+c<d,\\0&\text{otherwise},\end{cases}
 \qquad -d<e<d.
\]
Fix a distinguished alphabet symbol $0$. The Lean line view uses every
integer offset; for a legal receiver all offsets outside the displayed
window are identically zero. Restriction and zero extension therefore
give the same decoder interface.

\begin{proposition}[High-density line marker]\label{prop:line-marker}
\formal{CapacityLineMarker.lean}{LineMarker.exists_decoder,LineMarker.code_card,LineMarker.affineMarker_finrank}
\formal{CapacityLineMarkerRate.lean}{LineMarker.code_density,LineMarker.code_density_tendsto}
One total deterministic line decoder works on
$C=\{X:\text{the first column is nonzero}\}\cup\{0\}$, whose size is
\[
 |C|=(q^d-1)q^{d(d-1)}+1.
\]
Over a field, its affine subset $X_{0,0}=1$ has dimension $d^2-1$.
No optimality claim for this line construction is made.
\end{proposition}

\begin{outline}
In the view of a matrix with nonzero first column, the leftmost offset
with any nonzero symbol is exactly $-c$. If two such views agree, applying
this boundary test in both directions proves that their receivers agree;
evaluating at offsets $j-c$ then proves their matrices agree. A view
equal to the all-zero view comes only from the zero matrix. Thus all
required outputs are consistent, and one total decoder follows. A nonzero first column
has $q^d-1$ choices; all other columns are free. Fixing one entry to $1$
leaves an affine space of the asserted dimension. Its dimension exceeds
the one-class half bound for $d\ge3$, as checked by
\lean{LineMarker.affineMarker_exceeds_half}; the arbitrary set has density
$1-q^{-d}+q^{-d^2}$, tending to one.
\end{outline}

\begin{theorem}[Deterministic line rules on linear sources]\label{thm:line-deterministic}
\formal{CapacityLineDeterministic.lean}{IsDeterministicLineCode.two_mul_finrank_le,IsDeterministicLineCode.affine_two_mul_finrank_le}
Let $C\subseteq F^{d\times d}$ be a linear subspace of dimension $D$.
Suppose fixed deterministic decoders, constant on each of $M$ nonempty
consecutive row intervals, recover every demanded row from its line view.
The labels may repeat on different intervals. Then
\[
 2D\le d^2+Md.
\]
\end{theorem}

\begin{outline}
Let $T$ shift columns left with a zero final column, and let $Z$ be the
subspace with zero first column. Set $I=C\cap T^{-1}C\cap Z$.
If $\rank T=r$, then
$\dim T^{-1}C=d^2-r+\dim(C\cap\operatorname{im}T)\ge D$.
Hence $\dim I\ge2D-d^2-d$. For $X\in I$ the zero first column gives
the exact view identity
$\operatorname{view}_c(TX)=\operatorname{view}_{c+1}(X)$.
Within one interval, identical arguments to the same decoder imply
$x_{c+1}=Tx_c$. In an interval of length $\ell$, all rows are shifts of
its first row, and their zero first entries force its first $\ell$ entries
to vanish. At most $d-\ell$ entries remain free. Thus
$\dim I\le\sum(d-\ell)=Md-d$, giving the bound. Only equality of decoder
values, not their linearity, was used.
\end{outline}

This extends the linear-decoder line proof of
\cite[theorem ``The line form of Theorem TP'']{SeventeenthPaper2026}; the source remains linear.
On a cycle an affine source $x_0+V$ obeys the corresponding translate
intersection bound whenever the \emph{joint} intersection of the selected
rotated sources is nonempty. Choose a point of that intersection; its
direction is the intersection of the rotated copies of $V$, and the
same codimension and homogeneous covariance bounds apply. Pairwise
intersection does not suffice for three or more shifts. The affine line counterpart is also formalized by
\lean{IsDeterministicLineCode.affine_two_mul_finrank_le}. Its
sufficient affine premise is instead the nonempty threefold slice
$C\cap T^{-1}C\cap Z$, because the view identity needs the boundary zero.
The marker $X_{0,0}=1$ has $C\cap T^{-1}C\ne\varnothing$ for $d\ge2$
but cannot meet $Z$, demonstrating the missing premise.

If the decoders themselves are affine, subtracting their equations at
$x_0+v$ and $x_0$ produces a linear shared rule on $V$. That separate
argument needs no affine-intersection premise; it is invalid for arbitrary
deterministic decoders. The revised affine-entry advice theorem in
\cite[Corollary~6.16]{SeventeenthPaper2026} uses precisely such cancellation.
Our arbitrary-set constructions therefore do not contradict it.

These results sharpen the model boundary relevant to native extraction:
high entropy alone does not provide a linear source or the required affine
intersection. Conversely, recognizing the receiver from an entire static
view does not provide a paid native operation that reads that view.

\section{Diagonal defects and exact displayed kernels}\label{sec:block-defects}

The block flag bound can retain the actual defect on each diagonal block.
This is useful when a proposed column algebra has small invariant blocks,
but its diagonal compressions admit fewer transpose agreements than their
full dimensions would allow.

\begin{theorem}[Unequal block defects]\label{thm:block-defects}
\formal{ConstructionBlockDefects.lean}{upperBlockDescriptor_injective,finrank_upperBlockAgreement_le}
Let $V=\bigoplus_{j=1}^m V_j$, with $\dim V_j=s_j$, and use compatible
bases on the two tensor factors. In these coordinates suppose the linear
operator $K$ has block coefficients
\[
 (KX)_{ij}=\sum_{a,b}L^{ij}_{ab}(X_{ab}),\qquad
 L^{ij}_{ab}=0\quad\text{when }b<j.
\]
Write $K_j=L^{jj}_{jj}$ for its actual diagonal compression and put
$\delta_j=\dim\ker(\tau_{s_j}-K_j)$. Then
\[
 \dim\ker(\tau-K)\le \sum_{i<j}s_i s_j+\sum_j\delta_j.
\]
In particular this applies when the column factors preserve the prefix flag:
that hypothesis gives exactly the displayed coefficient vanishing.
For $b$ coupled lanes the first summand becomes $b\sum_{i<j}s_i s_j$;
each $\delta_j$ is the full diagonal-block defect including its lanes.
\end{theorem}

\begin{outline}
Retain the blocks strictly above the diagonal. On each diagonal block
retain a linear projection onto the kernel of its actual diagonal
compression. We show that this descriptor is injective on the agreement
space. Subtract two inputs with the same descriptor. Their upper blocks
vanish. Descend through the remaining block rows. The flag property
eliminates earlier columns; every other term is already zero by the
upper-block condition or the induction. The diagonal equation therefore
puts the current diagonal block in its diagonal kernel. Its retained
projection is zero, so the block itself is zero. The remaining equations
then set the entire current block row to zero. The descriptor has the
stated dimension.

The Lean statement allows different vector spaces on each rectangle,
explicit transpose equivalences, and all coefficients from later block
columns. It thus includes the tensor formulation and upper-triangular
coupling. The weaker block flag bound follows by replacing each
$\delta_j$ by $s_j^2$ (or $bs_j^2$).
\end{outline}

\subsection{The full kernel of the two-class cyclic lift}

The exhibited seed and lift also admit an exactness statement, stronger
than the existence of a subcode of the same dimension. It concerns these
particular operators; it does not optimize over seed decoders.

\begin{theorem}[Exact displayed cyclic-lift kernel]\label{thm:exact-lift-kernel}
\formal{ConstructionCyclicExact.lean}{binarySeed_full_rule,mem_cyclicLiftCode_iff_full_rule,binaryCyclicDefect_ker,finrank_binaryCyclicLaneDefect_ker}
For the published binary $6\times6$, two-class seed, the full
correctness set of its displayed linear decoders is its $25$-dimensional
code. For $d=6m$, $\gcd(m,6)=1$, and every width $b$, the entire agreement
kernel of the displayed cyclic lift, applied independently to the $b$
lanes, has dimension
\[
 b(18m^2+7m)=b\left(\frac{d^2}{2}+\frac{7d}{6}\right).
\]
More generally, the cyclic lifting construction preserves equality between
a seed code and the full correctness set of its deterministic decoders,
whenever its averaging scalar is invertible. This generic statement is
\lean{mem_liftedCode_iff_full_rule}; the cyclic-index counterpart cited above
uses linear seed decoders.
\end{theorem}

\begin{outline}
Let $A_6$ be the defect map whose $(c,r)$ coordinate is the demanded seed
entry minus the displayed decoder entry. The published generators already
provide a $25$-dimensional subspace of $\ker A_6$. Their private pivot
coordinates define a linear expansion $P$ into that subspace. A sparse
linear certificate $R$ satisfies
\[
 P+RA_6=I.
\]
Lean checks this identity on the $36$ standard basis matrices using
ordinary kernel evaluation. Hence $A_6X=0$ implies $X=PX$ and supplies the
reverse inclusion. The certificate generator does not belong to the
trusted proof; its output is checked again by Lean. The declarations
\texttt{binarySeed\_recovery\_identity} and
\texttt{finrank\_binarySeedDefect\_ker} record this step.

For the lift, let $E$ be the anti-diagonal averaging projector used in
\cref{sec:constructions}. The displayed rule writes a correct input as
\[
 X=\tau(I-E)X+\operatorname{inflate}(R(\operatorname{average}X)),
\]
where $R$ applies the displayed seed rules to every averaged slice.
Average this identity. The first term averages to zero, and averaging
undoes inflation. Every averaged slice therefore satisfies all seed
rules and belongs to the seed code. Subtract its inflation: $(I-E)X$
is symmetric and has zero average. These are precisely the two summands
of the constructed lifted code. The converse is the existing decoder
correctness proof. The Chinese remainder equivalence transfers both
directions to literal cyclic indices. The previous dimension theorem
therefore computes the full kernel, and independent lanes multiply its
dimension by $b$.
\end{outline}

For example, at $d=30,b=1$ the full kernel has dimension $485$,
compared with the two-interval ceiling $480$ and the identity kernel
dimension $465$. These are consequences of the general theorem, not
separate numerical evidence for it.

Reproduce the finite certificate with
\texttt{python3 tools/seed\_recovery\_certificate.py} from this paper's
directory, then run \texttt{lake build} from \texttt{lean/}.
The seed data are read from the sibling Paper~17 tools directory.
The paid native schedule in \cref{thm:paid-lift} is unchanged: exact
kernel equality adds no new operation or hidden decoder oracle.

\section{Mask-sensitive bounds and limits of the certificate methods}
\label{sec:mask-methods}

Throughout this section $C$ is a linear source of dimension $D$ over an
arbitrary field, entries have dimension $b$, and the fixed shared decoders
may be nonlinear. Write $N=bd^2$. Probabilities sample receiver labels,
never inputs to a machine or native algorithms. A \emph{certificate} below
is the numerical right side of a specified upper-bound inequality. Lower
bounds on these numbers are not lower bounds on achievable code dimension.

\subsection{Finite and structured masks}

For a binary mask $p$ let $R_r(s)$ count starts of monochromatic runs of
length at least $r$ along addition by $s$: a start $c$ satisfies
$p(c-s)\ne p(c)=p(c+s)=\cdots=p(c+(r-1)s)$. Constant translation orbits
have no such starts. Put $Q_s=R_1(s)$, the number of directed color changes.

\begin{proposition}[A finite alternative]\label{prop:finite-two-thirds}
\formal{CapacityMaskFinite.lean}{finite_weighted_run_budget,IsSharedCode.four_mul_finrank_le,IsSharedCode.six_mul_finrank_le}
For a binary partition and $d>8$,
\[
 6D\le4bd^2+5bd.
\]
\end{proposition}

\begin{outline}
The two-stage nested inequality with shifts $(1,3)$ gives
$4D\le2N+bd(R_1+R_4+3)$. With shifts $(2,1)$ it gives
$4D\le2N+bd(2R_2+R_3+4)$: use $Q_2=2R_2$, and retain the possible
two translation orbits of the step $2$. The run counts here have step $1$.
Two copies of the first inequality plus one of the second use the budget
\[
 2R_1+2R_2+R_3+2R_4\le2d.
\]
For completeness, charge the two length-one tickets of a run to offset
$0$ in lanes $0,1$, its two length-two tickets to offset $1$ in those
lanes, its length-three ticket to offset $2$ in lane $0$, and its two
length-four tickets to offsets $2,3$ in lanes $1,0$. Tickets occur only
when the charged offset is within that run. Distinct runs cannot collide,
and tickets within one run have distinct position--lane pairs. Thus they
inject into two copies of the cycle. Substitution gives
$12D\le8N+10bd$.
\end{outline}

\begin{theorem}[An explicit family of masks]\label{thm:structured-mask}
\formal{CapacityMaskStructured.lean}{CapacityMaskStructured.first_component_count,CapacityMaskStructured.second_component_count,CapacityMaskStructured.structured_finrank_le}
Let $d=8\ell$, $\ell\ge1$, write $c=a+\ell r$ with $0\le a<\ell$ and
$r\in\mathbb Z/8$, and let
\[
 p(c)=v(a)\mathbin{\mathrm{xor}}u(r),\qquad
 u=(0,1,1,0,1,0,0,1),
\]
where $v$ is arbitrary. Then
\[
 D\le\frac9{16}bd^2+\frac12bd.
\]
\end{theorem}

\begin{outline}
Use nested shifts $3\ell$ and $1$. The first graph has exactly $2\ell$
components. For each fixed $a$, its cyclic order of high indices is
$0,3,6,1,4,7,2,5$, whose labels are $0,0,0,1,1,1,1,0$ up to a common
complement. Each color therefore induces one four-vertex path.

The second shift set is $\{0,3\ell,1,3\ell+1\}$. Its differences retain
the first graph and include $3\ell+1$. For $a<\ell-1$, that step sends
$(a,r)$ to $(a+1,r+3)$. For every $z,\delta\in\{0,1\}$ there is an $r$
with $u(r)=z$ and $u(r+3)=z\mathbin{\mathrm{xor}}\delta$: use
$r=0,6,1,2$ for $(z,\delta)=(0,0),(0,1),(1,0),(1,1)$, respectively.
Taking $z$ to be the desired global color xor $v(a)$ and
$\delta=v(a)\mathbin{\mathrm{xor}}v(a+1)$ connects that color's two
neighboring low-index components. All low indices are thereby connected
within each color. No edge joins colors, so the final graph has exactly
two components. When $\ell=1$ this was already true of the first graph.

The nested bound is $D\le N/2+n_1/4+n_2/4$, with
$n_1=bd^2/4$ and $n_2=2bd$. Substitution proves the claim. In particular,
binary digit parity at $d=2^h$, $h\ge3$, belongs to this family: $u$ is
the parity of the three high bits and $v$ the parity of the remaining bits.
The formal proof constructs classifiers of the actual graph quotients,
using the displayed paths and bridges to prove their fibers are exactly
the connected components.
\end{outline}

\begin{proposition}[Exact direct envelope for digit parity]
\label{prop:digit-parity-envelope}
\formal{CapacityDigitParityEnvelope.lean}{CapacityDigitParity.digit_parity_envelope,CapacityDigitParity.digit_parity_floor_envelope}
For cyclic binary digit parity at $d=2^h$, $h\ge3$, the minimum over all
nonempty shift sets $\Delta$ of the original direct translate certificate
is exactly $5bd^2/8$. This is the numerical envelope of that formula, not
the capacity of the mask.
\end{proposition}

\begin{outline}
Put $s_h(c)=(-1)^{p(c)}$ and
$\eta_h(k)=d^{-1}\sum_c s_h(c)s_h(c+k)$. The formal mask is the literal
parity of the lowest $h$ binary digits, defined by repeated division by
two. Its periodicity identifies the finite correlation with the cyclic
mask statistic. Splitting the average into even and odd positions gives
\[
 \eta_h(2k)=\eta_{h-1}(k),\qquad
 \eta_h(2k+1)=-\tfrac12\bigl(\eta_{h-1}(k)+\eta_{h-1}(k+1)\bigr).
\]
Every correlation lies in $[-1,1]$, since its summands are signs. The
recurrences show that each adjacent sum at level $h\ge1$ is either
$(\eta_{h-1}(k)-\eta_{h-1}(k+1))/2$ or its negative. It therefore lies in
$[-1,1]$. Every odd-index correlation at level $h\ge2$ lies in
$[-1/2,1/2]$. Even indices reduce to the preceding level, and the
nonzero correlations at $h=2$ are $0,-1,0$. Induction gives
$\eta_h(t)\le1/2$ for every nonzero cyclic shift. Consequently
\[
 Q_t=\tfrac d2(1-\eta_h(t))\ge d/4.
\]
The actual first-graph inequality $\chi(\{0,t\})\ge Q_t$ suffices for
the lower bound. Translating a shift set leaves its link relation
unchanged, so every two-element set has certificate at least $5N/8$.
For $q=|\Delta|\ge3$, the ambient term alone gives $2N/3$; a singleton
set gives $N$.

For attainment, set $\ell=d/8$ and split the three high bits from the
remaining bits. The literal parity mask is exactly the preceding
structured family, with $u$ the three-bit parity word. Its actual first
graph for $t=3\ell$ has $2\ell=d/4$ components. This two-element shift
set therefore attains $5N/8$. The value is integral, so flooring any
other certificate cannot improve it. The structured nested bound improves
this entire envelope by $bd(d-8)/16$ when $d>8$.
\end{outline}

\subsection{Run spectra and random masks}

We use the bounded-differences inequality in its following precise form.
For independent coordinates, changing coordinate $i$ changes a real function
$f$ by at most $c_i$; then either tail at distance $t\ge0$ from its mean
has probability at most $\exp(-2t^2/\sum_i c_i^2)$ when $\sum_i c_i^2>0$.
This is McDiarmid's inequality\cite{McDiarmid1989}; the displayed form and
the lower-tail reduction by $-f$ are stated in
\cite[Theorem~1, equation~(2)]{Warnke2012}. No independence between the
individual summands of $f$ is required. For the finite uniform product laws
used here, the shared \texttt{lean-machine} modules
\texttt{FiniteUniformProbability}, \texttt{FiniteBoundedDifferences}, and
\texttt{FiniteUniformWindows} prove the needed inequality and restriction
identities. The exponential-moment proof applies the bounded-interval
estimate successively to actual finite conditional averages; it retains
the full interval length $c_i$ at each coordinate.

\begin{proposition}[The three-stage spectrum]\label{prop:run-spectrum-three}
\formal{CapacityMaskRunSpectrum.lean}{IsSharedCode.run_spectrum_three}
For $d>1$, putting $R_r=R_r(1)$ gives
\[
 D\le \frac N2+bd\left(\frac{R_1}{4}+\frac{R_4}{8}
                                  +\frac{R_{13}}8\right)+\frac{3bd}4.
\]
\end{proposition}

\begin{outline}
Use the actual shifts $1,3,9$ in the three-stage inequality. Their
successive difference sets cover radii $1,4,13$. The corresponding
component counts are at most $R_1+1,R_4+2,R_{13}+2$, respectively;
the first step has a single translation orbit. Their weights are
$1/4,1/8,1/8$, so the additive error is $3bd/4$.
\end{outline}

\begin{theorem}[Almost every binary mask]\label{thm:random-mask-spectrum}
\formal{CapacityMaskRunProbability.lean}{CapacityMaskProbability.spectrum_bad_probability,CapacityMaskProbability.SpectrumGood.source_bound}
Let $d>13$, $\varepsilon\ge0$, and sample the $d$ mask bits independently
and fairly. Except on probability at most
$\exp(-128\varepsilon^2d/529)$, simultaneously every valid source and
decoder family for the sampled mask satisfy
\[
 D\le\left(\frac{41473}{65536}+\varepsilon\right)bd^2+\frac{3bd}4.
\]
\end{theorem}

\begin{outline}
A run-start indicator of length $r<d$ examines $r+1$ distinct bits, with
exactly two favorable assignments. Thus $\mathbb E R_r=d/2^r$.
Changing one bit affects at most $r+1$ of these indicators. For
$F=R_1/4+R_4/8+R_{13}/8$, the coordinate-change constant is therefore
$2/4+5/8+14/8=23/8$, while
$1/2+\mathbb EF/d=41473/65536$. Apply bounded differences with
$t=\varepsilon d$ and then the preceding deterministic inequality.
The event depends only on the mask, so there is no union over fields,
widths, codes or decoders. Conditioning on both colors occurring divides
the failure bound by $1-2^{1-d}$; identifying complementary labelings
does not change any statistic.
\end{outline}

The same reasoning gives the full fixed-depth spectrum. For $k\ge2$,
$d>(3^k-1)/2$, set $r_j=(3^j-1)/2$ and use shifts $3^{j-1}$. Balanced-ternary
digits show by induction that their difference set covers every integer
from $-r_j$ to $r_j$: the preceding interval and its translates by
$\pm3^{j-1}$ are consecutive and exhaustive. Each difference digit is
the difference of two digits in $\{0,1\}$. For a color in cyclic order,
a gap greater than $r$ corresponds to an intervening opposite-color run
of length at least $r$; its components number the larger of one and its
number of long gaps. Summing both colors gives $\chi_r\le R_r+2$;
at radius one the sharper $R_1+1$ also covers constant masks. Consequently
\[
 D\le\frac N2+bd\left(\sum_{j=1}^{k-1}\frac{R_{r_j}}{2^{j+1}}
                  +\frac{R_{r_k}}{2^k}\right)+\frac{3bd}4.
\]
Replace each run count by its mean to obtain the constant
$c_k=1/2+\sum_{j=1}^{k-1}2^{-r_j-j-1}+2^{-r_k-k}$. The coordinate-change
constant is $L_k=\sum_{j=1}^{k-1}(r_j+1)/2^{j+1}+(r_k+1)/2^k$, so for
$\varepsilon\ge0$ the upper tail at $\varepsilon d$ is at most
$\exp(-2\varepsilon^2d/L_k^2)$. The formal endpoints are
\lean{CapacityTernarySpectrum.Good.source_bound} for the source inequality
and \lean{CapacityTernarySpectrum.bad_probability} for its probability bound.
Their proofs compute graph counts, means and sensitivities on the actual
recursively generated translate sets.

\subsection{What the numerical methods cannot certify}

For nonnegative ambient sizes $n_j$ define the exact nested certificate
\[
 B_k(n)=\sum_{j=0}^{k-1}\frac{n_j}{2^{j+1}}+\frac{n_k}{2^k},
 \qquad B_0(n)=n_0.
\]
If $n_1\le n_0$, every $k\ge0$ satisfies
$B_k(n)\ge n_0/2+n_1/4$. For $k\ge2$ these are the first two terms;
for $k=1$ the first ambient has weight $1/2$; for $k=0$ use $n_1\le n_0$.

\begin{theorem}[The nested-component floor]\label{thm:component-method-floor}
\formal{CapacityMaskComparisonProbability.lean}{CapacityMaskProbability.component_floor_probability,CapacityMaskProbability.ChangeGood.certificate_floor}
For $\eta>0$, outside probability
$(d-1)\exp(-\eta^2d/2)$ under independent fair binary labels, every
nested component certificate starting from $n_0=N$ is at least
$(5/8-\eta/4)N$. Shifts and chain lengths may depend on the whole mask.
Taking $\eta=\sqrt{6\log d/d}$ gives failure probability less than $d^{-2}$
for $d>1$.
\end{theorem}

\begin{outline}
For every nonzero $s$, $\mathbb E Q_s=d/2$. A bit change affects at most
the two directed comparisons at that bit and its predecessor, so the
coordinate-change constant is $2$, even for a step of order two.
Bounded differences and a union over the $d-1$ nonzero steps give one
event on which all $Q_s\ge(1/2-\eta)d$. Every first graph has
$\chi(\{0,s\})\ge Q_s$: on a nonconstant translation cycle its
monochromatic components equal its directed changes; a constant cycle
adds a component and no change. The formal proof obtains the same
inequality by counting the actual retained step edges. Thus
$n_1\ge(1/2-\eta)N$; for a zero shift $n_1=N$ instead. Apply the first-term
bound for every chain. Finally
$(d-1)e^{-\eta^2d/2}=(d-1)d^{-3}<d^{-2}$ for the stated $\eta$.
\end{outline}

\begin{theorem}[The full-ambient isolation floor]\label{thm:isolation-method-floor}
\formal{CapacityMaskMethodEndpoints.lean}{CapacityMaskProbability.isolation_rate_probability,CapacityMaskProbability.IsolationGood.rounded_floor}
For $\eta\ge0$, start the isolation recurrence from \emph{all} $d-2$ nonrepresentative
rows of a nonempty binary partition. Write $n_j$ for the number of these
rows isolated in the graph of the $j$th nested shift set, in row units,
and use precisely $B_k(n)$ with this terminal ambient. Outside probability
$(d-1)\exp(-2\eta^2d/9)$ under independent fair labels, every resulting
completed numerical bound $D_{\rm cert}=bd(2+B_k(n))$ satisfies
\[
 \frac{D_{\rm cert}}{N}\ge\frac9{16}-\frac\eta4+\frac1{2d}.
\]
Final integer flooring preserves the safe bound
$9/16-\eta/4-1/(2d)$ for every $b\ge1$. With
$\eta=\sqrt{(27/2)\log d/d}$ the unconditional failure probability is
less than $d^{-2}$. Conditioning on both classes being nonempty divides
it by $1-2^{1-d}$.
\end{theorem}

\begin{outline}
For $s\ne0$, a vertex is isolated in $G_{\{0,s\}}$ precisely when its
label differs from both neighbors at $\pm s$. If $s$ has order two the
neighbors coincide and its probability is $1/2$; otherwise three distinct
fair bits give probability $1/4$. Hence the total $I_s$ has mean at least
$d/4$. A bit change affects at most the indicators centered at that bit
and its two neighbors, so its coordinate-change constant is $3$.
Bounded differences and a union over shifts give, simultaneously,
$I_s\ge d/4-\eta d$. Deleting any two representatives loses at most two
isolated vertices. Thus $n_1\ge d/4-\eta d-2$, independently of how the
representatives were selected. The zero shift has $n_1=d-2$ and also
satisfies this inequality. Since $n_0=d-2$ and all later sizes are
nonnegative, the first-term bound gives
$B_k\ge9d/16-\eta d/4-3/2$. Add the baseline $2$, multiply by $bd$, and
divide by $N$. Final integer flooring of $D_{\rm cert}$ reduces its normalized value
by less than $1/N\le1/d$. The stated logarithmic choice gives
$(d-1)e^{-2\eta^2d/9}=(d-1)d^{-3}<d^{-2}$.
\end{outline}

Both floors are preserved under minima, infima over nonempty families,
convex averages, and pointwise limits of these same certificates. Replacing
counts by upper estimates only raises their right sides. Neither theorem
covers a smaller actual starting support, a sharper terminal estimate,
additional compatibility equations, or signed combinations of inequalities.

\begin{theorem}[The direct method with many classes]\label{thm:direct-method-floor}
\formal{CapacityMaskComparisonProbability.lean}{CapacityMaskProbability.direct_floor_probability,CapacityMaskProbability.EqualityGood.direct_gain_inverse_sqrt}
Let $d\ge2$, $1\le M\le d$, and give the $d$ positions independent uniform
labels in $\{1,\ldots,M\}$. For a nonempty shift set $\Delta$, put
\[
 g_\Delta=\frac{d-\chi(\Delta)}{|\Delta|d},\qquad
 g_{\rm direct}=\max_{\varnothing\ne\Delta\subseteq\mathbb Z/d}g_\Delta.
\]
For $\eta>0$, except on probability at most
$(d-1)e^{-\eta^2d/2}$,
\[
 g_{\rm direct}\le\sqrt{\frac{1/M+\eta}{2}},\qquad
 \min_\Delta N(1-g_\Delta)
       \ge N\left(1-\sqrt{\frac{1/M+\eta}{2}}\right).
\]
In particular, for $\eta_d=\sqrt{6\log d/d}$, the failure probability is
less than $d^{-2}$. If $M\eta_d\le1$, then $g_{\rm direct}\le1/\sqrt M$.
\end{theorem}

\begin{outline}
For each nonzero $s$, let $A_s=|\{c:p(c+s)=p(c)\}|$. The endpoints
are distinct, so $\mathbb EA_s=d/M$. A coordinate change affects at most
two comparisons, even for an order-two step. Bounded differences and a
union over nonzero steps give one event on which
$A_s\le d(1/M+\eta)$ for all of them.

Fix any mask on this event, then any shift set of size $q$. Choose one
orientation for each of its $\binom q2$ unordered distinct pairs. All
graph edges generated by that pair number at most $A_s$ for the chosen
difference $s$. An order-two step double-counts an edge and repeated
differences may count it again; both only increase this upper bound.
Every graph component admits a spanning tree, so
\[
 d-\chi(\Delta)\le |E(G_\Delta)|
       \le\binom q2 d(1/M+\eta).
\]
The formal proof gives the same deficit inequality by the rank of the
edge-incidence map: its kernel is exactly the functions constant on the
actual translate-link quotient. Put $a=1/M+\eta$. Dividing and also using
$d-\chi\le d$ gives
$g_\Delta\le\min((q-1)a/2,1/q)\le\sqrt{a/2}$. The last step follows
from $g_\Delta^2\le(q a/2)g_\Delta\le a/2$; the gain is nonnegative.
This controls all sets from one event, even when chosen after inspecting
the labels. Substitute $\eta_d$, and use $a\le2/M$ when $M\eta_d\le1$.
\end{outline}

An absent label must not be counted as a nonempty class. If $B$ is the event
that every label occurs, the union bound gives
$\Pr(B^c)\le M(1-1/M)^d\le M e^{-d/M}=:q_{\rm empty}$ (at $M=1$ the
first value is zero). Thus the last theorem and $B$ hold simultaneously
with probability at least $1-d^{-2}-q_{\rm empty}$. When
$q_{\rm empty}<1$, its conditional failure probability is at most
$d^{-2}/(1-q_{\rm empty})$; otherwise use the exact positive $\Pr(B)$
as denominator. In the growing range $M\eta_d\le1$, the empty-label bound
tends to zero, since $M\le\sqrt{d/(6\log d)}$ and hence
$M e^{-d/M}\le\sqrt d\,e^{-\sqrt{6d\log d}}$ eventually. Every unlabeled
nonempty partition has exactly $M!$ labelings, so the same conditional
fraction holds for uniform unlabeled partitions. The formal quotient uses
equality of label fibers; it constructs the unique relabeling permutation
and proves the uniform-fiber identity. The conditioning statements and
the growing-range limit are formalized in
\texttt{CapacityMaskUnlabeled.lean} and \texttt{CapacityMaskGrowth.lean}. Together with the
universal difference-cover guarantee, this makes the order $M^{-1/2}$
sharp for that direct certificate method in the stated growing range.
It is not an optimal constant, an attained code rate, or a restriction on
the stronger nested methods. Final flooring changes a normalized direct
certificate by less than $1/N$.

As another comparison, putting $\eta=1/6$ in the binary change-count event
makes every direct two-shift certificate at least $2N/3$, outside
probability $(d-1)e^{-d/72}$. Every direct set of at least three shifts
already has that lower bound from its ambient term. For any fixed
$0<\varepsilon<2/3-41473/65536$ and sufficiently large $d$, the run-spectrum
bound is strictly below this entire direct envelope, except on probability
at most $e^{-128\varepsilon^2d/529}+(d-1)e^{-d/72}$. For rounded bounds
choose $d$ large enough that the difference exceeds one dimension.
No comparison here extracts a shared code from a native computation.

\section{Extension-linear residuals with binary-linear decoders}
\label{sec:field-linear-residual}

This section concerns scalar binary entries, so $b=1$. Put
$E=\mathbb F_{2^d}$ and $\Phi(x)=x^2$. Use the reversed normal basis from
the interval construction to identify a binary row with $E$; literal column
rotation then acts on every row by $\Phi$. No prime or coprimality
hypothesis on $d\ge1$ is needed. Linearity of a decoder below means
$\mathbb F_2$-linearity on the complete binary view, not $E$-linearity.

Fix a partition $p:\mathbb Z/d\to\mathcal J$ into $M$ nonempty classes,
and a representative $r_J$ in every class. For a binary linear source $C$,
let $C_0=C\cap\ker P$, where $P$ is the actual class projection from the
completion theorem. Thus every representative row vanishes on $C_0$.
The additional hypothesis is that the image $K\subseteq E^d$ of
\emph{this intersection}, in the fixed normal basis, is an $E$-subspace.
It is not a hypothesis on the image $(I-P)C$, nor on $C$ or the baseline
separately. Define
\[
 S=\{c:\exists x\in K,\ x_c\ne0\},\qquad
 a=|S|,\quad a_J=|S\cap p^{-1}(J)|,\quad k=\dim_E K.
\]

\begin{lemma}[Frobenius components]\label{lem:frobenius-components}
\formal{CapacityFieldLinearComponents.lean}{CapacityFieldLinear.phaseBasis,CapacityFieldLinear.sum_components,CapacityFieldLinear.components_unique,CapacityFieldLinear.component_of_scalarTest,CapacityFieldLinear.component_zero_of_scalarTest_zero}
For every binary-linear map $f:E^S\to E$ there are unique $E$-linear
functionals $\ell_t:E^S\to E$, indexed by $t\in\mathbb Z/d$, such that
\[
 f(y)=\sum_t\Phi^t(\ell_t(y)).
\]
If $K\le E^S$, $c\in S$, and $f(\Phi^c x)=x_c$ for all $x\in K$, then on
$V_c=\Phi^cK$ the functional $\ell_{-c}$ is the coordinate $y\mapsto y_c$
and all other components vanish. For any phase $c\in\mathbb Z/d$,
if $f(\Phi^cx)=0$ for every $x\in K$, all components vanish on that view.
\end{lemma}

\begin{outline}
The $d$ distinct automorphisms $\Phi^t$ are linearly independent over $E$.
The space $\operatorname{Hom}_{\mathbb F_2}(E,E)$ has $E$-dimension $d$,
so they form a basis. Expand the binary-linear map
$\lambda\mapsto f(\lambda y)$ in this basis. Its coefficient at $t$,
followed by $\Phi^{-t}$, is $\ell_t(y)$. Uniqueness of coordinates proves
additivity and
$\ell_t(\lambda y)=\lambda\ell_t(y)$, hence $E$-linearity. Evaluating
at $\lambda=1$ gives the displayed decomposition and its uniqueness.

For $y=\Phi^c x\in V_c$, every $\lambda y$ is still in $V_c$ because
$K$ is $E$-linear. The decoder identity therefore gives
\[
 \sum_t\Phi^t(\lambda)\Phi^t(\ell_t(y))
       =\Phi^{-c}(\lambda)\Phi^{-c}(y_c).
\]
Independence of the automorphisms as functions of $\lambda$ forces the
claimed coefficient identities. This argument uses the actual row demand;
it supplies no receiver index to the decoder.
\end{outline}

\begin{theorem}[Active-row bound]\label{thm:field-linear-active}
\formal{CapacityFieldLinearResidual.lean}{CapacityFieldLinear.active_class_residual_bound,CapacityFieldLinear.original_source_class_bound,CapacityFieldLinear.finrank_encodedResidual}
Suppose $C$ has total binary-linear shared decoders and its actual residual
$K$ is $E$-linear as above. Then, for every class $J$,
\begin{equation}\label{eq:field-linear-class}
 a_J+k\le a,\qquad
 D+d a_J\le Md+da,\qquad D=\dim_{\mathbb F_2}C.
\end{equation}
Consequently
\[
 D\le Md+d\bigl(a-\max_J a_J\bigr)
       \le Md+d(1-1/M)a.
\]
\end{theorem}

\begin{outline}
Restrict inputs to their active rows $S$, extending by zero elsewhere.
This gives an $E$-linear isomorphism between $K$ and an $E$-subspace of
$E^S$. Restrict each actual decoder by this zero extension and the normal
basis identification. For $c\in S\cap p^{-1}(J)$ take its component
$L_c=\ell_{J,-c}$. By \cref{lem:frobenius-components}, $L_c$ equals the
$c$th coordinate on $V_c$ and vanishes on every $V_e$ with
$e\in p^{-1}(J)$, $e\ne c$. The coordinate is nonzero on $V_c$ because
$c$ is active. A relation $\sum_c u_cL_c=0$, restricted successively to
these views, therefore has every $u_c=0$.

The representative row is zero, so all $L_c$ annihilate $V_{r_J}$.
The latter has $E$-dimension $k$ in an ambient of dimension $a$; its
annihilator has dimension $a-k$. The independent family has size $a_J$,
giving the first inequality. In Lean the annihilating functionals are
conjugated by the representative Frobenius automorphism, so they all lie
in the annihilator of $K$ itself.

The normal-basis equivalence and the field tower give
$\dim_{\mathbb F_2}C_0=dk$. Rank--nullity for $P|_C$, whose image has
dimension at most $Md$, gives $D\le Md+dk$. This proves the second
inequality without assuming that the baseline is contained in $C$.
Take the largest $a_J$ for the first consequence; since
$\sum_J a_J=a$, its largest term is at least $a/M$.
\end{outline}

\begin{corollary}[One and two classes]\label{cor:field-linear-half}
\formal{CapacityFieldLinearCorollaries.lean}{CapacityFieldLinear.one_class_residual_eq_bot,CapacityFieldLinear.one_class_original_bound,CapacityFieldLinear.two_class_original_bound,CapacityFieldLinear.two_class_half_bound}
Under the same hypotheses, one class gives $K=0$ and $D\le d$.
For two nonempty classes $J_0,J_1$,
\[
 D\le d\bigl(1+\min(|J_0|,|J_1|)\bigr),\qquad
 2D\le d^2+2d.
\]
\end{corollary}

\begin{outline}
For one class $a_J=a$, so $k=0$. For two classes
$a=a_{J_0}+a_{J_1}$, and the active-row bound gives
$k\le\min(a_{J_0},a_{J_1})$. Each class contains its zero representative,
so $a_J\le |J|-1$. Substitute this into $D\le2d+dk$ and then use
$|J_0|+|J_1|=d$.
\end{outline}

This excludes above-half asymptotic rates in this construction class with
binary-linear decoders. Deterministic decoders for extension-linear graph
codes can be nonlinear, so those constructions do not contradict the
one-class conclusion. The theorem does not cover arbitrary binary sources,
alternative row identifications, or $b>1$, and it supplies no native
extraction or implementation claim.

\section{Measured structural restrictions}\label{sec:structural-bounds}

The parameters in this appendix describe an actual source subspace.
They are not running-time invoices. Neither result supplies its own
small-parameter hypothesis for an arbitrary tape computation.

\subsection{A robust bound for almost field-linear residuals}

Fix the reversed normal-basis identification of each binary row with
$E=\F_{2^d}$. Thus column rotation is coordinatewise Frobenius $\Phi$ on
the $d$ row vectors in $E^b$. For a two-class code $C$, let $P$ be the
representative-row projection of Section \ref{sec:structure}, and put
\[
 L=C\cap\ker P,\quad H=\operatorname{span}_E L,\quad
 \rho=\dim_{\F_2}H-\dim_{\F_2}L.
\]
The common $E$ scalar acts on all rows and lanes; no independent-lane
decomposition is assumed. Let $a$ be the number of potentially nonzero
residual rows, so $a\le d-2$.

\begin{theorem}[Scalar-closure defect]\label{thm:closure-defect}
\formal{CapacityClosureAsymptotic.lean}{CapacityClosure.residual_finite_bound, CapacityClosure.capacity_bound, CapacityClosure.rounds_eq_floor, CapacityClosure.error_ratio_tendsto, CapacityClosure.uniform_eventual_bound}
For $d\ge2$, $b\ge1$, arbitrary deterministic shared decoders, and every
integer $R\ge1$, set
\[
 \mathcal E_R(d)=\frac{d}{10R}+
 \frac{7\sqrt d}{5R}\frac{2^{R/2}-1}{\sqrt2-1}.
\]
Then
\begin{align*}
 \dim_{\F_2}L&\le bd\left(\frac a2+\frac d{10}+\mathcal E_R(d)\right)
                       +\frac{13}5\rho,\\
 \dim_{\F_2}C&\le\frac35N+bd+bd\mathcal E_R(d)+\frac{13}5\rho.
\end{align*}
Taking $R=\lfloor\tfrac12\log_2d\rfloor$ for $d\ge4$ gives
$\mathcal E_R(d)=o(d)$. In particular $\rho=o(N)$ implies rate at most
$3/5+o(1)$, uniformly over all widths and all lane coupling.
\end{theorem}
\begin{outline}
For a shift set $F$ containing zero write
$H_F=\bigcap_{e\in F}\Phi^{-e}H$ and
$L_F=\bigcap_{e\in F}\Phi^{-e}L$. The first is $E$-linear. The binary
map from $H_F$ to the direct sum of the quotients
$\Phi^{-e}H/\Phi^{-e}L$ has kernel $L_F$, so
\begin{equation}\label{eq:hull-intersection-defect}
 \operatorname{codim}_{\F_2}(L_F,H_F)\le |F|\rho.
\end{equation}
Delete loops from the same-class link graph for $F$. One joint
binary-linear covariance map is
\[
 B_F(X)=(X_v-\Phi^{v-u}X_u)_{u\to v},
\]
with one orientation for each edge. Actual decoder equality makes this
map zero on $L_F$. Hence its rank on $H_F$ is at most $|F|\rho$.

Choose a normal element $\lambda\in E$, so
$\lambda\ne\Phi^\delta\lambda$ for every nonzero $\delta\in\Z/d$.
Multiplication by $\lambda$ permutes $H_F$, so $B_F(\lambda X)$ has the
same rank on that space. For an edge of displacement $\delta$, the two
observations $A=X_v-\Phi^\delta X_u$ and
$A'=\lambda X_v-\Phi^\delta(\lambda)\Phi^\delta X_u$ give
\[
 A'-\lambda A=(\lambda-\Phi^\delta\lambda)\Phi^\delta X_u.
\]
They recover both entire endpoint vectors. Therefore the projection onto
all nonisolated rows has binary rank $t_F$ satisfying
\begin{equation}\label{eq:joint-hull-rank}
 t_F\le2|F|\rho.
\end{equation}
This uses one pair of joint maps, not a sum of edge ranks. It does not
assert decoder correctness on $H_F\setminus L_F$.

Let $n_F=bd|S\cap\operatorname{Iso}(F)|$, where $S$ is the active row set.
Projection gives $\dim_{\F_2}H_F\le n_F+t_F$ and, for any shift $s$,
\[
 \dim_{\F_2}(H_F+\Phi^{-s}H_F)\le n_F+2t_F.
\]
The two error images have separately paid dimensions $t_F$; neither need
be invariant. For a nested chain $F_i=F_{i-1}\cup(F_{i-1}+s_i)$,
$F_0=\{0\}$, Grassmann and this inequality yield
\[
 \dim_{\F_2}L\le \sum_{i<j}\frac{n_{F_i}}{2^{i+1}}
                         +\frac{n_{F_j}}{2^j}+(2j-1)\rho.
\]
Indeed unrolling contributes $\sum_{i=1}^jt_{F_i}/2^i\le2j\rho$
by \eqref{eq:joint-hull-rank}, and
$\dim_{\F_2}L=\dim_{\F_2}H-\rho$ subtracts the original defect once.

For $F_1=\{0,s\}$ and $F_2=\{0,s,3s,4s\}$ the full isolation counts obey
\[
 I_1\le Q_s-\tfrac12Q_{2s}+g_s,\qquad
 I_4\le R_4(s)+2g_s.
\]
For the first assertion, on a nonconstant $s$-cycle the isolated vertices
are precisely singleton runs, counted by $Q_s-R_2(s)$; constant or
one-point cycles are covered by $g_s$. For the second, a vertex isolated
in the radius-four graph must be a singleton run whose following
opposite-color run has length at least four; these following runs are
distinct. Cycles of length at most eight have at most two isolated
vertices, paid by $2g_s$. Use Lemma \ref{lem:binary-runs} and its weaker
budget $Q_s+R_2+2R_4\le d$. The one- and two-stage inequalities become
\begin{align*}
 \dim L&\le bd(a/2+Q_s/2-Q_{2s}/4+g_s/2)+\rho,\\
 \dim L&\le bd(a/2+d/8+Q_s/8-3Q_{2s}/16+3g_s/4)+3\rho.
\end{align*}
Their convex combination with weights $1/5,4/5$ is
\[
 \dim L\le bd\left(a/2+d/10+(Q_s-Q_{2s})/5+7g_s/10\right)+13\rho/5.
\]
Average $s=2^ru$ for uniform $u\in\Z/d$ and $0\le r<R$.
The $Q$ terms telescope; the initial average is at most $d/2$ and the
terminal average is nonnegative. The gcd estimate proved in Section
\ref{sec:capacity}, $\mathbb E\,g_{2^ru}\le2\sqrt d\,2^{r/2}$,
gives exactly $\mathcal E_R(d)$. Finally
$\dim C\le\dim L+2bd$ and $a\le d-2$ prove the second bound.
For the stated $R$, both terms of $\mathcal E_R(d)/d$ tend to zero.
\end{outline}

At $\rho=0$ the theorem covers genuinely $E$-linear residuals with
arbitrary deterministic decoders. It improves the general $41/64$
asymptotic bound whenever $\rho/N<1/64$ by a fixed margin. The parameter
can nevertheless be extensive for the cheapest partial transposer.

\begin{proposition}[Symmetric-copy control]\label{prop:symmetric-closure}
\formal{CapacityClosureControl.lean}{CapacityClosure.symmetric_isSharedCode, CapacityClosure.scalarHull_supportedSymmetric, CapacityClosure.symmetric_copy_defect}
On the lane-wise symmetric source, the identity output is correct for every
receiver partition. For two distinct representative rows its residual satisfies
\[
 \dim L=\frac{b(d-2)(d-1)}2,\qquad
 \dim_{\F_2}H=bd(d-2),\qquad
 \rho=\frac{b(d-2)(d+1)}2.
\]
Thus $\rho\sim N/2$, although straight streaming copy realizes this
partial correctness set.
\end{proposition}
\begin{outline}
Zeroing representative rows of a symmetric matrix also zeroes their columns.
Restriction to the remaining rows identifies the residual with symmetric
matrices of side $d-2$. Each remaining diagonal unit has a nonzero
normal-basis coordinate. Extension-field multiples of these units span every
row and lane of the zero-representative ambient. The dimension difference
gives the displayed defect. The native realization is the literal streaming
copy operation; the formal statements establish its symmetric correctness
set and the exact algebraic dimensions.
\end{outline}

\subsection{Retaining symmetric directions}

Let $\mathcal S=\operatorname{Fix}\tau$, put
$\mathcal S_0=C\cap\mathcal S$, and write
\[
 s_0=\dim\mathcal S=bd(d+1)/2,\qquad
 \eta=s_0-\dim\mathcal S_0.
\]
The following restriction uses actual demanded rows and no field-extension
structure. In particular it holds over every field and for nonlinear
decoders.

\begin{theorem}[Symmetric-retention bound]\label{thm:pairing-retention}
\formal{CapacityRetentionFinite.lean}{IsSharedCode.retention_bound}
Choose $k\ge0$ shifts $t_i$, let $g_i=\gcd(d,t_i)$, and let
$U=\{c:p(c)\ne p(c+t_i)\text{ for every }i\}$, $u=|U|$. Then
\begin{equation}\label{eq:pairing-finite}
 D\le s_0+(2k-1)\eta+\frac{bd}{2}\sum_i(g_i-1)
                         +\frac b2u(u-1).
\end{equation}
\end{theorem}
\begin{outline}
Put $A(X)=X-X^{\mathsf T}$ and $Z_t=\{c:p(c)=p(c+t)\}$.
If $X=T^tY\in C$ with $Y\in\mathcal S_0$, the same-view equations give
$\operatorname{row}_cX=\operatorname{row}_{c+t}Y$ for $c\in Z_t$.
Symmetry of $Y$ then gives $X_{c,j}=X_{j,c}$ in every lane. The restriction
of $A$ to rows $Z_t$ therefore vanishes on both $\mathcal S_0$ and
$C\cap T^t\mathcal S_0$. Grassmann and rank--nullity bound its rank by
$N-\dim(\mathcal S_0+T^t\mathcal S_0)$.

Stacking the $k$ row maps costs at most the sum of their ranks; no
independence is assumed. An antisymmetric image with all covered rows zero
is supported on $U\times U$, has zero diagonal, and is determined by
$bu(u-1)/2$ coordinates in every characteristic. Since $\ker(A|_C)=\mathcal S_0$,
this proves the stronger exact inequality
\begin{equation}\label{eq:retained-space-exact}
 D\le\dim\mathcal S_0+
 \sum_i\bigl[N-\dim(\mathcal S_0+T^{t_i}\mathcal S_0)\bigr]
 +\frac b2u(u-1).
\end{equation}
It is formalized as \lean{IsSharedCode.intrinsic_retention_bound}.

The remaining estimate uses the actual symmetric spaces. For
$g=\gcd(d,t)$, the two coordinate symmetries generate the dihedral action
whose rotations send $(i,j)$ to $(i+at,j-at)$ and whose reflections send
it to $(j-at,i+at)$. Exactly $g$ rotations fix all $d^2$ coordinates;
the other rotations fix none. Each of the $d$ reflections fixes $d$
coordinates. The space of invariant coordinate functions is the function
space on the orbits, so Burnside's formula gives
\[
 \dim(\mathcal S\cap T^t\mathcal S)=bd(g+1)/2.
\]
This counting proof does not divide by the group order in the field and
remains valid in characteristic two. Grassmann, and replacing each
symmetric summand by its codimension-$\eta$ subspace, give
\[
 N-\dim(\mathcal S_0+T^t\mathcal S_0)\le2\eta+bd(g-1)/2.
\]
Substitution in \eqref{eq:retained-space-exact} proves the theorem.
\end{outline}

\begin{corollary}[Averaged retention]\label{cor:pairing-average}
\formal{CapacityRetentionAverage.lean}{IsSharedCode.retention_average,IsSharedCode.retention_normalized}
Let the class sizes be $n_J$, put $a_d=d^{-1}\sum_t\gcd(d,t)$,
and let $k\ge1$. Then
\begin{equation}\label{eq:pairing-average}
 D\le s_0+(2k-1)\eta+\frac{bdk}{2}(a_d-1)
       +\frac b2\left[\sum_Jn_J(1-n_J/d)^{k/2}\right]^2.
\end{equation}
Consequently, with $h=\eta/N$ and $\delta=D/N-(d+1)/(2d)$,
\begin{equation}\label{eq:pairing-normalized}
 \delta\le(2k-1)h+\frac{k(a_d-1)}{2d}+\frac{2M^2}{(k+2)^2}.
\end{equation}
\end{corollary}
\begin{outline}
Average over the full set of $k$-tuples of shifts. A row in a class of
size $n_J$ is uncovered for exactly $(d-n_J)^k$ tuples. The finite
Cauchy--Schwarz inequality bounds the sum of squared uncovered-row counts
by the square of the sum of square roots of these marginal counts.
Grouping rows by their classes and using $u(u-1)\le u^2$ gives the
last term of \eqref{eq:pairing-average}; no independence between row
events is assumed. Each shift coordinate has the uniform marginal, so
its mean gcd cost is exactly $a_d-1$.
For $0\le x\le1$, Bernoulli's inequality gives
$x(1-x)^{k/2}\le2/(k+2)$. Applying it to every class and dividing by
$N$ proves \eqref{eq:pairing-normalized}.
\end{outline}

\begin{corollary}[A fixed rate gain requires a fixed loss]\label{cor:pairing-loss}
\formal{CapacityRetentionStability.lean}{RetentionAverage.stability,RetentionAverage.stability_bounded_classes,RetentionAverage.fixed_gain_requires_loss}
For fixed $M$, $\eta=o(N)$ implies $D\le(1/2+o(1))N$ as $d\to\infty$,
uniformly in the field and width. More generally this holds when
$M\eta/N\to0$ and $M(a_d-1)/d\to0$. For fixed $M$ and
$0<\varepsilon\le1/2$, if $\liminf D/N\ge1/2+\varepsilon$, then
\[
 \liminf\eta/N\ge\varepsilon^{3/2}/(12M).
\]
\end{corollary}
\begin{outline}
The divisor bound $a_d\le2\sqrt d$ makes the gcd term vanish for fixed
$k$. The formalization permits the field, lane space, dimension and
class set to vary along the sequence, retaining an actual shared code
at every index. For growing $M$, take $k=\lceil Mt\rceil$ with a fixed
real $t\ge1$. The normalized upper bound is at most
\[
 \frac1{2d}+4tM\left[\frac\eta N+\frac{a_d-1}{2d}\right]+\frac2{t^2}
\]
above one half. First let $d\to\infty$ and then let $t\to\infty$.
For the quantitative statement take $k=\lceil2M/\sqrt\varepsilon\rceil$.
The last term is at most $\varepsilon/2$, and
$k\le3M/\sqrt\varepsilon$. Take a subsequence attaining the lower limit
of $\eta/N$ and pass the finite bound to the limit; the rate lower limit
still applies on that subsequence. This gives the stated constant.
Both rates lie in $[0,1]$, so all lower limits used here have explicit
boundedness hypotheses in Lean.
\end{outline}

\begin{corollary}[Coprime partners and exact retention]\label{cor:coprime-retention}
\formal{CapacityRetentionCoprime.lean}{IsSharedCode.full_symmetric_retention,IsSharedCode.eq_symmetric_of_at_most_one_exception,IsSharedCode.prime_full_symmetric_retention}
Suppose $C\supseteq\mathcal S$. Let $V$ be the rows $c$ with no
same-class partner $c+t$ for which $\gcd(d,t)=1$, and put $v=|V|$.
Then, in every positive dimension,
\[
 D\le s_0+b\binom v2.
\]
In particular $v\le1$ implies $C=\mathcal S$. If $d$ is prime, $V$
is exactly the set of rows in singleton classes. Thus a partition into two
nonsingleton classes forces $C=\mathcal S$.
\end{corollary}
\begin{outline}
Use all coprime shifts in \cref{thm:pairing-retention}. Then $\eta=0$,
every gcd cost is zero, and the uncovered rows are exactly $V$. For
$v\le1$ the dimension bound and $\mathcal S\subseteq C$ give equality.
In prime dimension the nonzero shifts are precisely the coprime shifts.
\end{outline}

\begin{proposition}[Paired-lane control]\label{prop:paired-lane-control}
\formal{CapacityRetentionControl.lean}{RetentionControl.laneSwap_eq_transpose_iff,RetentionControl.paired_source_rate_and_loss}
For any field and even entry width $b$, swapping each paired group of
lanes entrywise is transpose-correct exactly on the source consisting of
lane pairs $(A,A^{\mathsf T})$. This source has
\[
 D=N/2,\qquad \eta=bd(d+1)/4.
\]
\end{proposition}
\begin{outline}
The source is the graph of transposition on one member of each lane pair,
so it is linearly equivalent to an arbitrary matrix with $b/2$ lanes.
Its intersection with the symmetric space consists of pairs $(A,A)$
with $A$ symmetric. The two explicit linear equivalences give the
dimensions. The formalization allows vector-valued lane groups.
\end{outline}

Over the binary alphabet, one fixed finite-state stream transducer realizes
the adjacent-pair interchange by buffering two input bits and emitting them
in reverse order. It repeats without computing the entry width. For every
even width its correctness source is \cref{prop:paired-lane-control}, with
$\eta/N=(d+1)/(4d)\to1/4$; infinitely many canonical widths are even.
This written implementation refutes an unqualified small-loss time supplier.
It does not address a supplier restricted to fixed excess above one half.

\section{Why a large linearly decoded subcode is not free}\label{sec:linearization}

The source remains a binary linear subspace throughout this appendix.
Only decoder linearity changes. The bounds allow a candidate subcode
and its decoder to mix all lanes; lane-wise upper bounds are not summed
without a joint argument.

\begin{lemma}[Projection with coupled lanes]\label{lem:lane-projection}
\formal{CapacityLinearization.lean}{IsLinearSharedCode.inr_comap,linear_subcode_lane_bound}
Let $G,H$ occupy disjoint lane groups with the same receiver partition.
Suppose every linearly decoded subcode of $H$ has dimension at most $L_H$.
Every linearly decoded subcode $C'\le G\oplus H$ has dimension at most
$\dim G+L_H$, even with arbitrary cross-lane linear decoding.
\end{lemma}
\begin{outline}
Project $C'$ onto the $G$ coordinates. Its image has dimension at most
$\dim G$. Its kernel is a subcode of $H$; restrict each full linear
decoder to zero $G$ input and project its output to the $H$ coordinates.
This is a legal linear shared decoder on the kernel. Rank-nullity gives
the assertion. No decoder on the projected $G$ image is needed.
\end{outline}

\begin{lemma}[A width-stable triple obstruction]\label{lem:triple-obstruction}
\formal{CapacityTripleBinary.lean}{CapacityTriple.exists_width_stable_hard_source}
Let $d=3s$ be odd. Partition $\{0,\ldots,s-1\}$ and enlarge each class
by its translates by $s$ and $2s$. There is a positive integer $\ell$
such that, for every width $B$ divisible by $\ell$, a binary linear source
$H_B$ of dimension $Bd^2/3$ has deterministic shared decoding, while
every binary-linearly decoded subcode has dimension at most $2Bd^2/9$.
\end{lemma}
\begin{outline}
Choose $E=\F_{2^\ell}$ containing a primitive $d$-th root $\zeta$;
such an extension exists since $d$ is odd. Write
$e_k(t)=\zeta^{kt}$, and let the allowed row space be
\[
 R=\operatorname{span}_E\{e_{3j+1}+e_{3j+2}:0\le j<s\}.
\]
All $d$ rows vary independently in $R$. The resulting binary source has
dimension $\ell ds=\ell d^2/3$. For distinct phases $c,e$, equality
between a vector in $\sigma^cR$ and one in $\sigma^eR$ compares the two
coefficients in each disjoint frequency pair. Their ratio is respectively
$\zeta^c$ and $\zeta^e$. Thus every pair coefficient is zero, and the two
row spaces intersect only in zero. The entire nonzero matrix views are
therefore disjoint, and the total phase-identifying decoder from the proof
of Theorem \ref{thm:interval-attainment} returns the actual demanded row.

Every row in $R$ satisfies $(I+\sigma^s+\sigma^{2s})x=0$, because the
two frequencies in each pair are nonzero modulo three. Thus every
source matrix satisfies $(I+T^s+T^{2s})X=0$. If a subcode has binary-linear
shared decoders, applying one decoder at a view $T^aX$, using the three
same-class receivers, forces
\[
 \operatorname{row}_aX+\operatorname{row}_{a+s}X+
 \operatorname{row}_{a+2s}X=0\qquad(0\le a<s).
\]
The joint row-sum map on the whole hard source is surjective onto $s$
independent copies of $R$: set just the first row of each triple to any
chosen vector. Its binary rank is $\ell s^2$, giving the claimed
$2\ell d^2/9$ upper bound. Repeat the field-coordinate blocks to width
$B$. The same joint map is surjective onto all repeated blocks, so its
rank scales by $B/\ell$ even when the subcode and decoders mix those
blocks arbitrarily.
\end{outline}

\begin{theorem}[Above-half sources can resist subcode promotion]\label{thm:subcode-barrier}
\formal{CapacitySubcodeMixture.lean}{CapacitySubcode.exists_above_half_barrier}
\formal{CapacitySubcodeBarrier.lean}{CapacitySubcode.above_half_seed_barrier}
For every odd $d\ge9$ divisible by three, at suitable binary widths and
a two-class partition there is a deterministically decoded linear source
of rate
\[
 \frac12+\frac1{8d+18}
\]
whose every linearly decoded subcode has rate at most
$\frac12-\frac1{8d+18}$.
More generally, a two-class seed at odd $s$ coprime to three with rate
$1/2+\varepsilon$, $\varepsilon>0$, gives at $d=3s$, after a width
enlargement, a source of rate $1/2+\delta$ whose every linearly decoded
subcode has rate at most $1/2-\delta$, where
\[
 \delta=\frac{\varepsilon}{18\varepsilon+12}.
\]
\end{theorem}
\begin{outline}
For the second statement apply Theorem \ref{thm:cyclic-lift} with $m=3$.
The good source has rate $\alpha=1/2+e$, $e=\varepsilon/3$, on the same
tripled partition used in Lemma \ref{lem:triple-obstruction}.
Repeat binary lane blocks to put it and the hard source at a common width
$B$; this repeats codes, and does not extend a nonlinear decoder's field.
Choose positive integers $g,h$ with
$h/(g+h)=e/(e+2/9)$, possible since all seed rates are rational.
Take $g$ good blocks and $h$ hard blocks. Their concatenated deterministic
decoders work on the direct sum. Its rate is
$(1-t)\alpha+t/3$; by the joint-width hard bound and Lemma
\ref{lem:lane-projection}, every linearly decoded subcode has rate at
most $(1-t)\alpha+2t/9$. Substitution of $t=e/(e+2/9)$ gives the stated
symmetric offsets $\delta=e/[18(e+2/9)]$.

For the unconditional statement use the symmetric good source, of rate
$1/2+1/(2d)$, at the tripled two-class partition. Take $g=4d$ good blocks
and $h=9$ hard blocks. The two rates are respectively
$(2d+5)/(4d+9)$ and at most $(2d+4)/(4d+9)$, which are the displayed
offsets from one half.
\end{outline}

The unconditional excess tends to zero. A family with a fixed positive
excess in the second statement is conditional on such a seed family;
it is not supplied here. Nor does the theorem separate the maximum
capacities over unrelated codes: its conclusion concerns subcodes of
one actual source. The enlarged widths can exceed both one and the
canonical logarithmic width, so neither restricted-width conclusion is
inferred.

\section{A reflection barrier at every binary word width}
\label{sec:prime-reflection}

The triple-relation family gives a direct linearization obstruction at its
specified cyclic dimensions and extension-compatible widths. The following
family has a different scope: every binary width $b\ge1$, odd prime side
length, and source rate tending to one half. It does not replace either the
above-half subcode construction or the triple relation-rank diagnostic.

Let $d=2m+1$ be an odd prime, and put
\[
 W_b=\{x\in(\mathbb F_2^b)^d:x_0=0,\ x_j=x_{-j}\ \forall j\},\qquad
 C_b=\{X:\text{every row of }X\text{ belongs to }W_b\}.
\]
All rotations act on entry columns and leave lane indices fixed.

\begin{theorem}[The reflection source]\label{thm:reflection-source}
\formal{CapacityReflectionSource.lean}{CapacityReflection.finrank_rowSpace,CapacityReflection.finrank_source,CapacityReflection.source_disjoint_views,CapacityReflection.source_shared}
The source $C_b$ has binary dimension $bdm$. Its distinct rotated views
$T^cC_b$, $T^eC_b$ intersect only in zero. It therefore admits one total
deterministic shared decoder for every receiver, and in particular the same
function can serve the two nonempty classes $\{0\}$ and
$\mathbb Z/d\setminus\{0\}$.
\end{theorem}

\begin{outline}
The $m$ pairs $\{j,-j\}$ of nonzero coordinates are free $b$-vectors, so
$\dim W_b=bm$ and $\dim C_b=dbm$. The formal equivalence chooses the
smaller positive representative of each pair and treats coordinate zero
separately.

A vector in $\sigma^cW_b$ is invariant under reflection
$j\mapsto-j-2c$ and vanishes at $-c$. Two such symmetries with distinct
centers give invariance under translation by $2(e-c)\ne0$. Because $d$ is
prime, this translation generates the whole cycle. Thus a vector in two
distinct translates is constant, and its zero coordinate forces it to be
zero. Apply this independently to every row and all its lanes. The formal
proof obtains any displacement as a multiple of the nonzero step in the
field $\mathbb Z/d$.

For a nonzero received view, disjointness makes its receiver phase and
underlying source matrix unique. Return that source's requested row.
Return zero at the zero view and choose arbitrary values off the union of
the views. This is one well-defined total function. It is not asserted
linear or efficiently computable.
\end{outline}

\begin{theorem}[Every-width linear subcode barrier]\label{thm:reflection-linear-barrier}
\formal{CapacityReflectionBarrier.lean}{CapacityReflection.one_class_finrank_le,CapacityReflection.two_class_finrank_le,CapacityReflection.two_class_codimension_lower,CapacityReflection.two_class_positive_deletion,CapacityReflection.deletion_rate_limit}
Let $C'\le C_b$ be any binary-linear subcode; it may couple every row and
every lane. If it has a single binary-linear shared decoder, then
\[
 \dim_{\mathbb F_2}C'\le b m(m+1).
\]
If instead it has two binary-linear shared decoders for the original
classes $\{0\}$ and $\mathbb Z/d\setminus\{0\}$, then
\begin{equation}\label{eq:reflection-two-class}
 \dim C'\le b\bigl(m(m+1)+d\bigr),\qquad
 \operatorname{codim}_{C_b}C'\ge b(m^2-d).
\end{equation}
The deletion bound is positive for odd primes $d\ge7$ and is asymptotic
to $bd^2/4$.
More generally, these linear-subcode bounds hold over every field of
characteristic two and at every odd side length $d=2m+1$. Primeness is
needed for the preceding deterministic-decoder construction.
\end{theorem}

\begin{outline}
First suppose there is one linear decoder $L$. Extend the coefficient
field to a splitting field $E$ for a primitive $d$th root of unity.
The formal construction takes the $E$-span of the embedded original basis
of $C'$ and proves that basis remains independent. The extended code
therefore has the same dimension. Extending the actual coordinate
coefficients of $L$ preserves its shared-decoder equations, and span
induction preserves the literal reflection and zero-coordinate equations
defining $C_b$. All row and lane coordinates are included. This step
extends only linear equations;
it makes no scalar-extension claim about the deterministic decoder of
\cref{thm:reflection-source}.

Since $d$ is odd, its image in the characteristic-two field is nonzero;
the formal primitive-character construction supplies the required
extension. Choose the Fourier basis $e_a(j)=\zeta^{aj}$,
$a\in\mathbb Z/d$. Its Vandermonde matrix is invertible. Apply this
\emph{same} change of basis to both physical indices, preserving the
transpose swap, and leave the lanes unchanged. Let $\mathcal F$ be the
one-index synthesis and define the column mode
$M_j(v)_{r,c,\ell}=\zeta^{jc}(\mathcal Fv)_{r,\ell}$.
Writing $X$ as the two-index synthesis of $Z$, its actual rotated view is
\[
 T^cX=\sum_j\zeta^{jc}M_j(\operatorname{col}_jZ).
\]
Apply $L$ and compare the independent characters in the receiver $c$
with the demanded row of $X$. This gives
$\mathcal F(\operatorname{row}_jZ)=L(M_j(\operatorname{col}_jZ))$.
Consequently, with the fixed linear maps
$B_j=\mathcal F^{-1}LM_j$, agreement takes the form
\begin{equation}\label{eq:reflection-frequency-equations}
 \operatorname{row}_j(Z)^{\mathsf T}
       =B_j\operatorname{col}_j(Z),\qquad j\in\mathbb Z/d,
\end{equation}
where $B_j:E^{db}\to E^{db}$ may mix every row and lane. No frequency
equation or dimension certificate is supplied as a premise: the formal
proof derives these equations from the original decoder.

In the Fourier basis a reflected row has equal coefficients at $a,-a$:
reflection exchanges those characters, and their independence forces
coefficient equality. Evaluating the row at physical coordinate zero sums
all coefficients. Equal nonzero pairs cancel in characteristic two, so
the zero-coordinate condition forces coefficient zero to vanish.
Conversely these coefficient conditions imply both original conditions.
Consequently
\[
 \operatorname{col}_0(Z)=0,\qquad
 \operatorname{col}_a(Z)=\operatorname{col}_{-a}(Z).
\]
The further Fourier change on the row index preserves these column
identities. Equation~\eqref{eq:reflection-frequency-equations} at zero now
forces $\operatorname{row}_0(Z)=0$ as well.

Choose representatives $1,\ldots,m$ of the nonzero frequency pairs.
Group the remaining rows and columns into those pairs. A $2$-by-$2$
block has duplicate columns and is specified by
$v_{a,b}\in E^{2b}$, its two entries in column $b$. The equations for
frequencies $a$ and $-a$ use the same input column, and show that the
whole compressed row $(v_{a,b})_b$ is a fixed linear function of the
compressed column $(v_{c,a})_c$. Equivalently, for fixed linear maps
$L_{a,b,c}:E^{2b}\to E^{2b}$,
\[
 v_{a,b}=\sum_c L_{a,b,c}(v_{c,a}).
\]
All cross-lane coefficients remain in these maps.

Supply the strict upper blocks $v_{a,b}$, $a<b$, and all diagonal blocks.
There are $\binom m2+m$ blocks, hence $bm(m+1)$ field coordinates.
Recover rows in descending order. In column $a$, the entries above the
diagonal were supplied, the diagonal was supplied, and those below it
belong to already recovered rows. The whole column is known and its fixed
linear maps determine row $a$. Thus the supplied-coordinate projection is
injective on the agreement space, proving its dimension bound. The formal
proof reuses the stronger descending block-recovery theorem, charging each
diagonal kernel by at most $2b$.

Finally suppose $C'$ has decoders $L_0,L_1$ for the two original classes.
In the \emph{original binary coordinates} impose
$L_1(X)=\operatorname{row}_0(X)$. This is one linear map to a $bd$-dimensional
row, so its kernel cuts $C'$ by at most $bd$ dimensions. On the resulting
subcode $C''$, the single decoder $L_1$ works also at receiver zero; it
already works at every nonzero receiver. The one-class bound gives
$\dim C'\le\dim C''+bd\le bm(m+1)+bd$. Subtract from
$\dim C_b=bdm$ and use $d=2m+1$ to get the stated deletion bound.
At $m\ge3$, $m^2-(2m+1)>0$, which covers the odd primes $d\ge7$.
\end{outline}

The source rate is $(d-1)/(2d)$, approaching one half from below, and the
statement covers every $b$ including $b=1$ and the canonical width $b(d)$.
Prime side length is used for the disjointness of rotated source views;
characteristic two is used for the Fourier zero coefficient. The
two-class partition in the linear bound is specifically singleton versus
complement. The linear bound itself needs only odd side length. The
result supplies neither a paid implementation of the nonlinear decoder
nor a full-cube or native impossibility theorem.

\section{A bounded extension to a single-site update}
\label{sec:single-site}

The balance argument also handles the other single-site endpoint of
the masked-neighbor family.  This extension is included to specify what
the pre-release strengthening check establishes and what it leaves open.

\begin{proposition}[One binary neighbor update]\label{prop:single-site}
\formal{CompositionSingleSiteEndpoint.lean}{single_at_most_two_classes_impossible}
Let $d\ge3$ and let $L=I+E_{0,1}$ act on every binary row/lane word,
where $E_{i,j}x$ has its only possibly nonzero coordinate, $x_j$, at $i$.
For every positive width, there is no complete factorization
$\tau=KBLH$ in which $B,H$ are column-equivariant total maps and $K$
has at most two receiver classes.
\end{proposition}
\begin{outline}
Absorb $B$ into $K$.  We prove adjacent-pair transitivity for
$Q=LTL^{-1}$ and then apply \cref{thm:balanced-pair}.
Put $U_i=I+E_{i,i+1}$, with cyclic indices.  Since $L=U_0=L^{-1}$,
conjugating $\langle Q,TQT^{-1}\rangle$ by $L^{-1}$ gives
$\langle T,A\rangle$, where
\[
 A=U_1U_0U_{-1}U_0.
\]
For $d\ge4$, direct multiplication gives $A=I+N$, with
\[
 N=E_{1,2}+E_{-1,0}+E_{-1,1},\qquad
 N^2=E_{-1,2},\qquad N^3=0.
\]
In characteristic two, $A^2=I+E_{-1,2}$.  Conjugating by rotations
therefore gives every transvection $I+E_{i,i+3}$.

If $d=4$, the directed step-three edges form one cycle.  The elementary
identity
\[
 [I+E_{i,j},I+E_{j,k}]=I+E_{i,k}
 \qquad(i,j,k\text{ distinct})
\]
along simple directed paths yields every $I+E_{i,j}$ with $i\ne j$.
For $d\ge5$, use the already obtained $E=I+E_{0,3}$.
Here $Ae_0=e_0+e_{-1}$, while $e_3^*A^{-1}=e_3^*$, because row three
is unchanged by $A$ and its inverse.  Thus
\[
 (AEA^{-1})E=I+E_{-1,3}.
\]
Rotations also give every step-four transvection.  Within each orbit of
the step-three rotation, a simple directed path from $i$ to $i-3$
gives the reverse step-three transvection by repeated commutators.
If that orbit has two vertices, the reverse edge is already a generator.
Since $d\ge5$, the indices $i,i+4,i+1$ are distinct, and
\[
 [I+E_{i,i+4},I+E_{i+4,i+1}]=I+E_{i,i+1}.
\]
The resulting step-one cycle gives every elementary transvection by
the same chain argument.  Such transvections
act transitively on nonzero binary words: preserve one nonzero coordinate
and add it to clear every other nonzero coordinate, then rotate.

When $d=3$, the coincident indices require a separate calculation.
The same product is
$A(a,b,c)=(a,a+c,a+b+c)$.  It sends $100$ to $111$ and then to $101$.
Their rotations include the three unit vectors, the three weight-two
vectors and the all-one vector, proving transitivity on all seven
nonzero words.  This completes the adjacent-pair argument for every
$d\ge3$.

Hence the partition of $K$ must properly color the cycle.  Odd $d$
is impossible with two colors.  For even $d\ge4$ a proper two-coloring
has period two.  But $L(1,\ldots,1)=(0,1,\ldots,1)$, so no nonzero
column rotation preserves $L\operatorname{Fix}(T)$.  The period
constraint in \cref{thm:period-transport} rules out this last case.
\end{outline}

For clarity, this does not settle the family
$L_m=I+\operatorname{diag}(m)T$ with an arbitrary binary mask $m$.
If $m$ has at least one zero, the weighted shift is nilpotent and
$L_m^{-1}=\sum_{j=0}^{d-1}(\operatorname{diag}(m)T)^j$.
The all-one mask is singular.  A period $r$ of $m$ makes $L_m$ commute
with $T^r$ and gives the class bound of \cref{thm:periodic-first}.
The masks with exactly one zero and exactly one one are covered by
\cref{cor:anchored-all,prop:single-site}, respectively.  The remaining
aperiodic masks, and half-period masks with two final classes, are not
excluded by these arguments.  All statements concern complete algebraic
maps; no native implementation or physical cost bound is asserted.

\section{A complete scattered first factor and its failure to compose}
\label{sec:completion}

A source decoder does not automatically extend to a bijection on all
matrices.  For the scattered normal-basis source, such an extension does
exist.  The particular extension below nevertheless forces every receiver
of the remaining factor to have its own class.  Both assertions concern
the complete finite state space; neither supplies an efficient evaluator
for the orbit permutation.

\subsection{The source and a total permutation}

Let $d\ge2$ and $b\ge1$, put
$V=(\F_2^b)^{d\times d}$, and let
$B=(\F_2^b)^d$ be the whole-column alphabet, of cardinality
$q=2^{db}$.  Use the reversed normal-basis identification
$\F_2^d\cong E=\F_{2^d}$ underlying
\cref{thm:interval-attainment}; a column rotation by $s$ acts on a row
as $\Phi^s$, where $\Phi(x)=x^2$.
Let $\beta$ be its normal element.  Its $d$ conjugates are distinct,
and none equals $0$ or $1$.
For $u\in B$, write $a_\ell\in E$ for the field encoding of
$(u_{j,\ell})_{j\in\Z/d}$ and define $w(u)$ by its encoded rows:
\[
 w(u)_{0,\cdot,\ell}=a_\ell,\qquad
 w(u)_{i,\cdot,\ell}=\beta a_\ell\quad(i\ne0).
\]
Set $W=\{w(u):u\in B\}$.  Also write $\bar u$ for the matrix whose
every column is $u$, and let $z(u)$ have column zero equal to zero and
all other columns equal to $u$.

\begin{theorem}[A complete scattered first factor]
\label{thm:scattered-completion}
\formal{CompositionScatteredSource.lean}{scatteredSource_finrank,scattered_orbit_injective,scattered_punctured_orbits_disjoint}
\formal{CompositionScatteredCompletion.lean}{scatteredCompletion_equivariant,scatteredFirstFactor_classUniform,scatteredFirstFactor_source,scatteredFirstFactor_fixed_preimage}
\formal{CompositionSparsePatch.lean}{columnPatch_involutive,columnPatch_nonzero_column}
The binary-linear source $W$ has dimension $db$.  There are total
permutations $H,R,K_1$ of $V$ such that $H$ commutes with column
rotation, $R$ changes only column zero, and $K_1=RH$ is class-uniform
for the two classes $\{0\}$ and $(\Z/d)\setminus\{0\}$.
They satisfy
\[
 H(w(u))=z(u),\qquad K_1(w(u))=\bar u\quad(u\in B),\qquad
 K_1^{-1}(\operatorname{Fix}(T))=W.
\]
Here $R$ is an involution defined on every input.
\end{theorem}
\begin{outline}
The row-zero projection recovers $u$, so $w$ is an injective binary-linear
map and $\dim W=db$.  If
$T^sw(u)=T^tw(v)$ with $u,v\ne0$, choose a nonzero lane parameter.
The two row types give
\[
 \Phi^s(a)=\Phi^t(a')\ne0,\qquad
 \Phi^s(\beta)\Phi^s(a)=\Phi^t(\beta)\Phi^t(a').
\]
Cancellation and the distinct conjugates of $\beta$ give $s=t$, then
all lane parameters agree.  Thus $(s,u)\mapsto T^sw(u)$ is injective
for $u\ne0$.  Similarly $(s,u)\mapsto T^sz(u)$ is injective: its
unique zero column determines $s$, and any other column determines $u$.

These two orbit families are disjoint.  Suppose a rotated $w(u)$ equals
a rotated $z(v)$.  Choose a lane whose row-zero field parameter is
nonzero.  On the punctured side, a row in this lane is either zero or
the common punctured binary word, according as the corresponding bit of
$v$ is zero or one.  Both rows zero and one on the scattered side are
nonzero.  They must therefore be equal on the punctured side.  In field
coordinates this would give
$\Phi^s(\beta)\Phi^s(a)=\Phi^s(a)$ with $a\ne0$, contradicting
$\Phi^s(\beta)\ne1$.

Define $H$ to exchange $T^sw(u)$ and $T^sz(u)$ for every $s$ and
every $u\ne0$, and fix every remaining state.  Injectivity and
disjointness make these disjoint pairs, so this is an actual total
involution.  Rotation permutes the pairs with their orientation
preserved, hence $HT=TH$.  Zero lies outside both families and is fixed.
This specializes the completion to the present source and needs no
assumed equivariant extension.

Define $R$ to exchange $z(u)$ and $\bar u$ for each $u\ne0$, fixing
all other inputs.  Each pair agrees outside column zero, so $R$ leaves
those columns unchanged.  The pairs are disjoint: a punctured state has
zero column zero, whereas a nonzero constant state does not.  Thus $R$
is again a total involution and $Rz(u)=\bar u$, including $u=0$.

For the exceptional class use the decoder
$Y\mapsto\operatorname{col}_0(RH(Y))$; for the other class use
$Y\mapsto\operatorname{col}_0H(Y)$.
If $c\ne0$, equivariance gives
$\operatorname{col}_0H(T^cX)=\operatorname{col}_cH(X)
 =\operatorname{col}_cK_1(X)$.
The exceptional decoder gives column zero by definition.  This proves
sharedness on every input.  Finally the $q$ matrices $\bar u$ exhaust
$\operatorname{Fix}(T)$, and bijectivity of $K_1$ gives the stated
preimage equality.
\end{outline}

\begin{corollary}[The exceptional class is necessary]
\label{cor:scattered-one-class}
\formal{CompositionScatteredCompletion.lean}{scattered_one_class_impossible}
No injective equivariant map sends all of $W$ into
$\operatorname{Fix}(T)$.
\end{corollary}
\begin{outline}
Choose $u\ne0$. Equivariance and the fixed image imply
\[K(Tw(u))=TK(w(u))=K(w(u)).\]  Injectivity forces $Tw(u)=w(u)$,
contradicting the free-orbit calculation above.  A one-class complete
map is equivariant, so it is excluded as well.
\end{outline}

\subsection{Complete-input disagreement}

The agreement estimates in \cref{prop:agreement} have the following
robust form.  The comparison need not be a permutation.

\begin{proposition}[Robust comparison with an equivariant map]
\label{prop:robust-equivariant-comparison}
\formal{CompositionRobustAgreement.lean}{agreement_error_cover,robust_full_input_count,robust_repeated_class_count}
Let $A$ be any finite alphabet with at least two elements and put
$q=|A|^d$.  Let $K_1,H\colon A^{d\times d}\to A^{d\times d}$ be
total maps, with $H$ equivariant, and let $K_2$ be class-uniform for a
partition $p$ with $M$ nonempty classes of sizes $n_J$.  Under the
uniform law on the complete input space, put
\[
 \gamma=\Pr[K_2K_1(X)\ne\tau X],\qquad
 \epsilon=\Pr[K_1(X)\ne H(X)].
\]
Then
\[
 \gamma+\epsilon\ \ge\
 1-\frac Md-\frac{\sum_Jn_J(n_J-1)}{2dq}.
\]
If a receiver class is repeated, also
$\gamma+\epsilon\ge(1-q^{-1})/2$.
\end{proposition}
\begin{outline}
The map $K_2H$ has the same receiver partition, with decoder
$f_J\circ H$.  If $K_2K_1$ is correct at $X$ and $K_1(X)=H(X)$,
then $K_2H$ is correct there.  The entire input space is therefore
covered by the correct inputs of $K_2H$, the failures of $K_2K_1$,
and the disagreements between $K_1$ and $H$.  Sum their cardinalities
and apply the two bounds of \cref{prop:agreement} to $K_2H$.
The formal statements retain the integer counts and multiply through
by the positive denominators.  They use no injectivity of $H$ or $K_1$.
\end{outline}

\begin{theorem}[Sparse completion forces a fully separated second factor]
\label{thm:scattered-second-factor}
\formal{CompositionSparsePatch.lean}{columnPatch_support_card}
\formal{CompositionSparseObstruction.lean}{scatteredFirstFactor_disagreement,scattered_remaining_factor_singletons}
For the binary construction of \cref{thm:scattered-completion},
\[
 |\{X:R(X)\ne X\}|=2(q-1),\qquad
 \Pr[K_1(X)\ne H(X)]=\frac{2(q-1)}{q^d}.
\]
If $d\ge3$ and a class-uniform total map $K_2$ satisfies
$K_2K_1=\tau$ on all inputs, every receiver class of $K_2$ is a
singleton.  Thus its uniquely required remaining factor
$\tau K_1^{-1}$ needs $d$ classes.
\end{theorem}
\begin{outline}
The support of $R$ consists exactly of the two states $z(u),\bar u$
for every $u\ne0$.  Bijectivity of $H$ transports this support count
to the disagreement count for $RH$ and $H$.
If $K_2$ had a repeated class, the previous proposition with
$\gamma=0$ would force
\[
 \frac{2(q-1)}{q^d}\ge\frac{q-1}{2q},
 \qquad\text{hence}\qquad q^{d-1}\le4.
\]
But $q=2^{db}\ge8$ when $d\ge3$ and $b\ge1$.
The formal proof uses the equivalent integer inequalities, avoiding
subtraction or division of probabilities.  No exclusion at $d=2$ is
asserted.
\end{outline}

This disagreement counts complete input states.  It is different from
the number of output columns modified by a map, a linear error rank,
or a time charge.  A change confined to one output column can affect
almost every input.

\subsection{Endpoint classes are not multiplicative}

The need to retain actual factors is already visible in a short
bijective control.  For binary matrices define rowwise maps
\[
 (B_1X)_{i,c}=\begin{cases}X_{i,0}+X_{i,1},&c=0,\\X_{i,c},&c\ne0,\end{cases}
 \qquad
 (B_2X)_{i,c}=X_{i,c}+\sum_jX_{i,j}.
\]

\begin{proposition}[Two classes followed by one can require all classes]
\label{prop:endpoint-class-blowup}
\formal{CompositionEndpointBlowup.lean}{endpointFirst_involutive,endpointSecond_involutive,endpointFirst_classUniform,endpointSecond_equivariant,endpoint_composite_getter,endpoint_composite_requires_singletons}
For every even $d\ge2$, the maps $B_1,B_2$ are total involutions.
The first uses the singleton/rest partition and the second is equivariant.
Nevertheless every class-uniform representation of $B_2B_1$ has $d$
singleton classes.
\end{proposition}
\begin{outline}
The first map adds the unchanged column one twice when squared.
For the second, let $s_i=\sum_jX_{i,j}$.  The row sum after the map
is $s_i+d s_i=s_i$ because $d$ is even in $\F_2$, so its square is
also the identity.  The sum is rotation-invariant, proving that $B_2$
is equivariant; the two stated decoders for $B_1$ are
$Y\mapsto Y_{\cdot,0}+Y_{\cdot,1}$ and $Y\mapsto Y_{\cdot,0}$.

The normalized getter of the composite at receiver $c$ is
\[
 g_c(Y)_i=Y_{i,0}+\sum_jY_{i,j}
   +\begin{cases}0,&c=0,\\Y_{i,1-c},&c\ne0.\end{cases}
\]
These functions are pairwise distinct.  Comparing $c=0$ with $e\ne0$,
use a row supported at $1-e$.  Comparing nonzero distinct $c,e$, use
a row supported at $1-c$; the extra term at $1-e$ vanishes.
A shared class would make the corresponding normalized getters equal
on every input, a contradiction.  The formal proof derives this getter
equality from the complete class-uniform interface.
\end{outline}

\section{Balanced fibres beyond transitivity}
\label{sec:orbit-residues}

Transitivity in \cref{thm:balanced-pair} is a sufficient way to obtain
an orbit divisor.  It is not necessary.  The statements below use an
actual invariant scalar-word action and derive the tuple counts through
stabilizers.  They do not assume transitivity on tuples or independent
behaviour of the row and lane actions.

\subsection{The fixed-word residue test}

Let a group $U$ act on a finite scalar-word set $\Omega$, and let $I$
be a finite index set.  The action on $\Omega^I$ is diagonal.  Write
$h=|\operatorname{Fix}_U(\Omega)|$ and $r=|I|$.

\begin{theorem}[Fixed-word residues]
\label{thm:fixed-word-residue}
\formal{CompositionOrbitResidues.lean}{tuple_orbit_divisor,fixedTupleEquiv,invariant_fiber_fixed_residue,balanced_fixed_word_residue}
Suppose $m$ divides the cardinality of every non-singleton $U$-orbit
in $\Omega$.  Let $f\colon\Omega^I\to Z$ be $U$-invariant, where
$Z$ is finite, and suppose every fibre has the same cardinality $F$.
Then
\[
 |Z|\,(F\bmod m)\le h^r.
\]
In particular, in the complete matrix interface over a scalar alphabet
of size $a$, with $r=db$, $|Z|=a^{db}$ and
$F=a^{db(d-1)}$, any such invariant action satisfies
\begin{equation}\label{eq:fixed-word-residue}
 h^{db}\ \ge\ a^{db}\bigl(a^{db(d-1)}\bmod m\bigr).
\end{equation}
The integer $m$ need not be prime.
\end{theorem}
\begin{outline}
If a tuple $x$ is not fixed, some coordinate $x_i$ is not fixed.
The stabilizer containment $U_x\le U_{x_i}$ shows that
$|U\cdot x_i|$ divides $|U\cdot x|$.  Therefore every non-singleton
tuple orbit has size divisible by $m$.

An invariant fibre is a union of tuple orbits.  Its cardinality modulo
$m$ equals the number of its fixed tuples modulo $m$.  Consequently
each fibre contains at least $F\bmod m$ fixed tuples.  Fixed tuples
are precisely tuples of fixed scalar words, of which there are $h^r$.
Summing over the disjoint fibres proves the inequality.  In the complete
matrix setting, bijectivity gives the stated fibre size, as in the proof
of \cref{thm:balanced-pair}.
The formal argument proves the class equation for an arbitrary common
orbit divisor and constructs the equivalence between fixed tuples and
tuples of fixed words.
\end{outline}

\begin{corollary}[Too few fixed words]
\label{cor:few-fixed-words}
\formal{CompositionOrbitResidues.lean}{balanced_fixed_word_small_impossible}
Under the preceding hypotheses, suppose $r>0$, $m>1$,
$\gcd(m,a)=1$, $|Z|=a^r$, and $F=a^n$ for some $n\ge0$.
If $h<a$, no such invariant balanced function exists.
In particular this applies to any prime divisor $m$ coprime to $a$.
\end{corollary}
\begin{outline}
Coprimality implies $m\nmid a^n$, so $a^n\bmod m\ge1$.
The residue inequality gives $a^r\le h^r$, whereas $h<a$ and
$r>0$ give the strict reverse inequality.
\end{outline}

There is a useful version requiring only one fibre, without counting
the fixed set.

\begin{proposition}[A semiregular divisor]
\label{prop:semiregular-balance}
\formal{CompositionOrbitResidues.lean}{no_balanced_semiregular,no_balanced_free_off_zero}
Suppose $U$ fixes $0\in\Omega$ and every orbit of a word other than
zero has cardinality divisible by $m>1$, with $\gcd(m,a)=1$.
If $f\colon\Omega^I\to Z$ is invariant, no fibre avoiding the
all-zero tuple has cardinality $a^n$, for any $n\ge0$.
In particular, a cyclic action of order $m$ that is free away from zero
has this obstruction.  Hence it cannot be an invariance subgroup of a
nonconstant balanced decoder whose fibre sizes are powers of $a$.
\end{proposition}
\begin{outline}
Every tuple in the selected fibre has a nonzero coordinate.  Its
stabilizer is contained in that coordinate's stabilizer, so its orbit
size is a multiple of $m$.  The invariant fibre is a disjoint union of
these orbits and also has size divisible by $m$.  This contradicts
$\gcd(m,a^n)=1$.  For a free action of a finite cyclic group of order
$m$, each nonzero scalar orbit has size exactly $m$, giving the stated
specialization.  The proof uses this orbit-size condition directly; it
does not assume that all nonzero scalar words belong to one orbit.
\end{outline}

\subsection{An actual nontransitive order-five example}

Set $d=4$ and $\Omega=\F_2^4$, with cyclic shift
$(Tx)_i=x_{i+1}$.  Define the literal word operation
\[
 L(x_0,x_1,x_2,x_3)
   =(x_2,\ x_2+x_3,\ x_1+x_2,\ x_0+x_2),
 \qquad Q_c=T^cLTL^{-1}T^{-c}.
\]
All equalities in the next result concern these displayed operations;
no assumed scalar-group identification is used.

\begin{theorem}[Balance excludes a nontransitive adjacent pair]
\label{thm:order-five-example}
\formal{CompositionOrderFive.lean}{orderFivePerm_zero,orderFivePerm_fixed,orderFivePerm_order,orderFive_conjugate_difference,orderFive_pair_nontransitive,orderFive_no_balanced_covariance,orderFive_adjacent_class_impossible}
The operation $L$ is a zero-fixing permutation of order five, and fixes
no other word.  The group $\langle Q_0,Q_1\rangle$ is not transitive
on the fifteen nonzero words.
Nevertheless, for every $b\ge1$, a complete factorization
$KLH=\tau$ on $\F_2^{4\times4\times b}$, with $H$ equivariant and
$K$ class-uniform, must satisfy $p(0)\ne p(1)$.
The same scalar $L$ is applied independently to every row and lane;
$H$ and $K$ may be nonlinear and may mix all lanes.
\end{theorem}
\begin{outline}
Substitution gives
\[
 L^4x=(x_0+x_3,\ x_0+x_2,\ x_0,\ x_0+x_1),
 \qquad L(L^4x)=L^4(Lx)=x.
\]
The fixed-point equations give successively $x_0=x_2$, $x_1=0$,
$x_3=0$, and $x_2=0$.  Thus $L$ fixes only zero, and since
$L^5=I$ and $L\ne I$, its order is five.  The actual conjugates are
\[
 \begin{aligned}
 Q_0x&=(x_0+x_1,\ x_1+x_3,\ x_1,\ x_1+x_2),\\
 Q_1x&=(x_0+x_2,\ x_0,\ x_0+x_1,\ x_0+x_3).
 \end{aligned}
\]
They satisfy
\[
 Q_0Q_1^{-1}=L.
\]
Write $e_i$ for the coordinate unit words and $v=(1,1,1,1)$.
On $\{e_0,e_1,e_2,e_3,v\}$, the first conjugate fixes $e_0$
and cycles $(e_1,v,e_2,e_3)$; the second fixes $e_3$ and cycles
$(e_0,v,e_1,e_2)$.  Thus both preserve this five-word set.  In particular no group element sends
a unit word to a word with exactly two nonzero coordinates, so the pair
group is not transitive on the nonzero words.  Lean checks these finite
identities on the sixteen literal words in its kernel.

If $p(0)=p(1)$, the common shared decoder $f$ has
$f(Q_0Y)=\sigma f(Y)=f(Q_1Y)$.  Substitute $Q_1^{-1}Y$ to obtain
$f(LY)=f(Y)$, and hence invariance under the cyclic group generated by
$L$.  Because $L$ has prime order five and fixes only zero, every
nonzero scalar orbit has size five.  Stabilizer containment therefore
makes every nonzero tuple orbit size divisible by five.

Exactness on the complete finite space makes $H$ and $K$ bijective.
The column decoder is balanced: each value has exactly
$2^{12b}$ preimages in the $4b$-tuple domain.
Choose an output value different from $f(0)$; it exists because the
column alphabet has $2^{4b}>1$ values.  Its fibre avoids zero and is
invariant under $L$, so the semiregular-divisor proposition says its
size is divisible by five.  The power $2^{12b}$ is not, a contradiction.
The formal endpoint derives covariance and balance from the full product
identity; neither property is supplied as a separate certificate.
\end{outline}

These tests are necessary conditions for an actual complete factorization.
Passing an orbit-assignment test for one shared function does not produce
jointly independent output columns or a bijective whole-matrix map.
The arguments also impose no cost bound on evaluating a scalar-word
permutation or an orbit representative.

\subsection{Field multipliers and two scope checks}

\begin{proposition}[Multiplier class costs]\label{prop:multiplier-classes}
\formal{CompositionMultiplierClasses.lean}{fieldMultiplierConjugate_eq,multiplier_class_residue,multiplier_class_count,multiplier_class_count_prefix}
\formal{CompositionMultiplierGetters.lean}{actual_multiplier_getters_card,multiplier_degree_le,multiplier_binary_degree_eq}
In the balanced-intertwiner interface, identify scalar words over $\F_q$
with $E=\F_{q^d}$ in a reversed normal basis, so rotation is
$\Phi(x)=x^q$. Let $L(x)=\lambda x$, $\lambda\ne0$, on each row and lane.
Put $\kappa=\lambda^{1-q}$, $h=[\F_q(\kappa):\F_q]$ and
$h_0=[\F_q(\lambda):\F_q]$. Every second-layer class is contained in
one residue class modulo $h$, so it requires at least $h$ classes.
More generally, among the first $n$ receivers there are at least
$\min(n,h)$ distinct class labels, for every $0\le n\le d$.
The complete operation $L$ itself requires exactly $h_0$ distinct class
getters, and $h\le h_0$. For $q=2$, $h=h_0$.
\end{proposition}
\begin{outline}
Writing $m_z$ for scalar multiplication, direct composition gives
\[
 Q_c=m_{\Phi^c(\kappa)}\Phi,\qquad
 Q_cQ_e^{-1}=m_{\Phi^c(\kappa)/\Phi^e(\kappa)}.
\]
If $c,e$ share a decoder, its covariance makes it invariant under the
second map. A multiplier $\mu\ne1$ has order $t>1$ dividing $q^d-1$;
its cyclic action is free away from zero. Since $\gcd(t,q)=1$,
\cref{prop:semiregular-balance} contradicts the required balance. Thus
$\Phi^c(\kappa)=\Phi^e(\kappa)$, equivalent to $c\equiv e\pmod h$.
The first $\min(n,h)$ receivers have distinct residues, proving the prefix
bound without requiring a full extension-degree set of receivers.

Let $\ell_0:E\to\F_q$ extract coordinate zero. The unique full-domain
getter for $L$ at receiver $c$ is, on each row and lane,
$x\mapsto\ell_0(\Phi^c(\lambda)x)$: rotate the received word back,
apply $L$, then rotate its demanded coordinate to zero. Two such maps
agree exactly when their scalar multipliers agree, since a nonzero
scalar difference is invertible and $\ell_0\ne0$. There are $h_0$
conjugates, giving both necessity and sufficiency. The inclusion
$\kappa\in\F_q(\lambda)$ gives $h\le h_0$; for $q=2$,
$\kappa=\lambda^{-1}$ gives equality. Thus this family does not hide a
new small-first-layer exclusion. For odd $q$, a nonzero trace-zero
$\lambda\in\F_{q^2}$ has $h_0=2$ but $\kappa=-1$ and $h=1$,
so the distinction is necessary.
\end{outline}

\begin{example}[Surjectivity is weaker than balance]\label{ex:ratio-unbalanced}
\formal{CompositionRatio.lean}{ratioDecoder_surjective,ratio_nonzero_fiber_card,ratio_zero_fiber_card,ratio_eight_fiber_cards}
\formal{CompositionMultiplierClasses.lean}{ratio_multiplier_conjugate_covariance}
For $E=\F_8$, $d=3$, $b=1$ and any $\lambda\ne0$, use the preceding
actual conjugates. On the complete row-word input $E^3$, the total map
\[
 f(x_0,x_1,x_2)=
 \begin{cases}x_1/x_0&x_0\ne0,\\0&x_0=0\end{cases}
\]
is surjective and satisfies $f(Q_cX)=\Phi(f(X))$ for every receiver.
Its seven nonzero fibres have size $56$ and its zero fibre has size $120$,
rather than the required common size $64$. More generally, over a field
of size $Q$, the corresponding fibre sizes are $(Q-1)Q$ and
$Q^2+(Q-1)Q$.
\end{example}
\begin{outline}
Each $Q_c$ applies Frobenius and the same nonzero multiplier to all three
rows, so the common multiplier cancels in the quotient; the zero branch
is preserved. The input $(1,z,0)$ attains any $z$. For nonzero $z$ there
are seven choices of $x_0$ and eight of $x_2$, with $x_1=zx_0$ fixed.
For $z=0$, add the $64$ inputs with $x_0=0$ to the $56$ with
$x_0\ne0,x_1=0$. These sum to $512$. No complete bijective transposer
or balanced decoder is constructed by this example.
\end{outline}

\begin{proposition}[Exact orbit-assignment criterion]\label{prop:orbit-assignment}
\formal{CompositionOrbitAssignment.lean}{OrbitAssignment.orbit_map_exists_iff,OrbitAssignment.allowed_iff_divides,OrbitAssignment.cyclic_balanced_map_iff_orbit_assignment}
Let a finite group $G$ act on a finite source set and let
$\theta:G\twoheadrightarrow\mathbb Z/d$ specify the cyclic action on a
finite target set. For each source orbit $O$, choose $x\in O$ and put
$n_O=|O|$ and $t_O=[\mathbb Z/d:\theta(G_x)]$.
For each individual target orbit $D$, put $\ell_D=|D|$.
An equivariant map with every target fibre of size $F$ exists if and only
if the whole source orbits can be assigned to individual target orbits,
with $\ell_D\mid t_O$ for every assignment and
\[
 \sum_{O\ \mathrm{assigned\ to}\ D}\frac{n_O}{\ell_D}=F
 \quad\text{for each }D.
\]
\end{proposition}
\begin{outline}
A map from the transitive orbit $O$ to $D$ is determined by the image $z$
of $x$. It is well defined precisely when
$\theta(G_x)\subseteq(\mathbb Z/d)_z$, equivalently
$\ell_D\mid t_O$ by the subgroup structure of a finite cyclic group.
When it exists, the map $gx\mapsto\theta(g)z$ is onto $D$ and every
fibre has size $n_O/\ell_D$: group translation bijects any two fibres,
and counting their disjoint union gives the quotient. The phase
stabilizer is independent of the source representative because the
phase group is abelian. Restricting a balanced map to each source
orbit proves necessity. Conversely choose one such orbit map for every
assignment; their disjoint union has the stipulated fibre sizes.
Whole source orbits cannot be split fractionally. This is an exact test
for one class function, not for joint bijectivity of all output columns.
\end{outline}

\section{Nonlinear corrections to a complete affine baseline}\label{sec:residuals}

An operator bound becomes useful for a computation only after its
\emph{complete} baseline is exhibited. This appendix gives that conditional
interface. It allows arbitrary nonlinear corrections and arbitrary correctness
sets, but supplies no cheap affine approximation to an arbitrary machine.
Fixed-advice fibre counting is already explicit in
\cite[Theorem~5.5]{ClassUniformMemo2026} and
\cite[Corollary~6.16]{SeventeenthPaper2026}. We use the same counting principle
with a chronological computation and retain the actual propagated correction
span. The distinction prevents charging one correction again at each consumer.

\subsection{The complete baseline and its joint record}

Let $F$ be a finite field of size $q$, and let $X\in F^N$ be the raw input.
All current payload, retained control and earlier output coordinates belong
to the state. A finite sequence of state spaces can be embedded in their
direct sum, so we use one finite-dimensional state space $V$. The baseline
initialization, updates and readout are global affine maps
\[
 v_0=JX+j,\qquad G_i(v)=L_i v+g_i,\qquad \operatorname{out}=Ov_m+o.
\]
Let the complete baseline output be
$K_0X+c_0$. At a marked update the actual map may be arbitrary. Its residual
on the actual old state is $E_i u_i(X)$, where $E_i$ is a fixed linear
insertion map. At other updates the residual is zero. The initialization must
also agree with the baseline, or its departure must be included as a marked
initial residual. All unmarked maps, including final readout, remain the same
\emph{global} affine maps on every reached state.

Choose one preimage $u_i(X)$ in a fixed linear section of $\operatorname{im}E_i$.
These values need not be independent or linear in $X$. Write $Z(X)$ for their
one joint tuple, and put
\[
 A=\tau-K_0,\qquad D=\dim\ker A,\qquad
 W=\sum_i\operatorname{im}(OL_{m-1}\cdots L_{i+1}E_i),
\]
where the empty product is the identity. These are comparison records and
subspaces; the definition does not give a native machine free access to them.

\begin{theorem}[Joint residual replay]\label{thm:residual-replay}
\formal{ResidualPropagation.lean}{residual_readout_sum,actual_readout_sum,actual_readout_insertions}
The complete-program counting endpoint is
\lean{actual_correct_card_le} in \texttt{ResidualProgram.lean}.
For every fixed tuple $z=(u_i)$, replay gives
\[
 \operatorname{out}_z(X)=K_0X+c_0+
       \sum_i OL_{m-1}\cdots L_{i+1}E_i u_i.
\]
On an arbitrary correctness set $C$, the actual tuple $Z(X)$ reproduces the
actual computation. Consequently
\begin{equation}\label{eq:residual-record}
 |C|\le |Z(C)|q^D.
\end{equation}
\end{theorem}

\begin{outline}
Induct in the original update order. After fixing the residual values, each
update is $v\mapsto L_iv+g_i+E_i u_i$, so its linear part is still $L_i$.
Expanding the recurrence propagates each inserted residual through exactly
the subsequent linear maps and then through $O$. With $z=Z(X)$ the
recurrence agrees with every actual update on $X$, including retained
outputs. On a fixed record fibre, correctness is the affine equation
$AX=c(z)$. Subtract one solution, if there is one, to inject the fibre into
$\ker A$. It therefore contains at most $q^D$ elements. Sum over the
\emph{attained} joint records. The formal proofs perform this recurrence and
translate-to-kernel injection; they do not assume that $C$ is a subspace.
\end{outline}

\begin{theorem}[Joint output-span bound]\label{thm:residual-span}
\formal{ResidualSpan.lean}{residualSpanDescriptor_injective,finrank_preimage_le_kernel_add_intersection,residual_span_card_le}
\formal{ResidualOutputSpan.lean}{actual_correct_span_card_le}
For the actual program, and more generally, if $C\subseteq F^N$ satisfies $AX-c\in W$ for every $X\in C$,
then
\begin{equation}\label{eq:residual-span}
 |C|\le q^{D+\dim(W\cap\operatorname{im}A)}
       \le q^{D+\dim W}.
\end{equation}
For the replay above, take $c=c_0$ and the displayed propagated span $W$.
\end{theorem}

\begin{outline}
Subtract one point of $C$ to inject it into $A^{-1}(W)$. Choose a linear
projection onto $\ker A$ that is the identity on that kernel. The map sending
$x\in A^{-1}(W)$ to its projection and $Ax$ is injective, with values in
$\ker A\times(W\cap\operatorname{im}A)$. This proves the dimension bound
and then the finite-field cardinality bound. In particular, the invoice
uses the joint propagated span, not a sum that charges dependent copies
repeatedly. Its dimension is at most the sum of the insertion ranks.
\end{outline}

\begin{corollary}[Information and sparse corrections]\label{cor:residual-information}
\formal{ResidualEntropy.lean}{residual_entropy_lower,residual_uniform_entropy_lower,actual_correct_entropy_lower,residual_span_information}
The sparse-count endpoint is \lean{actual_sparse_card_le} in
\texttt{ResidualSparse.lean}.
For any law supported on $C$, with entropy in base $q$,
\[
 H_q(Z)\ge H_q(X)-D.
\]
If $C$ is uniform and $\Delta=N-\log_q|C|$, this gives
$H_q(Z)\ge N-D-\Delta$ and $\dim W\ge N-D-\Delta$.
If each of $m$ public insertion ports has rank at most $r$, the correctly
handled inputs on which at most $s$ residuals are nonzero number at most
\begin{equation}\label{eq:sparse-residual-menu}
 q^D\sum_{j=0}^{\min(s,m)}\binom mj(q^r-1)^j.
\end{equation}
\end{corollary}

\begin{outline}
Each conditional input fibre has at most $q^D$ points, so
$H_q(X\mid Z)\le D$. Since $Z$ is a function of $X$, the entropy chain rule
gives the first inequality. The span bound gives the second assertion for
uniform $C$. For the sparse count, choose coordinates in each insertion
image, with the zero residual represented uniquely by zero. A support of
$j$ positions allows at most $(q^r-1)^j$ nonzero value tuples. Sum over
supports and apply \cref{eq:residual-record}. The locations may depend
arbitrarily on $X$; their support count is already included. The formal
count records each actual residual as a value in its insertion image,
so it charges the insertion rank and does not assume a particular
preimage selection. The shared support/value encoding also permits
different alphabets at different ports.
\end{outline}

For fixed $q,r$, $m=O(N)$ and $s=o(N)$, the logarithm of the sum in
\cref{eq:sparse-residual-menu} is $o(N)$. The formal endpoint is
\lean{LeanMachine.sparseMenu_log_div_tendsto_zero}. Its finite proof multiplies
the sum by $x^s$ and bounds the result by $(1+Ax)^m$ for $0<x\le1$,
where $A=q^r-1$.
Choose $x$ small enough that $m\log(1+Ax)$ costs at most half the desired
linear budget, then bound $s\log(1/x)$ by the other half. This treats
$s=0$ directly and requires no lower bound on $m/N$.

\subsection{What can be substituted, and what remains missing}

Suppose $N=bd^2$ and the \emph{complete} baseline has its column factors in
one common triangularizable algebra. The flag bound and its entry-width
extension give $D\le bd(d+1)/2$\cite[corollary ``One triangularizable side''
and remark ``Entries of $b$ bits'']{SeventeenthPaper2026}.
Thus correct inputs of deficit $\Delta$ require
\[
 H_q(Z)\ge \frac{bd(d-1)}2-\Delta,\qquad
 \dim W\ge \frac{bd(d-1)}2-\Delta.
\]
An independently established sharper operator bound may replace this $D$.
With fixed-width insertion events, sublinear nonzero corrections therefore
cover at most $q^{(1/2+o(1))N}$ inputs. A whole-row correction of rank $O(db)$
is not a fixed-width event.

The revised entry-rule bound of
\cite[Proposition~6.15 and Corollary~6.16]{SeventeenthPaper2026} is another
available specialization when its actual entry functionals are supplied.
It prices advice by its support and changes of an entry decoder by half a
symbol. A change of a whole row is $d$ entry changes. Additive advice alone
need not change any decoder. Our rank invoice and that change count are
different statistics; neither may be substituted for the other without
checking the complete rule.

An input-dependent branch or schema can be included in the joint record.
Then each branch has its own complete baseline and kernel bound $D_\theta$;
the support estimate sums the corresponding $q^{D_\theta}$ fibre bounds,
and entropy gives $H_q(Z)\ge H_q(X)-\mathbb E D_\theta$.
These attained-branch and entropy bounds are
\lean{residual_branch_card_le} and
\lean{residual_branch_entropy_lower} in \texttt{ResidualBranches.lean}. Unknown earlier
encodings, address histories, mutable flags or nonlinear unmarked readout
must be included in such a proved interface. They cannot be omitted by
checking only the zero-residual trajectory.

These statements exclude sparse corrections to an \emph{exhibited}
structured computation. They do not show that every fast transposer is
close to one, or that producing a correction of large rank costs superlinear
time. Supplying that structural comparison with a small joint record
remains the construction/impossibility bridge.

\subsection{Two sharp complete-baseline examples}

\begin{proposition}[Separable baseline defect]\label{prop:separable-defect}
\formal{ResidualSeparable.lean}{Separable.rank_defect_lower}
\formal{ResidualSeparableSharp.lean}{Separable.initialProjection_rank,Separable.rank_defect_sharp}
For $L(X)=AXB^{\mathsf T}$ over any field, put
$r=\min(\rank A,\rank B)$. Then
\[
 \rank(\tau-L)\ge d^2-\frac{r(r+1)}2.
\]
Equality is attained for $A=B=\operatorname{diag}(I_r,0)$ in every
characteristic. Thus this is a sharp bound for one separable baseline,
not for a sum of arbitrary tensor terms.
\end{proposition}
\begin{outline}
Writing $Z=X^{\mathsf T}$ turns agreement into $Z=AZ^{\mathsf T}B^{\mathsf T}$.
If $a=\rank A$, choose invertible $U,V$ with $UAV^{-1}=E_a$.
For $Y=UZV^{\mathsf T}$ and $D=VBU^{-1}$ the equation is
$Y=E_aY^{\mathsf T}D^{\mathsf T}$. The last $d-a$ rows vanish, and the
upper-left $a$-square block $Q$ determines the others and satisfies
$Q=Q^{\mathsf T}D_{11}^{\mathsf T}$. Its lower triangle, including the
diagonal, determines $Q$: it gives column zero and therefore row zero;
the recovered rows and the lower triangle then give the next column and
row, successively. The formal proof obtains the invertible factors and
the active coordinate reindexing from the actual matrix rank normal form.
The resulting restriction is a linear injection into the unordered-pair
coordinates of the active block, so the kernel has dimension at most
$a(a+1)/2$. Transposing the
agreement equation exchanges $A,B$ and gives the same bound with
$a=\rank B$. Rank--nullity proves the claim. For $A=B=E_r$, agreement
is exactly symmetry inside the leading $r$-square block and zero outside,
with $r(r+1)/2$ independent entries.
\end{outline}

\begin{proposition}[Nonlinear row-then-column agreement]\label{prop:row-column-agreement}
\formal{ResidualRowColumn.lean}{RowColumn.agreement_card_le,RowColumn.helper_agreement_card_le,RowColumn.identity_agreement_card}
\formal{ResidualRowColumnSharp.lean}{RowColumn.exists_sharp_helper,RowColumn.helper_bound_attained}
Over an alphabet of size $q\ge2$, first apply an arbitrary fixed map to
each row, then an arbitrary fixed map to each column. The resulting
complete map agrees with transpose on at most $q^{d(d+1)/2}$ inputs.
If one joint helper chooses all these maps, at most
$|\operatorname{range}H|q^{d(d+1)/2}$ inputs can be correct.
Both bounds are sharp as abstract counting statements.
\end{proposition}
\begin{outline}
Write the intermediate matrix as $Y_{j,:}=R_j(X_{j,:})$ and the output
column as $C_j(Y_{:,j})$. On agreement inputs,
$X_{j,:}=C_j(Y_{:,j})$ and therefore
$Y_{j,:}=R_j(C_j(Y_{:,j}))$. The lower triangle of $Y$ determines column
zero, hence row zero; then column one, hence row one; continue in order.
It determines all of $Y$, which in turn determines $X$. This injection
proves the first count without assuming either family is linear or
injective. Fix the helper and sum to prove the second count. Identity
row/column maps attain the first bound on symmetric inputs. For the
second, identify the alphabet with the cyclic group of order $q$ and
let the helper give $X_{ij}-X_{ji}$ for $i<j$. Identity row maps followed
by column maps adding or subtracting the corresponding differences
transpose every input. Every helper tuple is attained and there are
$q^{d(d-1)/2}$ tuples. Computing and delivering this helper is not free
in a native model.
\end{outline}

\section{Formalization boundary}\label{sec:formalization}

The library uses the earlier shared-decoder models and Mathlib, with the
repository's existing pins. Its entry type is a module, so vector lanes may
be mixed unless a theorem explicitly imposes additional structure. Arbitrary
deterministic shared rules and linear shared rules are separate predicates.
Complete factorization theorems quantify over total maps on the full state
space. Bijectivity follows from exactness and finite dimension for linear
factors, and from exactness and a finite state space for nonlinear factors.

The finite capacity proof proceeds through actual translated source spaces,
their component quotient, and cyclic run counts. Its arithmetic averaging
lemma is separate from those structural statements. The composition proof
similarly derives its forced row-class subspace from an actual factorization;
the common-message inequality is not an assumed surrogate for that step.
Finite progression extraction reuses Mathlib's Hales--Jewett theorem. The
small two-color certificate is evaluated by the kernel, without native
evaluation. Some finite checks use a larger heartbeat allowance.

All numbered mathematical statements, except the native implementation theorem,
are now checked. This includes exact arbitrary-set capacities and their cyclic
quotient counts, line-marker constructions, every structured and probabilistic
mask estimate, the actual scalar-extension/Fourier bridge for the reflection
barrier, multiplier and whole-orbit diagnostics, and sharp separable and nonlinear
row-then-column examples. The reflected linear-subcode bound holds at every odd
dimension; primality is retained only where the separate construction needs it.

The finite probability library constructs the actual uniform source law, proves
its bounded-differences inequalities and computes the required event expectations
and coordinate sensitivities. Conditioning, unlabelled partitions and simultaneous
run-spectrum events are included. The residual library evaluates the complete
chronological program, propagates its literal deviations, and proves joint
record, entropy, span and input-dependent sparse-support bounds. These statements
retain correlation between records and do not assume that a correctness source
is a subspace. The generated index lists the exact endpoints. Finite seed
coefficients are checked by the Lean kernel; the external generator is not trusted.

The native implementation has a distinct large obligation: implementing the
two-tape protocol as a transition table in the frozen model, with a complete
execution invariant and uniform time and state accounting. Its mathematical
protocol is proved in the text, but is not certified by an algebraic code or a
finite dimension calculation. That compiler is recorded as a substantial
formalization todo in the README. No additional axiom stands in for it.

\section{Index of Lean declarations}\label{sec:leanindex}

The following declarations are cited by formal-status lines. Names are relative to
\texttt{SharedDecoders}. A partial status retains the gap stated at its theorem.

\paragraph{\texttt{CapacityAmbient.lean}}
{\small\raggedright \lean{IsSharedCode.component_finrank_le}, \lean{covarianceSpace_rot}, \lean{finrank_covarianceSpace}.\par}

\paragraph{\texttt{CapacityClosureAsymptotic.lean}}
{\small\raggedright \lean{CapacityClosure.capacity_bound}, \lean{CapacityClosure.error_ratio_tendsto}, \lean{CapacityClosure.residual_finite_bound}, \lean{CapacityClosure.rounds_eq_floor}, \lean{CapacityClosure.uniform_eventual_bound}.\par}

\paragraph{\texttt{CapacityClosureControl.lean}}
{\small\raggedright \lean{CapacityClosure.scalarHull_supportedSymmetric}, \lean{CapacityClosure.symmetric_copy_defect}, \lean{CapacityClosure.symmetric_isSharedCode}.\par}

\paragraph{\texttt{CapacityCompletion.lean}}
{\small\raggedright \lean{IsSharedCode.completedCode}, \lean{completedCode_eq_of_range_le}, \lean{finrank_completedCode_add_rank}, \lean{finrank_le_completedCode}.\par}

\paragraph{\texttt{CapacityControls.lean}}
{\small\raggedright \lean{CapacityControls.exists_deterministic_not_linear}.\par}

\paragraph{\texttt{CapacityDifferenceCover.lean}}
{\small\raggedright \lean{IsSharedCode.difference_cover_real_bound}, \lean{IsSharedCode.saturated_cover_bound}, \lean{IsSharedCode.square_cover_bound}.\par}

\paragraph{\texttt{CapacityDigitParityEnvelope.lean}}
{\small\raggedright \lean{CapacityDigitParity.digit_parity_envelope}, \lean{CapacityDigitParity.digit_parity_floor_envelope}.\par}

\paragraph{\texttt{CapacityFieldLinearComponents.lean}}
{\small\raggedright \lean{CapacityFieldLinear.component_of_scalarTest}, \lean{CapacityFieldLinear.component_zero_of_scalarTest_zero}, \lean{CapacityFieldLinear.components_unique}, \lean{CapacityFieldLinear.phaseBasis}, \lean{CapacityFieldLinear.sum_components}.\par}

\paragraph{\texttt{CapacityFieldLinearCorollaries.lean}}
{\small\raggedright \lean{CapacityFieldLinear.one_class_original_bound}, \lean{CapacityFieldLinear.one_class_residual_eq_bot}, \lean{CapacityFieldLinear.two_class_half_bound}, \lean{CapacityFieldLinear.two_class_original_bound}.\par}

\paragraph{\texttt{CapacityFieldLinearResidual.lean}}
{\small\raggedright \lean{CapacityFieldLinear.active_class_residual_bound}, \lean{CapacityFieldLinear.finrank_encodedResidual}, \lean{CapacityFieldLinear.original_source_class_bound}.\par}

\paragraph{\texttt{CapacityIntervals.lean}}
{\small\raggedright \lean{CapacityIntervals.exists_binary_code_dimension}, \lean{CapacityIntervals.exists_optimal_interval_code}, \lean{IsSharedCode.two_mul_finrank_le_interval}.\par}

\paragraph{\texttt{CapacityIteration.lean}}
{\small\raggedright \lean{IsSharedCode.iterated_finrank_le}, \lean{componentCount_zero}.\par}

\paragraph{\texttt{CapacityLineDeterministic.lean}}
{\small\raggedright \lean{IsDeterministicLineCode.affine_two_mul_finrank_le}, \lean{IsDeterministicLineCode.two_mul_finrank_le}.\par}

\paragraph{\texttt{CapacityLineMarker.lean}}
{\small\raggedright \lean{LineMarker.affineMarker_finrank}, \lean{LineMarker.code_card}, \lean{LineMarker.exists_decoder}.\par}

\paragraph{\texttt{CapacityLineMarkerRate.lean}}
{\small\raggedright \lean{LineMarker.code_density}, \lean{LineMarker.code_density_tendsto}.\par}

\paragraph{\texttt{CapacityLinearization.lean}}
{\small\raggedright \lean{IsLinearSharedCode.inr_comap}, \lean{linear_subcode_lane_bound}.\par}

\paragraph{\texttt{CapacityMarker.lean}}
{\small\raggedright \lean{CapacityMarker.exists_affine_source_decoder}.\par}

\paragraph{\texttt{CapacityMaskComparisonProbability.lean}}
{\small\raggedright \lean{CapacityMaskProbability.ChangeGood.certificate_floor}, \lean{CapacityMaskProbability.EqualityGood.direct_gain_inverse_sqrt}, \lean{CapacityMaskProbability.component_floor_probability}, \lean{CapacityMaskProbability.direct_floor_probability}.\par}

\paragraph{\texttt{CapacityMaskFinite.lean}}
{\small\raggedright \lean{IsSharedCode.four_mul_finrank_le}, \lean{IsSharedCode.six_mul_finrank_le}, \lean{finite_weighted_run_budget}.\par}

\paragraph{\texttt{CapacityMaskMethodEndpoints.lean}}
{\small\raggedright \lean{CapacityMaskProbability.IsolationGood.rounded_floor}, \lean{CapacityMaskProbability.isolation_rate_probability}.\par}

\paragraph{\texttt{CapacityMaskRunProbability.lean}}
{\small\raggedright \lean{CapacityMaskProbability.SpectrumGood.source_bound}, \lean{CapacityMaskProbability.spectrum_bad_probability}.\par}

\paragraph{\texttt{CapacityMaskRunSpectrum.lean}}
{\small\raggedright \lean{IsSharedCode.run_spectrum_three}.\par}

\paragraph{\texttt{CapacityMaskStructured.lean}}
{\small\raggedright \lean{CapacityMaskStructured.first_component_count}, \lean{CapacityMaskStructured.second_component_count}, \lean{CapacityMaskStructured.structured_finrank_le}.\par}

\paragraph{\texttt{CapacityOptimization.lean}}
{\small\raggedright \lean{capacity_transversal_deficiency}, \lean{linear_capacity_eq_residual}, \lean{optimal_transversal_full_rank}, \lean{shared_capacity_eq_residual}.\par}

\paragraph{\texttt{CapacityReflectionBarrier.lean}}
{\small\raggedright \lean{CapacityReflection.deletion_rate_limit}, \lean{CapacityReflection.one_class_finrank_le}, \lean{CapacityReflection.two_class_codimension_lower}, \lean{CapacityReflection.two_class_finrank_le}, \lean{CapacityReflection.two_class_positive_deletion}.\par}

\paragraph{\texttt{CapacityReflectionSource.lean}}
{\small\raggedright \lean{CapacityReflection.finrank_rowSpace}, \lean{CapacityReflection.finrank_source}, \lean{CapacityReflection.source_disjoint_views}, \lean{CapacityReflection.source_shared}.\par}

\paragraph{\texttt{CapacityRetentionAverage.lean}}
{\small\raggedright \lean{IsSharedCode.retention_average}, \lean{IsSharedCode.retention_normalized}.\par}

\paragraph{\texttt{CapacityRetentionControl.lean}}
{\small\raggedright \lean{RetentionControl.laneSwap_eq_transpose_iff}, \lean{RetentionControl.paired_source_rate_and_loss}.\par}

\paragraph{\texttt{CapacityRetentionCoprime.lean}}
{\small\raggedright \lean{IsSharedCode.eq_symmetric_of_at_most_one_exception}, \lean{IsSharedCode.full_symmetric_retention}, \lean{IsSharedCode.prime_full_symmetric_retention}.\par}

\paragraph{\texttt{CapacityRetentionFinite.lean}}
{\small\raggedright \lean{IsSharedCode.retention_bound}.\par}

\paragraph{\texttt{CapacityRetentionStability.lean}}
{\small\raggedright \lean{RetentionAverage.fixed_gain_requires_loss}, \lean{RetentionAverage.stability}, \lean{RetentionAverage.stability_bounded_classes}.\par}

\paragraph{\texttt{CapacityRuns.lean}}
{\small\raggedright \lean{CapacityRuns.changeCount_double}, \lean{CapacityRuns.weighted_run_budget}.\par}

\paragraph{\texttt{CapacityRunsComponents.lean}}
{\small\raggedright \lean{componentCount_le_changes_add_gcd}, \lean{componentCount_le_runs_add_two_gcd}.\par}

\paragraph{\texttt{CapacitySetCounting.lean}}
{\small\raggedright \lean{SetCapacity.fixed_H_card}, \lean{SetCapacity.fixed_K_card}, \lean{SetCapacity.prime_capacity}.\par}

\paragraph{\texttt{CapacitySetOrbit.lean}}
{\small\raggedright \lean{SetCapacity.capacity_average_formula}, \lean{SetCapacity.capacity_orbit_formula}, \lean{SetCapacity.exact_capacity}.\par}

\paragraph{\texttt{CapacitySetRate.lean}}
{\small\raggedright \lean{SetCapacity.capacity_lower}, \lean{SetCapacity.logarithmic_rate_tendsto}, \lean{SetCapacity.relative_error}.\par}

\paragraph{\texttt{CapacitySubcodeBarrier.lean}}
{\small\raggedright \lean{CapacitySubcode.above_half_seed_barrier}.\par}

\paragraph{\texttt{CapacitySubcodeMixture.lean}}
{\small\raggedright \lean{CapacitySubcode.exists_above_half_barrier}.\par}

\paragraph{\texttt{CapacityTripleBinary.lean}}
{\small\raggedright \lean{CapacityTriple.exists_width_stable_hard_source}.\par}

\paragraph{\texttt{CapacityUniversal.lean}}
{\small\raggedright \lean{IsSharedCode.two_class_capacity_bound}.\par}

\paragraph{\texttt{CompositionAgreement.lean}}
{\small\raggedright \lean{full_input_agreement_probability}, \lean{repeated_class_agreement_probability}.\par}

\paragraph{\texttt{CompositionAnchoredEndpoint.lean}}
{\small\raggedright \lean{anchored_adjacent_classes_distinct}, \lean{anchored_at_most_two_classes_impossible}.\par}

\paragraph{\texttt{CompositionChannels.lean}}
{\small\raggedright \lean{commonMessage_interval_bound}.\par}

\paragraph{\texttt{CompositionEndpointBlowup.lean}}
{\small\raggedright \lean{endpointFirst_classUniform}, \lean{endpointFirst_involutive}, \lean{endpointSecond_equivariant}, \lean{endpointSecond_involutive}, \lean{endpoint_composite_getter}, \lean{endpoint_composite_requires_singletons}.\par}

\paragraph{\texttt{CompositionLinear.lean}}
{\small\raggedright \lean{exists_twoHop_cutoff}, \lean{twoHop_coprime_progression_dimension_le}, \lean{twoHop_dimension_le}, \lean{twoHop_progression_gcd}, \lean{twoHop_two_classes_le_sixteen}.\par}

\paragraph{\texttt{CompositionMultiplierClasses.lean}}
{\small\raggedright \lean{fieldMultiplierConjugate_eq}, \lean{multiplier_class_count}, \lean{multiplier_class_count_prefix}, \lean{multiplier_class_residue}, \lean{ratio_multiplier_conjugate_covariance}.\par}

\paragraph{\texttt{CompositionMultiplierGetters.lean}}
{\small\raggedright \lean{actual_multiplier_getters_card}, \lean{multiplier_binary_degree_eq}, \lean{multiplier_degree_le}.\par}

\paragraph{\texttt{CompositionOrbitAssignment.lean}}
{\small\raggedright \lean{OrbitAssignment.allowed_iff_divides}, \lean{OrbitAssignment.cyclic_balanced_map_iff_orbit_assignment}, \lean{OrbitAssignment.orbit_map_exists_iff}.\par}

\paragraph{\texttt{CompositionOrbitResidues.lean}}
{\small\raggedright \lean{balanced_fixed_word_residue}, \lean{balanced_fixed_word_small_impossible}, \lean{fixedTupleEquiv}, \lean{invariant_fiber_fixed_residue}, \lean{no_balanced_free_off_zero}, \lean{no_balanced_semiregular}, \lean{tuple_orbit_divisor}.\par}

\paragraph{\texttt{CompositionOrderFive.lean}}
{\small\raggedright \lean{orderFivePerm_fixed}, \lean{orderFivePerm_order}, \lean{orderFivePerm_zero}, \lean{orderFive_adjacent_class_impossible}, \lean{orderFive_conjugate_difference}, \lean{orderFive_no_balanced_covariance}, \lean{orderFive_pair_nontransitive}.\par}

\paragraph{\texttt{CompositionProgression.lean}}
{\small\raggedright \lean{commonMessage_finrank_le}.\par}

\paragraph{\texttt{CompositionRankField.lean}}
{\small\raggedright \lean{exists_twoHop_rank_threshold}, \lean{twoHop_rank_interval_any_field}, \lean{twoHop_rank_superlinear}.\par}

\paragraph{\texttt{CompositionRatio.lean}}
{\small\raggedright \lean{ratioDecoder_surjective}, \lean{ratio_eight_fiber_cards}, \lean{ratio_nonzero_fiber_card}, \lean{ratio_zero_fiber_card}.\par}

\paragraph{\texttt{CompositionRecurrence.lean}}
{\small\raggedright \lean{recurrence_finrank_le}.\par}

\paragraph{\texttt{CompositionRobustAgreement.lean}}
{\small\raggedright \lean{agreement_error_cover}, \lean{robust_full_input_count}, \lean{robust_repeated_class_count}.\par}

\paragraph{\texttt{CompositionScatteredCompletion.lean}}
{\small\raggedright \lean{scatteredCompletion_equivariant}, \lean{scatteredFirstFactor_classUniform}, \lean{scatteredFirstFactor_fixed_preimage}, \lean{scatteredFirstFactor_source}, \lean{scattered_one_class_impossible}.\par}

\paragraph{\texttt{CompositionScatteredSource.lean}}
{\small\raggedright \lean{scatteredSource_finrank}, \lean{scattered_orbit_injective}, \lean{scattered_punctured_orbits_disjoint}.\par}

\paragraph{\texttt{CompositionSingleSiteEndpoint.lean}}
{\small\raggedright \lean{single_at_most_two_classes_impossible}.\par}

\paragraph{\texttt{CompositionSparseObstruction.lean}}
{\small\raggedright \lean{scatteredFirstFactor_disagreement}, \lean{scattered_remaining_factor_singletons}.\par}

\paragraph{\texttt{CompositionSparsePatch.lean}}
{\small\raggedright \lean{columnPatch_involutive}, \lean{columnPatch_nonzero_column}, \lean{columnPatch_support_card}.\par}

\paragraph{\texttt{CompositionWordAction.lean}}
{\small\raggedright \lean{transitive_pair_impossible}.\par}

\paragraph{\texttt{ConstructionBlockDefects.lean}}
{\small\raggedright \lean{finrank_upperBlockAgreement_le}, \lean{upperBlockDescriptor_injective}.\par}

\paragraph{\texttt{ConstructionCapacity.lean}}
{\small\raggedright \lean{periodic_linear_capacity_algebraicClosure_upper}, \lean{periodic_linear_capacity_lower}, \lean{periodic_linear_capacity_split}.\par}

\paragraph{\texttt{ConstructionCyclicExact.lean}}
{\small\raggedright \lean{binaryCyclicDefect_ker}, \lean{binarySeed_full_rule}, \lean{finrank_binaryCyclicLaneDefect_ker}, \lean{mem_cyclicLiftCode_iff_full_rule}.\par}

\paragraph{\texttt{ConstructionLinearLift.lean}}
{\small\raggedright \lean{cyclic_lift}.\par}

\paragraph{\texttt{ConstructionSeed.lean}}
{\small\raggedright \lean{binaryCyclicCode_interval_excess}, \lean{binaryCyclicCode_linear}, \lean{binaryCyclicPartition_surjective}, \lean{finrank_binaryCyclicCode}.\par}

\paragraph{\texttt{Model.lean}}
{\small\raggedright \lean{ClassUniform}, \lean{IsSharedCode}.\par}

\paragraph{\texttt{PeriodTransport.lean}}
{\small\raggedright \lean{fixed_conjugate_iff}, \lean{fixed_image_preserved}, \lean{literal_period_complete}.\par}

\paragraph{\texttt{ResidualEntropy.lean}}
{\small\raggedright \lean{actual_correct_entropy_lower}, \lean{residual_entropy_lower}, \lean{residual_span_information}, \lean{residual_uniform_entropy_lower}.\par}

\paragraph{\texttt{ResidualOutputSpan.lean}}
{\small\raggedright \lean{actual_correct_span_card_le}.\par}

\paragraph{\texttt{ResidualPropagation.lean}}
{\small\raggedright \lean{actual_readout_insertions}, \lean{actual_readout_sum}, \lean{residual_readout_sum}.\par}

\paragraph{\texttt{ResidualRowColumn.lean}}
{\small\raggedright \lean{RowColumn.agreement_card_le}, \lean{RowColumn.helper_agreement_card_le}, \lean{RowColumn.identity_agreement_card}.\par}

\paragraph{\texttt{ResidualRowColumnSharp.lean}}
{\small\raggedright \lean{RowColumn.exists_sharp_helper}, \lean{RowColumn.helper_bound_attained}.\par}

\paragraph{\texttt{ResidualSeparable.lean}}
{\small\raggedright \lean{Separable.rank_defect_lower}.\par}

\paragraph{\texttt{ResidualSeparableSharp.lean}}
{\small\raggedright \lean{Separable.initialProjection_rank}, \lean{Separable.rank_defect_sharp}.\par}

\paragraph{\texttt{ResidualSpan.lean}}
{\small\raggedright \lean{finrank_preimage_le_kernel_add_intersection}, \lean{residualSpanDescriptor_injective}, \lean{residual_span_card_le}.\par}

\paragraph{\texttt{SubgroupClasses.lean}}
{\small\raggedright \lean{covariance_no_class_collision}, \lean{exact_periodic_first_class_bound}, \lean{exact_periodic_first_class_size_bound}, \lean{orbit_classes_injective}.\par}

\paragraph{\texttt{UniformSandwich.lean}}
{\small\raggedright \lean{uniform_sandwich_finite}, \lean{uniform_sandwich_injective_classes}.\par}

\paragraph{Other declarations named in the text}
{\small\raggedright \lean{CapacityTernarySpectrum.Good.source_bound}, \lean{CapacityTernarySpectrum.bad_probability}, \lean{IsSharedCode.intrinsic_retention_bound}, \lean{LineMarker.affineMarker_exceeds_half}, \lean{actual_correct_card_le}, \lean{actual_sparse_card_le}, \lean{binary_progression_nine_finite}, \lean{exists_progression_bound}, \lean{mem_liftedCode_iff_full_rule}, \lean{residual_branch_card_le}, \lean{residual_branch_entropy_lower}.\par}

\paragraph{Imported declarations}
{\small\raggedright \lean{LeanMachine.sparseMenu_log_div_tendsto_zero}.\par}

\section{Building and verifying}\label{sec:verification}
The repository must be cloned with its sibling paper directories. The Lean
package pins Lean and Mathlib to version 4.34.1 and imports earlier libraries
through relative paths. From the paper directory, run
\begin{verbatim}
python3 tools/lean_index.py
cd lean
lake exe cache get
lake build
lake env lean PaperAxioms.lean
cd ../paper
latexmk -pdf -interaction=nonstopmode -halt-on-error main.tex
\end{verbatim}
The generated index and axiom file list every declaration cited by a formal-status
line. The permitted axioms are \texttt{propext}, \texttt{Classical.choice} and
\texttt{Quot.sound}. There are no admitted proofs, extra axioms or native
computation shortcuts. Build products under \texttt{.lake/} are local and are
not part of the published sources. The README records the remaining formalization
obligations; a successful build does not certify the written-only parts.

One quotient-probability proof uses a local \texttt{maxHeartbeats} limit of
800000 for elaboration. This does not change its statement or kernel checking;
the remaining probability modules use the default limit.

\bibliographystyle{plainurl}
{\small\bibliography{references}}
\end{document}
